\chapter{Geodesics and Distance}
\label{chap:pet-geodesics-and-distance}

This chapter defines mixed partials, geodesics, curve length, and Riemannian distance;
proves the first variation formula, introduces exponential and normal coordinates,
and records local isometry, completeness, cut-locus, and injectivity-radius results.

\section{Mixed Partials}

Let $c:\Omega\to M$, where $\Omega\subset\mathbb R^m$ has coordinates $t^i$ and
$M$ has coordinates $x^i$.  The first partial $\partial c/\partial t^i$ is the
velocity obtained by varying $t^i$.  Tangent-valued second partials are required
to satisfy
\begin{enumerate}
\item[(1)] equality of mixed partials,
$\dfrac{\partial^2c}{\partial t^i\partial t^j}=\dfrac{\partial^2c}{\partial t^j\partial t^i}$;
\item[(2)] the product rule,
$\dfrac{\partial}{\partial r^k}g\Big(\dfrac{\partial c}{\partial r^i},\dfrac{\partial c}{\partial r^j}\Big)
=g\Big(\dfrac{\partial^2c}{\partial r^k\partial r^i},\dfrac{\partial c}{\partial r^j}\Big)
+g\Big(\dfrac{\partial c}{\partial r^i},\dfrac{\partial^2c}{\partial r^k\partial r^j}\Big)$.
\end{enumerate}
These conditions are the mixed-partial analogues of torsion-freeness and metric
compatibility.

\begin{lemma}[uniqueness of mixed partials]
  \label{lem:pet-ch5-mixed-partials-uniqueness}
  \lean{PetersenLib.mixedPartials_uniqueness}
  \source{petersen:page-0184}
  \leanok
  \dcref{ch5:5.1.1}
  There is at most one way of defining mixed partials of maps $c:\Omega\to M$ so
  that properties (1) and (2) above hold.
\end{lemma}
\begin{proof}
  \leanok
  \uses{lem:pet-ch5-mixed-partials-uniqueness}
  \sketch
  Applying the product rule cyclically and using symmetry gives
  \[
    2g\Big(\frac{\partial^2c}{\partial r^i\partial r^j},\frac{\partial c}{\partial r^k}\Big)
    =\frac{\partial}{\partial r^i}g\Big(\frac{\partial c}{\partial r^j},\frac{\partial c}{\partial r^k}\Big)
    +\frac{\partial}{\partial r^j}g\Big(\frac{\partial c}{\partial r^k},\frac{\partial c}{\partial r^i}\Big)
    -\frac{\partial}{\partial r^k}g\Big(\frac{\partial c}{\partial r^i},\frac{\partial c}{\partial r^j}\Big).
  \]
  The other four expanded terms cancel in symmetric pairs.  At an image point
  $p$, extend $c$ by a parameter $t^0$ with
  $\partial_{t^0}\tilde c|_p=v$ for arbitrary $v\in T_pM$.  Taking $k=0$ in the
  displayed identity determines
  $g(\partial_{t^i}\partial_{t^j}c,v)$ from first-order data.  Nondegeneracy of
  $g$ therefore determines the mixed partial itself.
\end{proof}

\begin{theorem}[existence of mixed partials]
  \label{thm:pet-ch5-mixed-partials-existence}
  \lean{PetersenLib.mixedPartials_existence}
  \source{petersen:page-0185}
  \uses{lem:pet-ch5-mixed-partials-uniqueness}
  \leanok
  \dcref{ch5:5.1.2}
  It is possible to define mixed partials of maps $c:\Omega\to M$ in a coordinate
  system so that properties (1) and (2) of \cref{lem:pet-ch5-mixed-partials-uniqueness}
  hold.
\end{theorem}
\begin{proof}
  \leanok
  \uses{thm:pet-ch5-mixed-partials-existence}
  \sketch
  Work in a coordinate neighborhood containing the image of $c$.  For parameters
  $s,t$, define
  \[
    \frac{\partial^2c}{\partial t\partial s}
    :=\Big(\frac{\partial^2c^k}{\partial t\partial s}+\frac{\partial c^i}{\partial s}\frac{\partial c^j}{\partial t}\Gamma^k_{ji}\Big)\partial_k.
  \]
  Ordinary mixed-partial symmetry and $\Gamma^k_{ij}=\Gamma^k_{ji}$ give (1).
  For a third parameter $u$, the identities
  $\partial_kg_{ij}=\Gamma_{k,ij}+\Gamma_{k,ji}$ and
  $g_{ij}\Gamma^i_{kl}=\Gamma_{j,kl}$ yield
  \begin{align*}
    \frac{\partial}{\partial t}g\Big(\frac{\partial c}{\partial s},\frac{\partial c}{\partial u}\Big)
    &=\frac{\partial}{\partial t}\Big(g_{ij}\frac{\partial c^i}{\partial s}\frac{\partial c^j}{\partial u}\Big)\\
    &=(\partial_kg_{ij})\frac{\partial c^k}{\partial t}\frac{\partial c^i}{\partial s}\frac{\partial c^j}{\partial u}
    +g_{ij}\frac{\partial^2c^i}{\partial t\partial s}\frac{\partial c^j}{\partial u}
    +g_{ij}\frac{\partial c^i}{\partial s}\frac{\partial^2c^j}{\partial t\partial u}\\
    &=g_{ij}\Big(\frac{\partial^2c^i}{\partial t\partial s}+\Gamma^i_{kl}\frac{\partial c^k}{\partial t}\frac{\partial c^l}{\partial s}\Big)\frac{\partial c^j}{\partial u}
    +g_{ij}\frac{\partial c^i}{\partial s}\Big(\frac{\partial^2c^j}{\partial t\partial u}+\Gamma^j_{kl}\frac{\partial c^k}{\partial t}\frac{\partial c^l}{\partial u}\Big)\\
    &=g\Big(\frac{\partial^2c}{\partial t\partial s},\frac{\partial c}{\partial u}\Big)
    +g\Big(\frac{\partial c}{\partial s},\frac{\partial^2c}{\partial t\partial u}\Big),
  \end{align*}
  which is (2).  Uniqueness from
  \cref{lem:pet-ch5-mixed-partials-uniqueness} makes this local formula
  coordinate independent.
\end{proof}

For $M\subset\tilde M$, write $v=v^\top+v^\perp$ for the tangential and normal
components of $v\in T_p\tilde M$.

\begin{proposition}[mixed partials in submanifolds]
  \label{prop:pet-ch5-mixed-partials-submanifold}
  \lean{PetersenLib.mixedPartialSubmanifoldProjection}
  \source{petersen:page-0187}
  \uses{thm:pet-ch5-mixed-partials-existence}
  \leanok
  \dcref{ch5:5.1.3}
  If $c:\Omega\to M\subset\tilde M$ and $\frac{\partial^2c}{\partial r\partial\theta}\in T_p\tilde M$
  is the mixed partial computed in $\tilde M$, then
  \[
    \left(\frac{\partial^2c}{\partial r\partial\theta}\right)^\top\in T_pM
  \]
  is the mixed partial computed in $M$.
\end{proposition}
\begin{proof}
  \leanok
  \uses{thm:pet-ch5-mixed-partials-existence,lem:pet-ch5-mixed-partials-uniqueness}
  \sketch
  The ambient Koszul identity is
  \[
    2\tilde g\!\left(\frac{\partial^2c}{\partial r\partial\theta},\frac{\partial c}{\partial k}\right)
    =\frac{\partial}{\partial r}\tilde g\!\left(\frac{\partial c}{\partial\theta},\frac{\partial c}{\partial k}\right)
    +\frac{\partial}{\partial\theta}\tilde g\!\left(\frac{\partial c}{\partial k},\frac{\partial c}{\partial r}\right)
    -\frac{\partial}{\partial k}\tilde g\!\left(\frac{\partial c}{\partial r},\frac{\partial c}{\partial\theta}\right).
  \]
  All first partials are tangent, so the right side uses only
  $g=\tilde g|_{TM}$, and pairing the ambient mixed partial with a tangent vector
  equals pairing its tangential projection.  As the last partial ranges over
  $T_pM$, uniqueness identifies that projection with the intrinsic mixed partial.
\end{proof}

\section{Geodesics}

\begin{definition}[acceleration of a curve]
  \label{def:pet-ch5-acceleration}
  \lean{PetersenLib.curveAcceleration}
  \source{petersen:page-0187,petersen:page-0188}
  \leanok
  For a $C^\infty$ curve $c:I\to M$ the \emph{acceleration} is $\ddot c=\frac{d^2c}{dt^2}$.
  In local coordinates,
  \[
    \ddot c=\frac{d^2c^k}{dt^2}\partial_k+\frac{dc^l}{dt}\frac{dc^j}{dt}\Gamma^k_{lj}\partial_k.
  \]
\end{definition}

\begin{definition}[geodesic; arclength parametrization]
  \label{def:pet-ch5-geodesic}
  \lean{PetersenLib.IsGeodesic,PetersenLib.geodesicCoordinateEquation,PetersenLib.IsChartGeodesicOn}
  \source{petersen:page-0188}
  \uses{def:pet-ch5-acceleration}
  \leanok
  A $C^\infty$ curve $c:I\to M$ with vanishing acceleration, $\ddot c=0$, is a
  \emph{geodesic}. If $c$ is a geodesic its speed $|\dot c|=\sqrt{g(\dot c,\dot c)}$
  is constant, since $\frac{d}{dt}g(\dot c,\dot c)=2g(\ddot c,\dot c)=0$; equivalently
  $c$ is parametrized proportionally to arc length. If $|\dot c|\equiv1$, $c$ is
  \emph{parametrized by arclength}. In local coordinates on $U\subset M$, $c:I\to U$
  is a geodesic iff its components satisfy the second-order system
  \[
    \ddot c^k(t)=-\dot c^l(t)\dot c^j(t)\,\Gamma^k_{lj}|_{c(t)},\qquad k=1,\dots,n,
  \]
  giving the initial-value problem $c(0)=q$, $\dot c(0)=\dot c^i(0)\partial_i|_q$ for
  specified $q\in U$ and initial velocity.
\end{definition}

\begin{remark}[distance functions produce geodesics]
  \label{rem:pet-ch5-distance-function-geodesics}
  \lean{PetersenLib.distanceFunctionIntegralCurvesAreGeodesics, PetersenLib.distanceFunctionIntegralCurvesAreGeodesicsOn}
  \source{petersen:page-0188}
  \uses{def:pet-ch5-geodesic, lem:pet-ch5-distance-function-segments, thm:pet-ch5-energy-minima-geodesics}
  \leanok
  \dcref{ch5:5.2.1}
  If $r:U\to\mathbb R$ is a distance function then the integral curves of
  $\partial_r=\nabla r$ are geodesics.  (Petersen derives this from the covariant
  identity $\nabla_{\partial_r}\partial_r=0$; the formalization instead argues
  variationally — an integral curve has unit speed and minimizes length in $U$ by
  \cref{lem:pet-ch5-distance-function-segments}, hence minimizes energy by
  Cauchy–Schwarz, hence is a geodesic by
  \cref{thm:pet-ch5-energy-minima-geodesics} — and does not formalize the
  covariant identity itself, which has no counterpart on the Lean side: the
  Hessian operator there takes an arbitrary affine connection as an argument, and
  no canonical Levi-Civita connection is attached to $g$.)  The formalized curve
  is $C^\infty$ on all of $\mathbb R$ and an integral curve of $\nabla r$ on
  $[a,b]$; the conclusion is the geodesic equation on $(a,b)$.
  The theory of geodesics is developed independently of distance functions and is
  ultimately used to establish their existence.
\end{remark}

\begin{theorem}[local uniqueness of geodesics]
  \label{thm:pet-ch5-local-uniqueness}
  \lean{PetersenLib.geodesic_local_uniqueness}
  \source{petersen:page-0189}
  \uses{def:pet-ch5-geodesic}
  \leanok
  \dcref{ch5:5.2.2}
  Let $I_1,I_2$ be intervals with $t_0\in I_1\cap I_2$. If $c_1:I_1\to U$ and
  $c_2:I_2\to U$ are geodesics in a chart $U\subset M$ with $c_1(t_0)=c_2(t_0)$ and
  $\dot c_1(t_0)=\dot c_2(t_0)$, then $c_1|_{I_1\cap I_2}=c_2|_{I_1\cap I_2}$.
\end{theorem}
\begin{proof}
  \leanok
  \uses{thm:pet-ch5-local-uniqueness}
  \sketch
  In coordinates set $z_i=(u_i,\dot u_i)$ with $u_i=\varphi\circ c_i$. Both solve
  $z'=F(z)$, $F(x,w)=(w,-\Gamma_x(w,w))$, and $F$ is locally Lipschitz. The
  equality set of $z_1,z_2$ is closed and, by local ODE uniqueness, open in
  $I_1\cap I_2$. It contains $t_0$, hence equals the whole interval.
\end{proof}

\begin{theorem}[local existence of geodesics]
  \label{thm:pet-ch5-local-existence}
  \lean{PetersenLib.geodesic_local_existence}
  \source{petersen:page-0189}
  \uses{def:pet-ch5-geodesic}
  \leanok
  \dcref{ch5:5.2.3}
  For each $p\in U$ and $v\in\mathbb R^n$ there are a neighborhood $V_1$ of $p$, a
  neighborhood $V_2$ of $v$, and $\varepsilon>0$ such that for each $q\in V_1$,
  $w\in V_2$, there is a geodesic $c_{q,w}:(-\varepsilon,\varepsilon)\to U$ with $c(0)=q$,
  $\dot c(0)=w^i\partial_i|_q$. Moreover $(q,w,t)\mapsto c_{q,w}(t)$ is $C^\infty$ on
  $V_1\times V_2\times(-\varepsilon,\varepsilon)$.
\end{theorem}
\begin{proof}
  \leanok
  \uses{thm:pet-ch5-local-uniqueness}
  \sketch
  By \cref{thm:pet-ch5-local-uniqueness}, the coordinate spray
  $F(x,w)=(w,-\Gamma_x(w,w))$ is smooth. Picard--Lindel\"of with parameters
  gives a flow $Z$ through a fixed time $T$, with smooth dependence on initial
  data in a ball around $z_0=(\varphi(p),0)$.

  Fibre homogeneity $F(x,aw)=(aw,-a^2\Gamma_x(w,w))$ gives the rescaling
  identity between velocity $w$ at time $t$ and velocity $\tfrac tT w$ at time $T$,
  \[
    Z\big((\varphi(q),\tfrac tT w),T\big)_1
      = Z\big((\varphi(q),\lambda w),\tfrac t\lambda\big)_1 ,
  \]
  Choose $\lambda>0$ with $\lambda(\lVert v\rVert+1)$ inside the flow ball,
  set $V_2=B(v,1)$, $V_1=\varphi^{-1}(B(\varphi(p),\nu))$, and
  $\varepsilon=T\lambda$. Then all rescaled data with $q\in V_1$, $w\in V_2$
  and $|t|<\varepsilon$ lie in that ball, so
  \[
    c_{q,w}(t) := \varphi^{-1}\Big(Z\big((\varphi(q),\tfrac tT w),T\big)_1\Big)
  \]
  is defined. Homogeneity shows that it solves the geodesic equation, with
  $c_{q,w}(0)=q$ and $\dot c_{q,w}(0)=w$. Smoothness in $(q,w,t)$ follows by
  composing the smooth rescaling, the initial-data flow, and evaluation at $T$.
\end{proof}

\begin{lemma}[global uniqueness of geodesics]
  \label{lem:pet-ch5-global-uniqueness}
  \lean{PetersenLib.geodesic_global_uniqueness}
  \source{petersen:page-0189}
  \uses{thm:pet-ch5-local-uniqueness}
  \leanok
  \dcref{ch5:5.2.4}
  Let $I_1,I_2$ be open intervals with $t_0\in I_1\cap I_2$. If $c_1:I_1\to M$ and
  $c_2:I_2\to M$ are geodesics with $c_1(t_0)=c_2(t_0)$ and $\dot c_1(t_0)=\dot c_2(t_0)$,
  then $c_1|_{I_1\cap I_2}=c_2|_{I_1\cap I_2}$.
\end{lemma}
\begin{proof}
  \uses{lem:pet-ch5-global-uniqueness}
  \sketch
  Let $A$ be the set where positions and velocities agree.  It is nonempty and
  closed, while local uniqueness in a chart makes it open.  Connectedness of
  $I_1\cap I_2$ gives $A=I_1\cap I_2$.
\end{proof}

\begin{lemma}[uniform short-time existence over a compact set]
  \label{lem:pet-ch5-uniform-time-existence}
  \lean{PetersenLib.geodesic_uniformShortTimeExistence}
  \source{petersen:page-0190}
  \uses{thm:pet-ch5-local-existence}
  \leanok
  If $\tilde K\subset TM$ is compact, there is $\varepsilon>0$ such that for each
  $(q,v)\in\tilde K$ there is a geodesic $c:(-\varepsilon,\varepsilon)\to M$ with
  $c(0)=q$, $\dot c(0)=v$.
\end{lemma}
\begin{proof}
  \uses{thm:pet-ch5-local-existence}
  \sketch
  Local existence gives each $(q,v)\in\tilde K$ a neighborhood in $TM$ and a
  positive existence time.  A finite subcover of compact $\tilde K$ and the
  minimum of the corresponding times provide one $\varepsilon>0$ for all data.
\end{proof}

\begin{proposition}[maximal geodesics leave every compact set]
  \label{prop:pet-ch5-leaves-compact-set}
  \lean{PetersenLib.maximalGeodesic_leavesCompactSet}
  \source{petersen:page-0190}
  \uses{lem:pet-ch5-uniform-time-existence, def:pet-ch5-geodesic}
  \leanok
  Let $c:I=(a,b)\to M$ be a geodesic on a maximal interval of definition, with
  $b<\infty$. Then for every compact $K\subset M$ there is $t_K<b$ such that
  $t_K<t<b$ implies $c(t)\notin K$.
\end{proposition}
\begin{proof}
  \uses{lem:pet-ch5-uniform-time-existence,lem:pet-ch5-global-uniqueness}
  \sketch
  If not, $c(t_j)\in K$ for $t_j\uparrow b$.  Constant speed and positive
  definiteness make the tangent vectors over $K$ with norm at most $|\dot c|$ a
  compact subset of $TM$ (use finitely many bundle charts).
  \Cref{lem:pet-ch5-uniform-time-existence} gives $\varepsilon>0$ for this set.
  For $b-t_j<\varepsilon/2$, the
  geodesic through $(c(t_j),\dot c(t_j))$ agrees with $c$ by
  \cref{lem:pet-ch5-global-uniqueness} and extends it past $b$, contradiction.
\end{proof}

\begin{corollary}[compact manifolds are geodesically complete]
  \label{cor:pet-ch5-compact-manifold-complete}
  \lean{PetersenLib.compactManifold_geodesicallyComplete}
  \source{petersen:page-0190}
  \uses{prop:pet-ch5-leaves-compact-set}
  \leanok
  \dcref{ch5:5.2.5}
  If $M$ is a compact Riemannian manifold, then for each $p\in M$, $v\in T_pM$
  there is a geodesic $c:\mathbb R\to M$ with $c(0)=p$, $\dot c(0)=v$: geodesics exist
  for all time.
\end{corollary}
\begin{proof}
  \uses{prop:pet-ch5-leaves-compact-set,lem:pet-ch5-global-uniqueness,thm:pet-ch5-local-existence}
  \sketch
  Patch the local solutions through $(p,v)$ using
  \cref{lem:pet-ch5-global-uniqueness}; their union is maximal on $(a,b)$.
  If $b<\infty$, \cref{prop:pet-ch5-leaves-compact-set} with $K=M$ would force
  $c(t)\notin M$ near $b$, impossible.  Thus $b=\infty$, and the reversed
  geodesic gives $a=-\infty$.
\end{proof}

\begin{definition}[geodesically complete]
  \label{def:pet-ch5-geodesically-complete}
  \lean{PetersenLib.IsGeodesicallyComplete}
  \source{petersen:page-0190}
  \uses{def:pet-ch5-geodesic}
  \leanok
  A Riemannian manifold on which every geodesic is defined for all time (i.e.\ on
  all of $\mathbb R$) is called \emph{geodesically complete}.
\end{definition}

\begin{lemma}[uniform extension on a compact interval]
  \label{lem:pet-ch5-uniform-neighborhood}
  \lean{PetersenLib.geodesic_uniformExtensionOnCompactInterval}
  \source{petersen:page-0190,petersen:page-0191}
  \uses{thm:pet-ch5-local-existence}
  \leanok
  \dcref{ch5:5.2.6}
  Suppose $c:[a,b]\to M$ is a geodesic on a compact interval. There is a
  neighborhood $V$ in $TM$ of $\dot c(a)$ such that if $v\in V$, then there is a
  geodesic $c_v:[a,b]\to M$ with $\dot c_v(a)=v$.
\end{lemma}
\begin{proof}
  \uses{thm:pet-ch5-local-existence,lem:pet-ch5-global-uniqueness}
  \sketch
  Subdivide $[a,b]$ so local existence supplies neighborhoods $V_i$ of the velocity
  at $b_i$ valid on $[b_i,b_{i+1}]$.  Continuity of the flow permits successive
  shrinkings $U_i\subset U_{i-1}$ with velocities at $b_{i+1}$ in $V_{i+1}$;
  the final $U_{k-1}$ is the required neighborhood.
\end{proof}
\begin{example}[jump in maximal interval]
  \label{ex:pet-ch5-punctured-plane-geodesic}
  \lean{PetersenLib.example_punctured_plane_geodesic_interval}
  \source{petersen:page-0191}
  \uses{def:pet-ch5-geodesically-complete}
  \leanok
  \dcref{ch5:5.2.7}
  In the punctured plane $\mathbb R^2-\{(0,0)\}$ the unit-speed geodesic from
  $(-1,0)$ with tangent $(1,0)$ is defined only on $(-\infty,1)$, but nearby
  geodesics from $(-1,0)$ with tangent $(1+\varepsilon_1,\varepsilon_2)$,
  $\varepsilon_2\ne0$ small, are defined on all of $\mathbb R$: maximal intervals can
  jump up in size but, by \cref{lem:pet-ch5-uniform-neighborhood}, not down.
\end{example}

\begin{example}[endpoints escaping to infinity]
  \label{ex:pet-ch5-region-geodesic-jump}
  \lean{PetersenLib.example_region_geodesic_endpoint_limit}
  \source{petersen:page-0191}
  \uses{lem:pet-ch5-uniform-neighborhood}
  \dcref{ch5:5.2.8}
  On $\{(x,y)\mid -1<xy\}$, $t\mapsto(t,0)$ is a geodesic on all of $\mathbb R$ that
  is a limit as $\varepsilon\to0$ of unit-speed geodesics $t\mapsto(t,-\varepsilon)$,
  each defined only on a finite interval whose endpoints tend to infinity, as
  required by \cref{lem:pet-ch5-uniform-neighborhood}.
\end{example}

\begin{example}[great circles on spheres]
  \label{ex:pet-ch5-sphere-great-circles}
  \lean{PetersenLib.sphereGeodesicsAreGreatCircles}
  \source{petersen:page-0191}
  \uses{def:pet-ch5-geodesic, prop:pet-ch5-mixed-partials-submanifold}
  \dcref{ch5:5.2.9}
  On $S^n(R)\subset\mathbb R^{n+1}$, the acceleration of $c:I\to S^n(R)$ equals the
  Euclidean acceleration in $\mathbb R^{n+1}$ projected onto $S^n(R)$
  (\cref{prop:pet-ch5-mixed-partials-submanifold}), so $c$ is a geodesic iff $\ddot c$
  is normal to $S^n(R)$, i.e.\ proportional to $c$. The great circles
  $c(t)=p\cos(at)+v\sin(at)$ with $p,v\in\mathbb R^{n+1}$, $|p|=|v|=R$, $p\perp v$,
  satisfy this, with $c(0)=p$, $\dot c(0)=av$: for every initial value problem there
  is such a geodesic.
\end{example}
\begin{proof}
  \uses{ex:pet-ch5-sphere-great-circles}
  \sketch
  The formula gives $\ddot c=-a^2c$.  Since $c\cdot c=R^2$, the vector $c$ is
  normal to the sphere; hence the ambient acceleration has zero tangential
  projection.  By \cref{prop:pet-ch5-mixed-partials-submanifold}, the intrinsic
  acceleration vanishes.
\end{proof}

\begin{example}[hyperbolas on hyperbolic space]
  \label{ex:pet-ch5-hyperbolic-hyperbolas}
  \lean{PetersenLib.hyperbolicSpaceGeodesicsAreHyperbolas}
  \source{petersen:page-0192}
  \uses{def:pet-ch5-geodesic, prop:pet-ch5-mixed-partials-submanifold}
  \dcref{ch5:5.2.10}
  On $H^n(R)\subset\mathbb R^{n,1}$, the acceleration is again the Minkowski
  projection of the ambient (Euclidean-formula) acceleration, and
  $T_pH^n(R)=\{v\mid v\perp p\}$ in the Minkowski sense. The curves
  $c(t)=p\cosh(at)+v\sinh(at)$, $p,v\in\mathbb R^{n,1}$, $|p|^2=-R^2$, $|v|^2=R^2$,
  $p\perp v$ (Minkowski), satisfy $\ddot c=a^2c$, proportional to $c$ hence normal to
  $H^n(R)$, so these hyperbolas are the geodesics.
\end{example}

\begin{example}[left-invariant geodesics on Lie groups]
  \label{ex:pet-ch5-lie-group-geodesics}
  \lean{PetersenLib.leftInvariantGeodesicFields}
  \source{petersen:page-0192}
  \uses{def:pet-ch5-geodesic}
  \dcref{ch5:5.2.11}
  On a Lie group $G$ with left-invariant metric, the integral curves of a
  left-invariant vector field $X$ are geodesics iff $\nabla_XX\equiv0$; the
  upper-half-plane Lie group model does not satisfy this, but for a bi-invariant
  metric and left-invariant $X$, $\nabla_XX=\tfrac12[X,X]=0$; all compact Lie groups
  admit bi-invariant metrics.
\end{example}

\section{The Metric Structure of a Riemannian Manifold}

\begin{definition}[piecewise $C^\infty$ curve]
  \label{def:pet-ch5-piecewise-smooth-curve}
  \lean{PetersenLib.IsPiecewiseSmoothCurve}
  \source{petersen:page-0193}
  \leanok
  A curve $c:[a,b]\to M$ is \emph{piecewise $C^\infty$} if it is continuous and there
  is a partition $a=a_1<a_2<\dots<a_k=b$ such that $c|_{[a_i,a_{i+1}]}$ is $C^\infty$
  for $i=1,\dots,k-1$.
\end{definition}

\begin{definition}[length of a curve]
  \label{def:pet-ch5-curve-length}
  \lean{PetersenLib.curveLength,PetersenLib.ContMDiffOn.intervalIntegrable_sqrt_curveSpeedSq,PetersenLib.IsPiecewiseSmoothCurve.intervalIntegrable_sqrt_curveSpeedSq,PetersenLib.curveLength_comp_mul_add,PetersenLib.curveLength_comp_const_sub}
  \source{petersen:page-0193,petersen:page-0194}
  \uses{def:pet-ch5-piecewise-smooth-curve}
  \leanok
  For $c:[a,b]\to M$ piecewise $C^\infty$, the \emph{length} is
  \[
    L(c)=\int_a^b|\dot c(t)|\,dt=\int_a^b\sqrt{g(\dot c(t),\dot c(t))}\,dt,
  \]
  a well-defined finite nonnegative number, invariant under reparametrization by the
  chain and substitution rules. $c$ is \emph{parametrized by arc length} if
  $L(c|_{[a,t]})=t-a$ for all $t\in[a,b]$, equivalently $|\dot c(t)|=1$ for all $t$.
\end{definition}

\begin{proposition}[arclength reparametrization of regular curves]
  \label{prop:pet-ch5-arclength-reparametrization}
  \lean{PetersenLib.regularCurve_arclengthReparametrization,PetersenLib.IsPiecewiseRegularCurve,PetersenLib.piecewiseRegularCurve_arclengthReparametrization,PetersenLib.curveLength_eq_of_unitSpeed_offFinset,PetersenLib.contDiffAt_curveSpeedSq,PetersenLib.curveSpeedSq_nonneg}
  \source{petersen:page-0194}
  \uses{def:pet-ch5-curve-length}
  \leanok
  A regular curve $c:[a,b]\to M$ (velocity never vanishing) admits a reparametrization
  $c\circ\varphi^{-1}:[0,L(c)]\to M$ by an arclength parameter, piecewise smooth with
  unit speed everywhere.
\end{proposition}
\begin{proof}
  \leanok
  \uses{prop:pet-ch5-arclength-reparametrization}
  \sketch
  Put $\varphi(t)=\int_a^t|\dot c(\tau)|\,d\tau$.  Regularity makes
  $\varphi$ strictly increasing and piecewise smooth, with a piecewise smooth
  inverse.  For $\bar c=c\circ\varphi^{-1}$,
  $\dot{\bar c}(s)=\dot c(t)/\varphi'(t)$ at $t=\varphi^{-1}(s)$, hence
  $|\dot{\bar c}|=1$.
\end{proof}

\begin{definition}[Riemannian distance]
  \label{def:pet-ch5-distance}
  \lean{PetersenLib.riemannianDistance,PetersenLib.riemannianDistance_comm,PetersenLib.riemannianDistance_triangle,PetersenLib.riemannianDistance_triangle_of_connected,PetersenLib.exists_isPiecewiseSmoothCurve_connecting}
  \source{petersen:page-0194}
  \uses{def:pet-ch5-curve-length}
  \leanok
  For $p,q\in M$ let $\Omega_{p,q}=\{c:[0,1]\to M\mid c\text{ piecewise }C^\infty,\
  c(0)=p,\ c(1)=q\}$ and define the \emph{distance} $d(p,q)=|pq|=\inf\{L(c)\mid
  c\in\Omega_{p,q}\}$. Immediately $|pq|=|qp|$ and $|pq|\le|pr|+|rq|$; that $|pq|=0$
  only for $p=q$ is established in \cref{thm:pet-ch5-metric-topology}. When the
  metric needs to be specified one writes $|pq|_g$.
\end{definition}

\begin{definition}[metric balls and set distance]
  \label{def:pet-ch5-metric-balls}
  \lean{PetersenLib.metricBall,PetersenLib.metricClosedBall,PetersenLib.setDistance}
  \source{petersen:page-0194}
  \uses{def:pet-ch5-distance}
  \leanok
  $B(p,r)=\{x\in M\mid|px|<r\}$, $\bar B(p,r)=\{x\in M\mid|px|\le r\}$. For
  $A,B\subset M$, $d(A,B)=|AB|=\inf\{|pq|\mid p\in A,q\in B\}$, and
  $B(A,r)=\{x\in M\mid|Ax|<r\}$, $\bar B(A,r)=D(A,r)=\{x\in M\mid|Ax|\le r\}$.
\end{definition}

\begin{example}[infimum not realized]
  \label{ex:pet-ch5-infimum-not-realized}
  \lean{PetersenLib.example_infimum_not_realized}
  \source{petersen:page-0194,petersen:page-0195}
  \uses{def:pet-ch5-distance}
  \leanok
  \dcref{ch5:5.3.1}
  In $\mathbb R^2-\{(0,0)\}$ with the induced metric, $|(-1,0)(1,0)|=2$, but no curve
  of length $2$ from $(-1,0)$ to $(1,0)$ avoids $(0,0)$, so this distance is not
  realized by any curve; in the sense developed later, $\mathbb R^2-\{(0,0)\}$ is
  incomplete.
\end{example}

\begin{definition}[segment]
  \label{def:pet-ch5-segment}
  \lean{PetersenLib.IsSegment,PetersenLib.segmentNotation}
  \source{petersen:page-0195}
  \uses{def:pet-ch5-distance, def:pet-ch5-curve-length}
  \leanok
  A curve $\sigma\in\Omega_{p,q}$ is a \emph{segment} if $L(\sigma)=|pq|$ and $\sigma$
  is parametrized proportionally to arc length. The notation $\overline{pq}$ denotes
  a specific segment on $[0,|pq|]$ with $\overline{pq}(0)=p$, $\overline{pq}(|pq|)=q$.
\end{definition}

\begin{lemma}[distance functions bound length below]
  \label{lem:pet-ch5-distance-function-length-bound}
  \lean{PetersenLib.distanceFunction_curveLength_ge}
  \source{petersen:page-0195}
  \uses{def:pet-ch5-curve-length, def:pet-ch3-distance-function}
  \leanok
  \dcref{ch5:5.3.2a}
  Let $r:U\to\mathbb R$ be a smooth distance function on an open set
  $U\subset(M,g)$. For every $C^\infty$ curve $c:[a,b]\to M$ whose image lies in
  $U$,
  \[
    r(c(b))-r(c(a))\le L(c)|_a^b.
  \]
\end{lemma}
\begin{proof}
  \leanok
  \uses{lem:pet-ch5-distance-function-length-bound}
  \sketch
  The eikonal equation and Cauchy--Schwarz give
  $(r\circ c)'=g(\nabla r,\dot c)\le|\dot c|$. Therefore
  \[
    L(c)=\int_a^b|\dot c|\,dt\ge\int_a^b g(\nabla r,\dot c)\,dt
      =\int_a^b (r\circ c)'\,dt=r(c(b))-r(c(a)),
  \]
    by the fundamental theorem of calculus.
\end{proof}

\begin{lemma}[distance functions and segments]
  \label{lem:pet-ch5-distance-function-segments}
  \lean{PetersenLib.distanceFunction_integralCurvesAreSegments,PetersenLib.distanceFunction_integralCurve_curveLength}
  \source{petersen:page-0195}
  \uses{def:pet-ch5-segment, lem:pet-ch5-distance-function-length-bound}
  \leanok
  \dcref{ch5:5.3.2b}
  If $r:U\to\mathbb R$ is a smooth distance function on open $U\subset(M,g)$, the
  integral curves of $\nabla r$ are \emph{segments in $U$}: an integral curve $c$
  of $\nabla r$ on $[a,b]$ realises $L(c)|_a^b=r(c(b))-r(c(a))=b-a$ and minimises
  length among curves in $U$ with the same endpoints.
\end{lemma}
\begin{proof}
  \leanok
  \uses{lem:pet-ch5-distance-function-segments}
  \sketch
  Along $\dot c=\nabla r$, one has $|\dot c|=1$ and $dr(\dot c)=1$, so both
  $L(c)|_a^b$ and $r(c(b))-r(c(a))$ equal $b-a$.  The bound in
  \cref{lem:pet-ch5-distance-function-length-bound} makes every competitor at
  least this long.
\end{proof}

\begin{remark}[rigidity of length-minimisers]
  \label{rem:pet-ch5-distance-function-rigidity}
  \lean{PetersenLib.distanceFunction_velocity_eq_speed_smul_gradient_of_curveLength_eq}
  \source{petersen:page-0195}
  \uses{lem:pet-ch5-distance-function-length-bound, lem:pet-ch5-distance-function-segments}
  \leanok
  Let $r:U\to\mathbb R$ be a smooth distance function on an open set
  $U\subset(M,g)$, and let $c:[a,b]\to U$ be a smooth curve with $a<b$ attaining
  the bound of \cref{lem:pet-ch5-distance-function-length-bound}, i.e.\
  $L(c)|_a^b=r(c(b))-r(c(a))$. Then
  \[
    \dot c(t)=|\dot c(t)|\,\nabla r|_{c(t)}\qquad\text{for every }t\in[a,b].
  \]
  (Petersen states this for $c\in\Omega_{p,q}(U)$, the case $a=0$, $b=1$; the
  hypothesis $a<b$ is automatic there. It is not removable in general: for $a=b$
  the length hypothesis degenerates to $0=0$ while the conclusion fails.)
\end{remark}
\begin{proof}
  \leanok
  \uses{lem:pet-ch5-distance-function-length-bound}
  \sketch
  Cauchy--Schwarz gives, for every $t\in[a,b]$,
  \[
    (r\circ c)'(t)=dr(\dot c(t))=g(\nabla r|_{c(t)},\dot c(t))\le|\dot c(t)|,
  \]
  Equality holds exactly when $\dot c(t)=|\dot c(t)|\nabla r|_{c(t)}$. Integrating,
  \[
    r(c(b))-r(c(a))=\int_a^b(r\circ c)'\,dt\le\int_a^b|\dot c|\,dt=L(c)|_a^b,
  \]
  The endpoint hypothesis makes the two integrals equal.  Continuity and $a<b$
  then force equality pointwise, giving the claimed velocity formula.
\end{proof}

\begin{remark}[a minimiser is a reparametrised integral curve of $\nabla r$]
  \label{rem:pet-ch5-minimizer-is-reparametrized-integral-curve}
  \lean{PetersenLib.distanceFunction_minimizer_eq_integralCurve_comp}
  \source{petersen:page-0195}
  \uses{rem:pet-ch5-distance-function-rigidity, lem:pet-ch5-distance-function-segments}
  \leanok
  In the situation of \cref{rem:pet-ch5-distance-function-rigidity}, put
  $L=L(c)|_a^b$ and $\varphi(s)=L(c)|_a^s=\int_a^s|\dot c|\,dt$, and suppose
  $\sigma:[0,L]\to U$ is an integral curve of $\nabla r$ with $\sigma(0)=c(a)$.
  Then $c=\sigma\circ\varphi$ on $[a,b]$: a curve realising the bound is a
  reparametrisation of an integral curve of $\nabla r$.

\end{remark}
\begin{proof}
  \leanok
  \uses{rem:pet-ch5-distance-function-rigidity}
  \sketch
  By \cref{rem:pet-ch5-distance-function-rigidity}, $\dot c=\varphi'\nabla r$;
  also $\varphi(a)=0$ and $\varphi(b)=L$.  The chain rule gives
  \[
    (\sigma\circ\varphi)'(s)=\varphi'(s)\,\dot\sigma(\varphi(s))
      =\varphi'(s)\,\nabla r|_{\sigma(\varphi(s))}.
  \]
  Thus $c$ and $\sigma\circ\varphi$ solve the same locally Lipschitz equation
  $y'=\varphi'\nabla r(y)$ and agree at $a$; Grönwall uniqueness gives
  $c=\sigma\circ\varphi$.  The map $\varphi$ is only nondecreasing, so rests of
  $c$ are allowed.
\end{proof}

\begin{example}[Euclidean segments]
  \label{ex:pet-ch5-euclidean-segments}
  \lean{PetersenLib.euclideanSegmentsAreStraightLines}
  \source{petersen:page-0195}
  \uses{lem:pet-ch5-distance-function-segments}
  \dcref{ch5:5.3.3}
  In $\mathbb R^n$, constant-speed straight segments $t\mapsto p+tv$ are segments, by
  \cref{lem:pet-ch5-distance-function-segments} applied to $r(x)=v\cdot x$ for a unit
  vector $v$. Each pair of points is joined by a segment unique up to
  reparametrization.
\end{example}

\begin{example}[circle segments]
  \label{ex:pet-ch5-circle-segment}
  \lean{PetersenLib.circleSegmentsBoundedByPi}
  \source{petersen:page-0195}
  \uses{lem:pet-ch5-distance-function-segments}
  \dcref{ch5:5.3.4}
  For $M=S^1$, $U=S^1-\{(1,0)\}$, the distance function $r(\theta)=\theta$,
  $\theta\in(0,2\pi)$, shows any $c(\theta)=(\cos\theta,\sin\theta)$, $\theta\in
  I\subset(0,2\pi)$, is a segment in $U$; but if the length of $I$ exceeds $\pi$,
  $c$ is not a segment in $S^1$.
\end{example}

\begin{example}[segments on the round sphere]
  \label{ex:pet-ch5-sphere-segments}
  \lean{PetersenLib.sphereSegmentsCharacterization}
  \source{petersen:page-0195,petersen:page-0196}
  \uses{lem:pet-ch5-distance-function-segments}
  \dcref{ch5:5.3.5}
  On $S^n(1)\subset\mathbb R^{n+1}$ with the warped product structure
  $ds_n^2=dr^2+\sin^2(r)ds_{n-1}^2$ centered at $p$: antipodal $p,q$ correspond to
  $r=0,\pi$; otherwise $p=(a,x_0)$, $q=(b,x_0)\in(0,\pi)\times S^{n-1}$, and $r$ is
  smooth on $U\simeq(0,\pi)\times S^{n-1}$. In either case any minimizing curve has
  length $\ge\pi$ resp.\ $\ge|r(p)-r(q)|$, realized by great-circle arcs.
\end{example}
\begin{proof}
  \uses{ex:pet-ch5-sphere-segments}
  \sketch
  For antipodal endpoints, apply \cref{lem:pet-ch5-distance-function-segments} on
  $[\epsilon,b-\epsilon]$ and let $\epsilon\to0$ to obtain $L(c)\ge\pi$;
  great-circle arcs attain equality.  For nonantipodal endpoints, the integral
  curve of $\nabla r$ minimizes inside $U$.  A curve leaving $U$ passes through a
  pole and has length at least
  $\min\{r(p)+r(q),2\pi-r(p)-r(q)\}\ge|r(p)-r(q)|$.
\end{proof}

\begin{example}[segments in hyperbolic space]
  \label{ex:pet-ch5-hyperbolic-segments}
  \lean{PetersenLib.hyperbolicSpaceAllGeodesicsAreSegments}
  \source{petersen:page-0196}
  \uses{ex:pet-ch5-sphere-segments}
  \dcref{ch5:5.3.6}
  The same strategy as for the sphere shows that all geodesics in hyperbolic space
  are segments.
\end{example}

\begin{example}[punctured plane again]
  \label{ex:pet-ch5-punctured-plane-no-segment}
  \lean{PetersenLib.example_punctured_plane_no_segment}
  \source{petersen:page-0196}
  \uses{def:pet-ch5-segment, def:pet-ch5-distance}
  \leanok
  \dcref{ch5:5.3.7}
  In $\mathbb R^2-\{(0,0)\}$, as already noted, not every pair of points is joined by
  a segment.
\end{example}

\begin{theorem}[metric and manifold topologies agree]
  \label{thm:pet-ch5-metric-topology}
  \lean{PetersenLib.riemannianDistance_inducesManifoldTopology,PetersenLib.riemannianMetricSpace,PetersenLib.isOpen_iff_riemannianDistance,PetersenLib.eq_of_riemannianDistance_eq_zero,PetersenLib.exists_riemannianDistance_le_sqrt_mul,PetersenLib.exists_sqrt_mul_le_riemannianDistance,PetersenLib.sqrt_mul_norm_sub_le_curveLength,PetersenLib.sqrt_mul_le_curveLength_of_exit}
  \source{petersen:page-0196}
  \uses{def:pet-ch5-distance}
  \leanok
  \dcref{ch5:5.3.8}
  The metric topology on $M$ induced by $|\cdot|$ coincides with the manifold
  topology.
\end{theorem}
\begin{proof}
  \uses{thm:pet-ch5-metric-topology}
  \leanok
  \sketch
  Choose coordinates at $p$ with $g_{ij}(p)=\delta_{ij}$ and let $g_0$ be the
  coordinate Euclidean metric.  On a sufficiently small Euclidean ball $U$,
  continuity gives positive $\lambda,\mu$, tending to $1$ at $p$, such that
  $\lambda|v|_{g_0}\le|v|_g\le\mu|v|_{g_0}$.

  The Euclidean segment from $p$ to $x\in U$ gives
  $|px|_g\le\mu_0|px|_{g_0}$.  A competitor contained in $U$ has
  $L_g(c)\ge\lambda_0L_{g_0}(c)$; one leaving $U$ has the same lower bound up to
  its first exit and hence at least $\lambda_0\varepsilon$.  Taking infima gives
  \[
    \lambda_0|px|_{g_0}\le |px|_g\le\mu_0|px|_{g_0}.
  \]
  These local comparisons identify the topologies and separate distinct points.
  Since $\lambda,\mu\to1$ at $p$, also
  $|px|_g/|px|_{g_0}\to1$.
\end{proof}

\begin{corollary}[compact manifolds are metrically complete]
  \label{cor:pet-ch5-compact-complete-metric}
  \lean{PetersenLib.compactManifold_metricallyComplete}
  \source{petersen:page-0198}
  \uses{thm:pet-ch5-metric-topology}
  \leanok
  \dcref{ch5:5.3.9}
  If $M$ is compact and $g$ a Riemannian metric on $M$, then $(M,|\cdot|_g)$ is a
  complete metric space.
\end{corollary}
\begin{proof}
  \uses{cor:pet-ch5-compact-complete-metric}
  \leanok
  \sketch
  By \cref{thm:pet-ch5-metric-topology}, the distance topology is the compact
  manifold topology.  Every compact metric space is complete.
\end{proof}

\begin{corollary}[approximation by constant-speed curves]
  \label{cor:pet-ch5-regular-curve-approx}
  \lean{PetersenLib.approximateByConstantSpeedCurve,PetersenLib.IsConstantSpeedCurve}
  \source{petersen:page-0198}
  \uses{thm:pet-ch5-metric-topology, prop:pet-ch5-arclength-reparametrization}
  \leanok
  \dcref{ch5:5.3.10}
  For $c\in\Omega_{p,q}$ and $\epsilon>0$ there is a constant-speed
  $\bar c\in\Omega_{p,q}$ with $L(\bar c)\le(1+\epsilon)L(c)$.
\end{corollary}
\begin{proof}
  \leanok
  \uses{cor:pet-ch5-regular-curve-approx}
  \sketch
  In finitely many comparison charts from \cref{thm:pet-ch5-metric-topology},
  smooth polygonal approximations give a regular $\bar c$ with
  \[
    L_g(\bar c)\le\mu_0L_{g_0}(\bar c)\le\mu_0L_{g_0}(c)\le\frac{\mu_0}{\lambda_0}L_g(c).
  \]
  Shrink the charts so $\mu_0/\lambda_0<1+\epsilon$.  Finally apply
  \cref{prop:pet-ch5-arclength-reparametrization}, which preserves length and makes
  the regular approximation constant speed.
\end{proof}

\begin{lemma}[distance via constant-speed curves]
  \label{lem:pet-ch5-constant-speed-distance}
  \lean{PetersenLib.riemannianDistance_eq_sInf_constantSpeedDistanceSet}
  \source{petersen:page-0198}
  \uses{def:pet-ch5-distance, cor:pet-ch5-regular-curve-approx}
  \leanok
  Let $(M,g)$ be a Riemannian manifold and $p,q\in M$. Then
  \[
    |pq|=\inf\{L(\gamma)\mid\gamma\in\Omega_{p,q}\text{ of constant speed}\},
  \]
  i.e.\ the infimum defining the Riemannian distance is unchanged if the competing
  curves are required to have constant speed. (When $p$ and $q$ lie in different
  components both infima are over the empty set, and the identity is the
  degenerate one carried by the convention for $|pq|$ in that case; no
  connectedness hypothesis is needed.)
\end{lemma}
\begin{proof}
  \leanok
  \uses{cor:pet-ch5-regular-curve-approx}
  \sketch
  Restricting the competitor class can only increase the infimum.  Conversely,
  \cref{cor:pet-ch5-regular-curve-approx} replaces any $c\in\Omega_{p,q}$ by a
  constant-speed $\bar c$ with $L(\bar c)\le(1+\epsilon)L(c)$.  Let
  $\epsilon\downarrow0$ and then infimize over $c$.  Both classes are empty
  simultaneously when the endpoints lie in different components.
\end{proof}

\begin{remark}[other classes of curves]
  \label{rem:pet-ch5-absolutely-continuous-curves}
  \source{petersen:page-0198}
  \uses{lem:pet-ch5-constant-speed-distance, rem:pet-ch5-ex-29}
  \dcref{ch5:5.3.11}
  The distance theory can equivalently be developed using absolutely continuous
  curves; by \cref{cor:pet-ch5-regular-curve-approx} one could also have restricted
  to piecewise smooth constant-speed curves throughout.

  The constant-speed assertion is
  \cref{lem:pet-ch5-constant-speed-distance}. The absolutely-continuous assertion
  belongs to the program of \cref{rem:pet-ch5-ex-29} and is not developed here.
\end{remark}

\begin{definition}[functional distance]
  \label{def:pet-ch5-functional-distance}
  \lean{PetersenLib.functionalDistance,PetersenLib.HasGradientLeOne,PetersenLib.functionalDistance_nonneg}
  \source{petersen:page-0199}
  \uses{def:pet-ch5-distance}
  \leanok
  The \emph{functional distance} is
  \[
    d_F(p,q)=\sup\{|f(p)-f(q)|\mid f:M\to\mathbb R,\ |\nabla f|\le1\text{ on }M\}.
  \]
  Always $d_F\le d$, and $d_F$ also generates the manifold topology; once smooth
  distance functions are established the two distances agree for $p,q$ sufficiently
  close.
\end{definition}

\section{First Variation of Energy}

\begin{definition}[energy functional]
  \label{def:pet-ch5-energy-functional}
  \lean{PetersenLib.energyFunctional,PetersenLib.energyFunctional_nonneg,PetersenLib.IsPiecewiseSmoothCurve.intervalIntegrable_curveSpeedSq}
  \source{petersen:page-0199}
  \uses{def:pet-ch5-curve-length}
  \leanok
  The \emph{arclength functional} is $L(c)=\int_0^1|\dot c|\,dt$, $c\in\Omega_{p,q}$;
  its minima (if they exist) have minimal length but not necessarily the correct
  (constant-speed) parametrization, and $L$ is invariant under reparametrization. The
  \emph{energy functional}
  \[
    E(c)=\frac12\int_0^1|\dot c|^2\,dt,\qquad c\in\Omega_{p,q},
  \]
  depends on parametrization and measures kinetic energy.
\end{definition}

\begin{proposition}[$E$- and $L$-minimizers coincide]
  \label{prop:pet-ch5-energy-length-minimizers}
  \lean{PetersenLib.energyMinimizer_of_constantSpeed_lengthMinimizer,PetersenLib.lengthMinimizer_of_energyMinimizer,PetersenLib.curveLength_sq_le_two_mul_energyFunctional,PetersenLib.IsConstantSpeedCurve.two_mul_energyFunctional}
  \source{petersen:page-0199}
  \uses{def:pet-ch5-energy-functional, cor:pet-ch5-regular-curve-approx}
  \leanok
  \dcref{ch5:5.4.1}
  If $\sigma\in\Omega_{p,q}$ has constant speed and minimizes $L$, then $\sigma$
  minimizes $E$. Conversely, if $\sigma$ minimizes $E$, then $\sigma$ minimizes $L$.
\end{proposition}
\begin{proof}
  \leanok
  \uses{prop:pet-ch5-energy-length-minimizers}
  \sketch
  Cauchy--Schwarz gives $L(c)^2\le2E(c)$, with equality precisely at constant
  speed.  Thus a constant-speed length minimizer $\sigma$ satisfies
  $E(\sigma)=\tfrac12L(\sigma)^2\le E(c)$.

  If instead $\sigma$ minimizes energy, use
  \cref{cor:pet-ch5-regular-curve-approx} to choose a constant-speed $c_\epsilon$
  with $L(c_\epsilon)\le(1+\epsilon)L(c)$. Then
  \[
    L(\sigma)\le\sqrt{2E(\sigma)}\le\sqrt{2E(c_\epsilon)}=L(c_\epsilon)\le(1+\epsilon)L(c),
  \]
  and $\epsilon\downarrow0$ proves length minimality.
\end{proof}

\begin{definition}[variation of a curve]
  \label{def:pet-ch5-variation}
  \lean{PetersenLib.CurveVariation,PetersenLib.IsProperVariation,PetersenLib.CurveVariation.curve,PetersenLib.CurveVariation.isPiecewiseSmoothCurve}
  \source{petersen:page-0200}
  \leanok
  A \emph{variation} of $c:I\to M$ is $\bar c:(-\varepsilon,\varepsilon)\times[a,b]\to
  M$ with $\bar c(0,t)=c(t)$ for $t\in[a,b]$; it is \emph{piecewise smooth} if
  continuous and $[a,b]$ splits into $[a_i,a_{i+1}]$, $i=0,\dots,m-1$, on which
  $\bar c$ is smooth. Write $c_s(t)=\bar c(s,t)$. The \emph{velocity field}
  $\partial\bar c/\partial t$ is well defined on each $[a_i,a_{i+1}]$, with one-sided
  values $\partial\bar c/\partial t^\pm(s,a_i)$ at the breaks. The \emph{variational
  field} is $\partial\bar c/\partial s$, well defined and continuous everywhere,
  smooth on each $(-\varepsilon,\varepsilon)\times[a_i,a_{i+1}]$. When $a=0$, $b=1$,
  $\bar c(s,0)=p$, $\bar c(s,1)=q$ for all $s$, all $c_s\in\Omega_{p,q}$ and the
  variation is called \emph{proper}.
\end{definition}

\begin{lemma}[first variation formula]
  \label{lem:pet-ch5-first-variation-formula}
  \lean{PetersenLib.firstVariationOfEnergy,PetersenLib.hasDerivAt_pieceEnergy,PetersenLib.hasDerivAt_windowEnergy,PetersenLib.hasDerivAt_windowEnergy_chart,PetersenLib.chartReading_acceleration_transfer,PetersenLib.firstVariation_integrand_chart,PetersenLib.hasDerivAt_halfSpeedSq_variation}
  \source{petersen:page-0200,petersen:page-0201}
  \uses{def:pet-ch5-variation, def:pet-ch5-energy-functional}
  \leanok
  \dcref{ch5:5.4.2}
  For a piecewise smooth variation $\bar c:(-\varepsilon,\varepsilon)\times[a,b]\to M$,
  \begin{align*}
    \frac{dE(c_s)}{ds}
    &=-\int_a^bg\!\left(\frac{\partial^2\bar c}{\partial t^2},\frac{\partial\bar c}{\partial s}\right)dt
    +g\!\left(\frac{\partial\bar c}{\partial t},\frac{\partial\bar c}{\partial s}\right)\Big|_{(s,b)}
    -g\!\left(\frac{\partial\bar c}{\partial t},\frac{\partial\bar c}{\partial s}\right)\Big|_{(s,a)}\\
    &\quad+\sum_{i=1}^{m-1}g\!\left(\frac{\partial\bar c}{\partial t^-}-\frac{\partial\bar c}{\partial t^+},
      \frac{\partial\bar c}{\partial s}\right)\Big|_{(s,a_i)}.
  \end{align*}
  (The Lean statement is at the base parameter $s=0$ --- the case every
  application uses; the general $s$ follows by re-basing the variation --- and
  records the boundary contribution in the equivalent per-piece form
  $\sum_i\big(g(\partial_s\bar c,\partial_t^-\bar c)|_{(0,a_{i+1})}
  -g(\partial_s\bar c,\partial_t^+\bar c)|_{(0,a_i)}\big)$, from which the
  displayed endpoint-plus-break-term shape arises by regrouping the finite
  sum.)
\end{lemma}
\begin{proof}
  \leanok
  \uses{lem:pet-ch5-first-variation-formula}
  \sketch
  Split the energy over the smooth pieces; it is therefore enough to compute for
  a smooth variation and sum the endpoint terms at the breaks.  Metric compatibility
  and mixed-partial symmetry (\cref{prop:pet-ch5-mixed-partials-submanifold}) give
  \begin{align*}
    \frac{dE(c_s)}{ds}
    &=\frac{d}{ds}\frac12\int_a^bg\!\left(\frac{\partial\bar c}{\partial t},\frac{\partial\bar c}{\partial t}\right)dt
    =\frac12\int_a^b\frac{\partial}{\partial s}g\!\left(\frac{\partial\bar c}{\partial t},\frac{\partial\bar c}{\partial t}\right)dt\\
    &=\int_a^bg\!\left(\frac{\partial^2\bar c}{\partial s\partial t},\frac{\partial\bar c}{\partial t}\right)dt
    =\int_a^bg\!\left(\frac{\partial^2\bar c}{\partial t\partial s},\frac{\partial\bar c}{\partial t}\right)dt\\
    &=\int_a^b\frac{\partial}{\partial t}g\!\left(\frac{\partial\bar c}{\partial s},\frac{\partial\bar c}{\partial t}\right)dt
      -\int_a^bg\!\left(\frac{\partial\bar c}{\partial s},\frac{\partial^2\bar c}{\partial t^2}\right)dt\\
    &=-\int_a^bg\!\left(\frac{\partial\bar c}{\partial s},\frac{\partial^2\bar c}{\partial t^2}\right)dt
      +g\!\left(\frac{\partial\bar c}{\partial s},\frac{\partial\bar c}{\partial t}\right)\Big|_{(s,b)}
      -g\!\left(\frac{\partial\bar c}{\partial s},\frac{\partial\bar c}{\partial t}\right)\Big|_{(s,a)},
  \end{align*}
  Integration by parts on each piece and the fundamental theorem of calculus give
  the displayed endpoint and velocity-jump contributions; summing the pieces yields
  the formula.
\end{proof}

\begin{theorem}[local minima of energy are geodesics]
  \label{thm:pet-ch5-energy-minima-geodesics}
  \lean{PetersenLib.energyLocalMinimum_isGeodesic}
  \source{petersen:page-0201}
  \uses{def:pet-ch5-energy-functional, def:pet-ch5-geodesic}
  \leanok
  \dcref{ch5:5.4.3}
  If $c\in\Omega_{p,q}$ is a local minimum for $E:\Omega_{p,q}\to[0,\infty)$, then
  $c$ is a smooth geodesic.
\end{theorem}
\begin{proof}
  \leanok
  \uses{lem:pet-ch5-first-variation-formula,def:pet-ch5-geodesic}
  \sketch
  Stationarity holds for every compactly supported chart variation.  With
  variational field $V$, \cref{lem:pet-ch5-first-variation-formula} gives
  \[
    0=\frac{dE(c_s)}{ds}\Big|_{s=0}
    =-\int_a^bg(\ddot c,V)\,dt+\sum_{i=1}^{m-1}g\!\left(\frac{dc}{dt^-}(a_i)-\frac{dc}{dt^+}(a_i),V(a_i)\right).
  \]
  Choosing $V$ supported inside one smooth piece and directed along $\ddot c$
  forces $\ddot c=0$ there.  Then choose $V(a_i)$ along each velocity jump; the
  remaining term is
  \[
    0=\sum_{i=1}^{m-1}\left|\frac{dc}{dt^-}(a_i)-\frac{dc}{dt^+}(a_i)\right|^2,
  \]
  so every jump vanishes.  Position and velocity now agree across each break; the
  geodesic ODE $\ddot u=-\Gamma_u(\dot u,\dot u)$ matches all higher one-sided
  derivatives.  Hence the broken curve is a smooth geodesic.
\end{proof}

\begin{corollary}[characterization of segments]
  \label{cor:pet-ch5-segments-are-geodesics}
  \lean{PetersenLib.segment_isGeodesic}
  \source{petersen:page-0203}
  \uses{thm:pet-ch5-energy-minima-geodesics, prop:pet-ch5-energy-length-minimizers, def:pet-ch5-segment}
  \leanok
  Any piecewise smooth segment is a geodesic.
\end{corollary}
\begin{proof}
  \leanok
  \uses{thm:pet-ch5-energy-minima-geodesics,prop:pet-ch5-energy-length-minimizers,def:pet-ch5-segment}
  \sketch
  A segment is a constant-speed length minimizer.  By
  \cref{prop:pet-ch5-energy-length-minimizers} it minimizes energy, and
  \cref{thm:pet-ch5-energy-minima-geodesics} makes it a smooth geodesic.
\end{proof}

\begin{remark}[energy minima need not be minimal geodesics]
  \label{rem:pet-ch5-sphere-not-locally-minimizing}
  \lean{PetersenLib.longSphereGeodesicsNotLocallyMinimizing}
  \source{petersen:page-0203}
  \uses{thm:pet-ch5-energy-minima-geodesics}
  Local minima of $E$ characterize geodesics, but not every geodesic is a local
  minimum. In $\mathbb R^n$ all geodesics are integral curves of globally defined
  distance functions $u(x)=v\cdot x$; on the unit sphere, however, no geodesic of
  length $>\pi$ is locally minimizing, since it forms part of a great circle whose
  complementary arc has length $<\pi$, giving a shorter competitor.
\end{remark}

\section{Riemannian Coordinates}

\subsection{The Exponential Map}

\begin{definition}[exponential map; homogeneity property]
  \label{def:pet-ch5-exponential-map}
  \lean{PetersenLib.expMap,PetersenLib.geodesicHomogeneity}
  \source{petersen:page-0204}
  \uses{def:pet-ch5-geodesic}
  \leanok
  For $v\in T_pM$ let $c_v$ be the unique geodesic with $c_v(0)=p$, $\dot c_v(0)=v$,
  and $[0,L_v)$ the nonnegative part of its maximal domain. Uniqueness of geodesics
  gives the \emph{homogeneity property} $c_{\alpha v}(t)=c_v(\alpha t)$ for
  $\alpha>0$, $t<L_{\alpha v}$, so $L_{\alpha v}=\alpha^{-1}L_v$. Let $O_p\subset
  T_pM$ be the vectors $v$ with $1<L_v$. The \emph{exponential map} at $p$,
  $\exp_p:O_p\to M$, is $\exp_p(v)=c_v(1)$; equivalently $\exp_p(tv)=c_v(t)$, so
  $\exp_p(v)$ is reached from $p$ by going ``distance'' $|v|$ in direction $v/|v|$.
  Collecting the $\exp_p$ over $p\in M$ gives $\exp:O=\bigcup_pO_p\to M$. By
  \cref{lem:pet-ch5-uniform-neighborhood} and
  \cref{thm:pet-ch5-local-existence}, $O$ is open in
  $TM$ and $\exp$ is smooth; likewise each $O_p\subset T_pM$ is open and $\exp_p$ is
  smooth.
\end{definition}

\begin{proposition}[$\exp_p$ and $E$ are local diffeomorphisms]
  \label{prop:pet-ch5-exp-diffeomorphism-properties}
  \lean{PetersenLib.expMap_localDiffeomorphism,PetersenLib.expDiagonalMap_localDiffeomorphism}
  \source{petersen:page-0204,petersen:page-0205}
  \uses{def:pet-ch5-exponential-map}
  \dcref{ch5:5.5.1a}
  \begin{enumerate}
  \item For $p\in M$, $D\exp_p:T_0(T_pM)\to T_pM$ is nonsingular at the origin, so
  $\exp_p$ is a local diffeomorphism near $0\in T_pM$.
  \item Define $E:O\to M\times M$ by $E(v)=(\pi(v),\exp v)$, $\pi(v)$ the base point
  of $v$. For each $p\in M$, $DE:T_{0_p}(TM)\to T_{(p,p)}(M\times M)$ is nonsingular,
  so $E$ is a diffeomorphism from a neighborhood of the zero section of $TM$ onto a
  neighborhood of the diagonal in $M\times M$.
  \end{enumerate}

  \emph{Formalization status.} Part (1) is proved. Part (2) is proved
  \emph{locally and in coordinates}: \texttt{expDiagonalMap\_localDiffeomorphism}
  gives, for each $p$, a $C^\infty$ diffeomorphism with $C^\infty$ inverse and open
  image for the chart reading $(y,w)\mapsto(y,\exp_y w)$ of $E$ on a product ball
  around $(\varphi_p(p),0)$, together with the $[[I,0],[I,I]]$ shear derivative.
  The node is therefore not yet \texttt{\textbackslash leanok}: what is missing is
  the \emph{global} half — patching the per-$p$ charts into one diffeomorphism from
  a neighbourhood of the zero section of $TM$ onto a neighbourhood of the diagonal
  (the Lean statement mentions no $O\subseteq TM$, no zero section and no
  \texttt{Diffeomorph}) — and the identification of the second component with
  $\exp_q v$, which is an unwritten consequence of geodesic uniqueness.
\end{proposition}
\begin{proof}
  \uses{prop:pet-ch5-exp-diffeomorphism-properties}
  \sketch
  (1) For the canonical isomorphism
  $I_0(v)=\frac{d}{dt}(tv)|_{t=0}:T_pM\to T_0T_pM$, homogeneity gives
  \[
    D\exp_p(I_0(v))=\frac{d}{dt}\exp_p(tv)\Big|_{t=0}
    =\frac{d}{dt}c_{tv}(1)\Big|_{t=0}=\frac{d}{dt}c_v(t)\Big|_{t=0}=\dot c_v(0)=v,
  \]
  Thus $D\exp_p\circ I_0=\mathrm{id}$, and the inverse function theorem applies.

  (2) Under the standard product identifications, varying the base point and fibre
  shows
  \[
    DE|_{(p,0_p)}=\begin{bmatrix}I&0\\ *&I\end{bmatrix}.
  \]
  This is invertible by part (1).  The local inverse-function-theorem neighborhoods
  patch along the zero section, which $E$ maps diffeomorphically to the diagonal.
\end{proof}

\begin{lemma}[$D\exp_p$ is nonsingular at the origin]
  \label{lem:pet-ch5-exp-differential-nonsingular}
  \lean{PetersenLib.expMap_localDiffeomorphism}
  \source{petersen:page-0204,petersen:page-0205}
  \uses{def:pet-ch5-exponential-map}
  \leanok
  \dcref{ch5:5.5.1b}
  For every $p\in M$ there is $\rho>0$ such that $B_\rho(0)\subset\mathcal O_p$,
  $\exp_p$ is injective on $B_\rho(0)$, its chart reading $w\mapsto\varphi_p(\exp_p
  w)$ has strict Fréchet derivative the identity at $0$ (so $D\exp_p$ is nonsingular
  at the origin), and $\exp_p$ pushes the neighbourhood filter of $0\in T_pM$ onto
  that of $p$.  With the inverse function theorem this is part (1) of
  \cref{prop:pet-ch5-exp-diffeomorphism-properties}: $\exp_p$ is a local
  diffeomorphism near $0$.  (Part (2), the map $E=(\pi,\exp)$ near the diagonal, is
  not yet formalized.  It does \emph{not} need the joint $(q,v,t)$-smoothness of the
  geodesic flow, as previously recorded here: $E$ evaluates at a fixed time, and its
  chart reading is already a $C^1$ diffeomorphism onto a neighbourhood of the
  diagonal; only the upgrade to a smooth one remains.)
\end{lemma}
\begin{proof}
  \leanok
  \uses{def:pet-ch5-exponential-map}
  \sketch
  Homogeneity $\exp_p(tv)=c_v(t)$ identifies the strict Fréchet derivative at
  $0$ with $\mathrm{id}$.  The strict inverse function theorem supplies the small
  injective ball and the neighborhood-filter identity.
\end{proof}

\begin{definition}[injectivity radius]
  \label{def:pet-ch5-injectivity-radius}
  \lean{PetersenLib.injectivityRadius}
  \source{petersen:page-0206}
  \uses{prop:pet-ch5-exp-diffeomorphism-properties}
  \leanok
  The \emph{injectivity radius} at $p$, $\operatorname{inj}_p$, is the largest
  $\varepsilon>0$ such that $\exp_p:B(0,\varepsilon)\to M$ is defined and a
  diffeomorphism onto its image.
\end{definition}

\begin{lemma}[positivity of the injectivity radius]
  \label{lem:pet-ch5-injectivity-radius-positive}
  \lean{PetersenLib.injectivityRadius_pos}
  \source{petersen:page-0206}
  \uses{def:pet-ch5-injectivity-radius, lem:pet-ch5-exp-differential-nonsingular}
  \leanok
  For every $p\in M$ the injectivity radius is strictly positive,
  $\operatorname{inj}_p>0$. In particular the supremum defining $\operatorname{inj}_p$
  is a genuine positive number and Riemannian normal coordinates exist at every point.
\end{lemma}
\begin{proof}
  \leanok
  \uses{def:pet-ch5-injectivity-radius,lem:pet-ch5-exp-differential-nonsingular}
  \sketch
  By \cref{lem:pet-ch5-exp-differential-nonsingular}, $\exp_p$ is defined and
  injective on a model-norm ball of radius $\rho>0$.  Norm equivalence gives
  $\|v\|\le\sqrt c,|v|_g$.  Thus
  $\varepsilon=\rho/(\sqrt c+1)$ has
  $B_g(0,\varepsilon)\subset B_{\|\cdot\|}(0,\rho)$, so
  $\operatorname{inj}_p\ge\varepsilon>0$.
\end{proof}

\begin{remark}[Riemannian normal coordinates]
  \label{rem:pet-ch5-riemannian-normal-coordinates}
  \lean{PetersenLib.riemannianNormalCoordinates}
  \source{petersen:page-0206}
  \uses{def:pet-ch5-injectivity-radius}
  \leanok
  Identifying $T_pM$ with $\mathbb R^n$ via a linear isomorphism and using that
  $\exp_p$ is a diffeomorphism near $0$ gives \emph{exponential} or \emph{Riemannian
  normal coordinates} at $p$, unique up to the choice of identification; requiring a
  linear isometry makes them unique up to $O(n)$. Part (2) of
  \cref{prop:pet-ch5-exp-diffeomorphism-properties} says that points $q_1,q_2$
  sufficiently near $p$ are joined by a unique short geodesic $\exp_{q_1}(tv)$,
  $|v|$ small.
\end{remark}

\begin{corollary}[uniform injectivity radius on a compact set]
  \label{cor:pet-ch5-uniform-injectivity-radius-compact}
  \lean{PetersenLib.compactSet_uniformInjectivityRadius,PetersenLib.compactSet_uniformCInftyDiffeo}
  \source{petersen:page-0206}
  \uses{prop:pet-ch5-exp-diffeomorphism-properties}
  \leanok
  Let $K\subset(M,g)$ be compact. There is $\varepsilon>0$ such that for every
  $p\in K$, $\exp_p:B(0,\varepsilon)\to M$ is defined and a diffeomorphism onto its
  image.

  \emph{Formalization status.} Proved in full.
  \texttt{compactSet\_uniformCInftyDiffeo} produces a single $\varepsilon>0$ such
  that for every $q\in K$: $\exp_q v$ is defined for every $v\in T_qM$ with
  $|v|_g<\varepsilon$; the map $v\mapsto\exp_q v$ is injective and $C^\infty$ on
  $\{v\mid|v|_g<\varepsilon\}$; its image is open; and it admits a $C^\infty$ left
  inverse on that image.  ``Diffeomorphism onto its image'' is rendered by these
  unbundled clauses rather than by a bundled diffeomorphism object: $T_qM$ carries
  no charted-space structure in the formalization, so the smoothness clause is
  stated on the model space with the definitional identification of the two.  A
  left inverse on the image is also a right inverse there, since every image point
  is of the form $\exp_q w$.  The quantifier structure ($\forall K$ compact,
  $\exists\varepsilon>0$, $\forall p\in K$, $\forall v$ in the ball), the norm
  ($|\cdot|_g$, the book's) and the hypotheses are Petersen's, and the Lean proof is
  the book's proof.  One caveat:
  \begin{enumerate}
    \item The Lean is phrased with the \emph{intrinsic} moving-foot exponential,
      not with \texttt{PetersenLib.expDomain}.  This is forced, not a convenience:
      $\texttt{expDomain}(g,p)$ is defined by integral curves of the geodesic field
      of the \emph{single chart at $p$}, which the trivialization makes $0$ off
      that chart's source, so an \texttt{expDomain} geodesic can never leave the
      chart at its own initial point.  Chart sources admit no uniform lower bound
      over a compact set, so this corollary stated at \texttt{expDomain} is
      \emph{false}, not merely unproven.  See
      \cref{rem:pet-ch5-injectivity-radius-chart-artifact}.
  \end{enumerate}
\end{corollary}
\begin{proof}
  \leanok
  \uses{cor:pet-ch5-uniform-injectivity-radius-compact}
  \sketch
  Near each $(x,0_x)$, part (2) of
  \cref{prop:pet-ch5-exp-diffeomorphism-properties} makes $E$ a diffeomorphism.
  In a bundle trivialization, continuity of the fibre metric gives one intrinsic
  ball $B(0_p,\varepsilon_x)$ contained in this neighborhood for all $p$ near $x$.
  A finite cover of compact $K$ and the minimum of the $\varepsilon_x$ give the
  asserted uniform radius.
\end{proof}

\begin{remark}[the Lean \texttt{injectivityRadius} is a chart artifact]
  \label{rem:pet-ch5-injectivity-radius-chart-artifact}
  \uses{def:pet-ch5-injectivity-radius}
  The Lean \texttt{PetersenLib.injectivityRadius} is defined through
  \texttt{PetersenLib.expDomain}, which is anchored to the chart at the initial
  point (see the formalization note on
  \cref{cor:pet-ch5-uniform-injectivity-radius-compact}).  It is therefore
  \emph{not} the injectivity radius of this text: on the round circle it is bounded
  by the size of the chart at $p$ rather than by $\pi$.
  \Cref{lem:pet-ch5-injectivity-radius-positive} is unaffected — it asserts only
  $\mathrm{inj}_p>0$, and the chart anchoring only shrinks the set whose supremum
  is taken — but any downstream statement that consumes $\mathrm{inj}_p$
  \emph{quantitatively} must be read against the intrinsic moving-foot exponential
  instead.  This is a known defect of the vendored geodesic engine, whose authors
  record the same point: the chart-anchored witnesses solve the junk-extended
  equation of a single chart.

  The defect propagates to every object defined through \texttt{expMap}.  In
  particular \texttt{PetersenLib.segmentDomain} is defined as
  $\{v\mid t\mapsto\texttt{expMap}(g,p,tv)\text{ is a segment on }[0,1]\}$, and
  \texttt{segmentDomainStarInterior} and \texttt{cutLocus} are defined from it, so
  all three are chart artifacts as well: a geodesic witnessing membership of
  $\operatorname{seg}(p)$ can never leave the chart at $p$, whereas the predicate
  \texttt{IsSegment} that it must satisfy is chart-free.  Consequently
  \cref{lem:pet-ch5-cut-locus-characterization}, read against these definitions, is
  \emph{false}, not merely unproven: see
  \cref{rem:pet-ch5-segment-domain-chart-artifact}.  The repair is the same one
  used for \cref{cor:pet-ch5-uniform-injectivity-radius-compact} and for part~(2)
  of \cref{prop:pet-ch5-local-isometry-properties} — restate at the intrinsic
  maximal geodesic — and it is not yet carried out for the segment domain.
\end{remark}

\begin{definition}[normal exponential map]
  \label{def:pet-ch5-normal-exponential-map}
  \lean{PetersenLib.normalSpace,PetersenLib.normalBundle,PetersenLib.isCompl_tangentSpaceImage_normalSpace,PetersenLib.normalExpMap}
  \source{petersen:page-0206,petersen:page-0207}
  \uses{def:pet-ch5-exponential-map}
  \leanok
  Let $N\subset M$ be a properly embedded submanifold; its normal bundle is
  $T^\perp N=\{v\in T_pM\mid p\in N,\ v\in(T_pN)^\perp\}$, so $T_pM=T_pN\oplus
  (T_pN)^\perp$ orthogonally for $p\in N$. The \emph{normal exponential map}
  $\exp^\perp:O\cap TN^\perp\to M$ restricts $\exp$ to $TN^\perp$.
\end{definition}

\begin{corollary}[tubular neighborhoods]
  \label{cor:pet-ch5-tubular-neighborhood}
  \lean{PetersenLib.tubularNeighborhoodTheorem,PetersenLib.exists_tangentNormalSumEquiv,PetersenLib.exists_ball_normalExp_defined_injOn_contMDiffOn}
  \source{petersen:page-0207}
  \uses{def:pet-ch5-normal-exponential-map, prop:pet-ch5-exp-diffeomorphism-properties}
  \dcref{ch5:5.5.3}
  $D\exp^\perp$ is nonsingular at $0_p$ for all $p\in N$, and there is an open
  neighborhood $U$ of the zero section in $TN^\perp$ on which $\exp^\perp$ is a
  diffeomorphism onto its image $\exp^\perp(U)\subset M$, a \emph{tubular
  neighborhood} of $N$.
\end{corollary}
\begin{proof}
  \uses{cor:pet-ch5-tubular-neighborhood}
  \sketch
  Restrict the block derivative
  $DE|_{(p,0_p)}=\begin{bmatrix}I&0\\ *&I\end{bmatrix}$ from
  \cref{prop:pet-ch5-exp-diffeomorphism-properties} to
  $T_pN\oplus T_p^\perp N$.  It remains nonsingular, so the inverse function
  theorem applies to $\exp^\perp$ at every $0_p$.  Proper embeddedness of $N$
  allows these local inverses to patch over a neighborhood of the zero section.
\end{proof}

\subsection{Short Geodesics Are Segments}

\begin{theorem}[short geodesics are segments]
  \label{thm:pet-ch5-short-geodesics-segments}
  \lean{PetersenLib.expMap_isSegment_unique}
  \source{petersen:page-0207}
  \uses{def:pet-ch5-exponential-map, def:pet-ch5-segment, lem:pet-ch5-short-geodesics-segments-normal}
  \dcref{ch5:5.5.4a}
  Let $(M,g)$ be Riemannian, $p\in M$, $\varepsilon>0$ such that
  $\exp_p:B(0,\varepsilon)\to U\subset M$ is a diffeomorphism onto its image $U$.
  Then $U=B(p,\varepsilon)$, and for each $v\in B(0,\varepsilon)$ the geodesic
  $\exp_p(tv)$, $t\in[0,1]$, is the unique segment of speed $|v|$ from $p$ to
  $\exp_p(v)$ in $M$.

  Not \texttt{\textbackslash leanok}: the Lean furnishes its own $\varepsilon$
  rather than accepting the hypothesised diffeomorphism radius; everything else is
  \cref{lem:pet-ch5-short-geodesics-segments-normal}, including the two clauses
  that were previously open (the ball clause and the uniqueness clause now measure
  the radius in the same intrinsic norm, and the uniqueness is pointwise equality
  with the radial geodesic rather than equality up to reparametrization).  Closing
  the residue needs the $\exists\rho$ Gauss/minimizing-equality chain restated in
  hypothesised-$\rho$ form; the first link of that chain,
  \texttt{PetersenLib.Exponential.gauss\_lemma\_ball\_at}, is done, and the
  single-chart hypothesis it still takes is discharged by the chart-anchoring of
  \texttt{expDomain} recorded in \texttt{ExpChartConfinement.lean}.
\end{theorem}

\begin{lemma}[radial geodesics realize distance and are segments]
  \label{lem:pet-ch5-radial-geodesic-segment}
  \lean{PetersenLib.expMap_riemannianDistance_eq, PetersenLib.exists_expMap_isSegment}
  \source{petersen:page-0207}
  \uses{def:pet-ch5-exponential-map, def:pet-ch5-segment}
  \leanok
  For every $p\in M$ there is $\varepsilon>0$ such that for all $v\in T_pM$ with
  $|v|_g<\varepsilon$ one has $v\in\mathcal O_p$ and
  \[
    d(p,\exp_p v)=|v|_g=\sqrt{g_p(v,v)},
  \]
  and the radial geodesic $\gamma(t)=\exp_p(tv)$, $t\in[0,1]$, is a segment from
  $p$ to $\exp_p v$ of speed $|v|_g$. This is the metric core (the existence half)
  of \cref{thm:pet-ch5-short-geodesics-segments}; the uniqueness of the segment and
  the identity $\exp_p(B(0,\varepsilon))=B(p,\varepsilon)$ are its remaining
  clauses.
\end{lemma}
\begin{proof}
  \leanok
  \uses{lem:pet-ch5-gauss-lemma}
  \sketch
  On a small normal ball, \cref{lem:pet-ch5-gauss-lemma} gives the radial lower
  bound $L(c)\ge|v|_g$ for every curve from $p$ to $\exp_pv$.  The radial geodesic
  $\gamma(t)=\exp_p(tv)$ is smooth by the geodesic-spray ODE and has constant speed
  $|v|_g$, so it realizes the opposite inequality.  Hence
  $d(p,\exp_pv)=L(\gamma)=|v|_g$.  Homogeneity gives
  $L(\gamma)|_0^t=|tv|_g=t|v|_g$, proving proportional arclength and the segment
  assertion.
\end{proof}

\begin{lemma}[uniqueness of the short radial segment, smooth competitors]
  \label{lem:pet-ch5-radial-geodesic-unique}
  \lean{PetersenLib.exists_expMap_smoothSegment_unique}
  \source{petersen:page-0207,petersen:page-0208}
  \uses{lem:pet-ch5-radial-geodesic-segment, lem:pet-ch5-gauss-lemma}
  \leanok
  For every $p\in M$ there is $\varepsilon>0$ such that $B_\varepsilon(0)\subset
  \mathcal O_p$ and, for every $v$ with $|v|<\varepsilon$, every $C^\infty$ curve
  $\sigma:[0,1]\to M$ from $p$ to $\exp_p v$ whose length realizes the radial
  length, $L(\sigma)|_0^1=|v|_g=\sqrt{g_p(v,v)}$, is a monotone reparametrization
  of the radial geodesic: there is a continuous nondecreasing $s:[0,1]\to[0,1]$
  with $s(0)=0$, $s(1)=1$ and $\sigma(t)=\exp_p\bigl(s(t)\,v\bigr)$ for all
  $t\in[0,1]$. Thus, up to reparametrization, $t\mapsto\exp_p(tv)$ is the
  \emph{unique} smooth segment of speed $|v|_g$ from $p$ to $\exp_p v$ — the
  uniqueness clause of \cref{thm:pet-ch5-short-geodesics-segments} restricted to
  globally-$C^\infty$ competitors.
\end{lemma}
\begin{proof}
  \leanok
  \uses{lem:pet-ch5-gauss-lemma,lem:pet-ch5-radial-geodesic-segment}
  \sketch
  In normal coordinates, \cref{lem:pet-ch5-gauss-lemma} gives
  $|\dot\sigma|\ge dr(\dot\sigma)$, with equality exactly in the radial direction.
  A length-realizing competitor cannot escape the normal ball and must have
  equality everywhere.  It therefore traces $\exp_p(sv)$, where $s$ is the
  nondecreasing fraction of accumulated radial length.  The identity
  $\operatorname{pathELength}\sigma=\int_0^1|\dot\sigma|_g$ supplies the formal
  bridge to the same equality argument.
\end{proof}

\begin{lemma}[the uniqueness clause for piecewise-smooth competitors]
  \label{lem:pet-ch5-radial-geodesic-unique-piecewise}
  \lean{PetersenLib.shortGeodesicsAreUniqueSegments}
  \source{petersen:page-0207,petersen:page-0208}
  \uses{lem:pet-ch5-radial-geodesic-unique, def:pet-ch5-piecewise-smooth-curve}
  \leanok
  The strengthening of \cref{lem:pet-ch5-radial-geodesic-unique} from
  globally-$C^\infty$ competitors to the piecewise-$C^\infty$ class
  (\cref{def:pet-ch5-piecewise-smooth-curve}) that
  \cref{def:pet-ch5-segment} and \cref{def:pet-ch5-distance} actually range over.
  For every
  $p\in M$ there is $\varepsilon>0$ such that $B_\varepsilon(0)\subset\mathcal
  O_p$ and, for every $v$ with $|v|<\varepsilon$, every \emph{piecewise}-$C^\infty$
  curve $\sigma:[0,1]\to M$ from $p$ to $\exp_p v$ with
  $L(\sigma)|_0^1=|v|_g$ is a monotone reparametrization of the radial geodesic —
  continuous nondecreasing $s:[0,1]\to[0,1]$, $s(0)=0$, $s(1)=1$,
  $\sigma(t)=\exp_p(s(t)\,v)$ — and moreover the images agree,
  $\sigma([0,1])=\exp_p([0,1]\cdot v)$.

  This is the uniqueness clause of \cref{thm:pet-ch5-short-geodesics-segments} in
  Petersen's own regularity class. The parent theorem is still not
  \texttt{\textbackslash leanok}, but for one reason only: it quantifies
  \emph{universally} over every $\varepsilon$ for which $\exp_p$ is a
  diffeomorphism on $B(0,\varepsilon)$, whereas this lemma supplies its own
  $\varepsilon$.  The norm discrepancy this lemma used to carry (its radius in the
  model norm, the clause $U=B(p,\varepsilon)$ of
  \cref{cor:pet-ch5-metric-ball-image} in $|v|_g$) is resolved in
  \cref{lem:pet-ch5-short-geodesics-segments-normal}, which restates every radius
  intrinsically and upgrades the monotone reparametrization $s$ below to
  $s=\mathrm{id}$.
\end{lemma}
\begin{proof}
  \leanok
  \uses{lem:pet-ch5-radial-geodesic-unique}
  \sketch
  Apply the equality argument of \cref{lem:pet-ch5-radial-geodesic-unique} on
  each smooth piece.  Gauss-radius deficits and the length identities telescope
  over the partition, so every piece is radial with one common nondecreasing
  parameter; additivity also gives
  $\operatorname{pathELength}\sigma=\int|\dot\sigma|_g$.
\end{proof}

\begin{lemma}[short geodesics on a normal ball]
  \label{lem:pet-ch5-short-geodesics-segments-normal}
  \lean{PetersenLib.expMap_isSegment_unique}
  \source{petersen:page-0207}
  \uses{lem:pet-ch5-radial-geodesic-segment, lem:pet-ch5-radial-geodesic-unique-piecewise, cor:pet-ch5-metric-ball-image, def:pet-ch5-segment}
  \leanok
  \dcref{ch5:5.5.4b}
  For every $p\in M$ there is $\varepsilon>0$ such that
  \begin{enumerate}
    \item $\{v:|v|_g<\varepsilon\}\subset\mathcal O_p$;
    \item $\exp_p(\{v:|v|_g<\delta\})=B(p,\delta)$ for every $0<\delta\le\varepsilon$;
    \item for every $v$ with $|v|_g<\varepsilon$, the curve $t\mapsto\exp_p(tv)$,
      $t\in[0,1]$, \emph{is} a segment, and every segment $\sigma:[0,1]\to M$ from
      $p$ to $\exp_p(v)$ satisfies $\sigma(t)=\exp_p(tv)$ for all $t\in[0,1]$.
  \end{enumerate}
  This is \cref{thm:pet-ch5-short-geodesics-segments} verbatim \emph{except} that
  $\varepsilon$ is furnished rather than hypothesised.  Every radius is measured in
  the intrinsic $|\cdot|_g$, so clauses (ii) and (iii) agree on what the ball is.
  Clause (iii) is in fact stronger than the book's: it quantifies over every
  segment competitor — not only those declared to have speed $|v|$, since a
  segment's speed is forced to be $|v|_g$ by
  \cref{lem:pet-ch5-radial-geodesic-segment} — and concludes pointwise equality
  with the radial geodesic rather than equality up to reparametrization.
\end{lemma}
\begin{proof}
  \leanok
  \uses{lem:pet-ch5-radial-geodesic-unique-piecewise,
        lem:pet-ch5-radial-geodesic-segment,cor:pet-ch5-metric-ball-image}
  \sketch
  Choose $\varepsilon$ below the radii in
  \cref{cor:pet-ch5-metric-ball-image},
  \cref{lem:pet-ch5-radial-geodesic-segment}, and
  \cref{lem:pet-ch5-radial-geodesic-unique-piecewise}, using norm coercivity to
  pass from intrinsic to model radii.  For a segment $\sigma$ to $\exp_pv$,
  proportional arclength gives $L(\sigma)|_0^t=t|v|_g$, while uniqueness gives
  $L(\sigma)|_0^t=s(t)|v|_g$.  If $v\ne0$, cancellation yields $s=t$; for $v=0$
  both curves are constant.
\end{proof}

\begin{definition}[radial distance function in exponential coordinates]
  \label{def:pet-ch5-radial-distance-function}
  \lean{PetersenLib.radialDistanceFunction}
  \source{petersen:page-0207}
  \uses{thm:pet-ch5-short-geodesics-segments}
  \leanok
  On $U=\exp_p(B(0,\varepsilon))$ let $r(x)=|\exp_p^{-1}(x)|$, the Euclidean
  distance from the origin in $B(0,\varepsilon)\subset T_pM$, continuously extended
  by $r|_{\partial U}=\varepsilon$. In Cartesian coordinates on $T_pM$,
  $\nabla r=\partial_r=\tfrac1rx^i\partial_i$.
\end{definition}

\begin{lemma}[the Gauss Lemma]
  \label{lem:pet-ch5-gauss-lemma}
  \lean{PetersenLib.gaussLemma}
  \source{petersen:page-0207}
  \uses{def:pet-ch5-radial-distance-function}
  \leanok
  \dcref{ch5:5.5.5}
  On $(U,g)$ the function $r$ has gradient $\nabla r=\partial_r$, where
  $\partial_r=D\exp_p(\partial_r)$ (the pushforward of the flat radial field).
  Equivalently, and as formalized: on a normal ball $B(0,\rho)\subset T_pM$,
  read in the chart at $p$, for all $v$ with $\|v\|<\rho$ and all $w$,
  \[
    \big\langle D\exp_p|_v(v),\,D\exp_p|_v(w)\big\rangle_{\exp_p(v)}=g_p(v,w),
  \]
  where $D\exp_p|_v(v)$ is the radial pushforward $\partial_r$; this is exactly
  $dr(w)=g(\partial_r,w)$ for all $w$, i.e. $\nabla r=\partial_r$.
\end{lemma}
\begin{proof}
  \uses{lem:pet-ch5-gauss-lemma, thm:pet-ch5-short-geodesics-segments, lem:pet-ch5-distance-function-segments}
  \sketch
  Transport an orthonormal basis of $T_pM$ by $\exp_p$.  Then
  $r^2=\sum_i(x^i)^2$, $\partial_r=r^{-1}x^i\partial_i$, and
  $dr=r^{-1}x^idx^i$.  Radial and rotational fields
  $J=-x^i\partial_j+x^j\partial_i$ span each tangent space.  Radial geodesics have
  unit speed, so $dr(\partial_r)=g(\partial_r,\partial_r)=1$, while directly
  $dr(J)=0$.

  Since $[\partial_r,J]=0$, torsion-freeness and the geodesic equation give
  \[
    \partial_rg(\partial_r,J)
    =g(\partial_r,\nabla_J\partial_r)
    =\tfrac12Jg(\partial_r,\partial_r)=0.
  \]
  Thus $g(\partial_r,J)$ is constant on a radial ray.  Near its origin,
  \[
    |g(\partial_r,J)|\le|\partial_r|\,|J|=|J|\le|x^i|\,|\partial_i|+|x^j|\,|\partial_j|
    \le r(x)\big(|\partial_i|+|\partial_j|\big).
  \]
  and the coordinate fields stay bounded, so this constant tends to $0$.  Hence
  $g(\partial_r,J)=0$, proving $dr=g(\partial_r,\cdot)$ on the spanning fields.
\end{proof}
\begin{proof}[Proof of \cref{thm:pet-ch5-short-geodesics-segments}]
  \uses{thm:pet-ch5-short-geodesics-segments,lem:pet-ch5-gauss-lemma,lem:pet-ch5-distance-function-segments}
  \sketch
  By \cref{lem:pet-ch5-gauss-lemma}, $r$ is a distance function on $U-\{p\}$ and
  its integral curves are the unit-speed radial geodesics.  Hence
  \cref{lem:pet-ch5-distance-function-segments} makes them segments and gives
  $U\subset B(p,\varepsilon)$.  For any competitor $c$ from $p$, let $b_0$ be its
  first exit time from $U$ (or its endpoint).  Then
  \[
    L(c|_{[0,b_0)})=\int_0^{b_0}|\dot c|\,dt\ge\int_0^{b_0}dr(\dot c)\,dt=r(c(b_0)),
  \]
  An exit through $\partial U$ costs at least $\varepsilon$; otherwise equality
  holds only for a radial velocity, so every minimizer reparametrizes the radial
  geodesic.  Finally, a point of $B(p,\varepsilon)$ is joined to $p$ by a curve of
  length below $\varepsilon$, which cannot exit $U$ by the same estimate.  Thus
  $B(p,\varepsilon)\subset U$.
\end{proof}

\begin{corollary}[exponential image of balls]
  \label{cor:pet-ch5-metric-ball-image}
  \lean{PetersenLib.expMap_ballImage,PetersenLib.expMap_closedBallImage}
  \source{petersen:page-0209}
  \uses{lem:pet-ch5-radial-geodesic-segment, def:pet-ch5-metric-balls}
  \leanok
  \dcref{ch5:5.5.6}
  There is $\varepsilon>0$ such that for each $\delta$ with $0<\delta\le\varepsilon$
  the exponential map carries the $g$-metric ball onto the Riemannian metric ball,
  both in the open and the closed form,
  \[
    \exp_p\bigl(\{v\in T_pM\mid |v|_g<\delta\}\bigr)=B(p,\delta),\qquad
    \exp_p\bigl(\{v\in T_pM\mid |v|_g\le\delta\}\bigr)=\bar B(p,\delta).
  \]
\end{corollary}
\begin{proof}
  \leanok
  \uses{lem:pet-ch5-radial-geodesic-segment,def:pet-ch5-metric-balls}
  \sketch
  By \cref{lem:pet-ch5-radial-geodesic-segment},
  $d(p,\exp_pv)=|v|_g$, proving the forward inclusion.  Conversely, the normal-ball
  escape estimate gives $q=\exp_pw$ for $q\in B(p,\delta)$, and the same identity
  forces $|w|_g<\delta$.  Replacing strict by non-strict inequalities, with the
  radius below the escape threshold, proves the closed-ball formula.
\end{proof}

\subsection{Properties of Exponential Coordinates}

\begin{remark}[radial isometry reformulation of the Gauss Lemma]
  \label{rem:pet-ch5-radial-isometry}
  \lean{PetersenLib.radialIsometryCondition}
  \source{petersen:page-0210}
  \uses{lem:pet-ch5-gauss-lemma}
  \leanok
  The Gauss Lemma says $\exp_p:B(0,\varepsilon)\to B(p,\varepsilon)$ is a
  \emph{radial isometry}: $g(D\exp_p(\partial_r),D\exp_p(v))=g_p(\partial_r,v)$.
  Indeed $\tfrac1rg_{ij}v^iv^j=g(\tfrac1rx^i\partial_i,v^j\partial_j)=
  \tfrac1r\delta_{ij}x^iv^j=dr(v)$, so this is equivalent to $\nabla r=\partial_r=
  \tfrac1rx^i\partial_i$, i.e.\ to $g_{ij}v^j=\delta_{ij}v^j$ for all $v$.
\end{remark}

\begin{lemma}[exponential coordinates are flat to first order]
  \label{lem:pet-ch5-normal-coords-jet}
  \lean{PetersenLib.expCoordinates_jetFlatness}
  \source{petersen:page-0210,petersen:page-0211}
  \uses{rem:pet-ch5-radial-isometry, lem:pet-ch5-exp-differential-nonsingular}
  \leanok
  \dcref{ch5:5.5.7a}
  In exponential coordinates at $p$ the metric is differentiable at the origin, its
  value there is $g_p$ (that is, $g_{ij}|_p=\delta_{ij}$ in a $g_p$-orthonormal
  frame), and all of its first partial derivatives vanish: $\partial_kg_{ij}|_p=0$.
\end{lemma}
\begin{proof}
  \leanok
  \uses{rem:pet-ch5-radial-isometry,lem:pet-ch5-exp-differential-nonsingular}
  \sketch
  Let $g_x$ denote the metric in exponential coordinates.  Its $C^1$ regularity
  follows from \cref{lem:pet-ch5-exp-differential-nonsingular}, and
  $D\exp_p|_0=\mathrm{id}$ gives $g_0=g_p$.  The radial identity in
  \cref{rem:pet-ch5-radial-isometry} says $g_x(x,w)=g_p(x,w)$.  Substituting
  $x=tv$ gives
  \[
    t\,g_{tv}(v,w)=g_{tv}(tv,w)=g_p(tv,w)=t\,g_p(v,w),
  \]
  so $g_{tv}(v,w)$ is constant in $t$.  Put
  $D_u(a,b)=\partial_u(g_x(a,b))|_0$.  Differentiation gives $D_v(v,w)=0$;
  polarizing yields $D_a(b,w)=-D_b(a,w)$.  Symmetry in the last two arguments then
  gives
  \[
    D_a(b,c)=-D_b(a,c)=-D_b(c,a)=D_c(b,a)=D_c(a,b)=-D_a(c,b)=-D_a(b,c),
  \]
  hence $D=0$, equivalently $\partial_kg_{ij}|_p=0$.
\end{proof}

\begin{lemma}[first-order flatness of exponential coordinates]
  \label{lem:pet-ch5-normal-coords-first-order}
  \lean{PetersenLib.expCoordinates_firstOrderFlatness}
  \source{petersen:page-0210}
  \uses{rem:pet-ch5-radial-isometry, lem:pet-ch5-normal-coords-jet}
  \leanok
  \dcref{ch5:5.5.7b}
  In exponential coordinates, $g_{ij}=\delta_{ij}+O(r^2)$.
\end{lemma}
\begin{proof}
  \leanok
  \uses{lem:pet-ch5-normal-coords-jet}
  \sketch
  By \cref{lem:pet-ch5-normal-coords-jet}, $g_{ij}(0)=\delta_{ij}$ and
  $\partial_kg_{ij}(0)=0$.  Differentiating the radial identity
  $g_{ij}x^j=\delta_{ij}x^j$ once gives
  \[
    \delta_{ik}=\partial_k\Big(\sum_jg_{ij}x^j\Big)=(\partial_kg_{ij})x^j+g_{ik}.
  \]
  A second derivative gives
  \[
    0=\partial_l\big((\partial_kg_{ij})x^j\big)+\partial_lg_{ik}
    =(\partial_l\partial_kg_{ij})x^j+\partial_kg_{il}+\partial_lg_{ik},
  \]
  At $0$, three index permutations imply
  \[
    2\partial_kg_{ij}|_p=(\partial_kg_{ij}+\partial_jg_{ik})+(\partial_kg_{ij}+\partial_ig_{jk})
    -(\partial_ig_{jk}+\partial_jg_{ik})\Big|_p=0,
  \]
  With $C^2$ regularity, Taylor's theorem now gives
  $g_{ij}=\delta_{ij}+O(r^2)$. The formal proof packages this as a uniform bound
  on the second derivative over a smaller closed ball.
\end{proof}

\begin{remark}[asymptotics of the Hessian of $r$]
  \label{rem:pet-ch5-hessian-r-asymptotics}
  \lean{PetersenLib.hessianRadialFunction_asymptotics}
  \source{petersen:page-0211}
  \uses{lem:pet-ch5-normal-coords-first-order}
  Writing the metric in polar form $g=dr^2+g_r$ (restriction of $g$ to the level
  spheres of $r$) near $p$, and comparing with the Euclidean
  $\delta_{ij}=dr^2+r^2ds_{n-1}^2$, first-order agreement
  (\cref{lem:pet-ch5-normal-coords-first-order}) gives $\lim_{r\to0}g_r=0$ and
  $\lim_{r\to0}\big(\partial_rg_r-\partial_r(r^2ds_{n-1}^2)\big)=0$. Since
  $\partial_rg_r=2\operatorname{Hess}r$, this yields
  $\lim_{r\to0}\big(\operatorname{Hess}r-\tfrac1rg_r\big)=0$.
\end{remark}

\begin{theorem}[local rigidity of constant curvature]
  \label{thm:pet-ch5-riemann-constant-curvature-local}
  \lean{PetersenLib.constantCurvature_locallyIsometricToSpaceForm}
  \source{petersen:page-0211}
  \uses{rem:pet-ch5-hessian-r-asymptotics}
  \dcref{ch5:5.5.8}
  If a Riemannian $n$-manifold $(M,g)$ has constant sectional curvature $k$, every
  point of $M$ has a neighborhood isometric to an open subset of the space form
  $S^n_k$.
\end{theorem}
\begin{proof}
  \uses{thm:pet-ch5-riemann-constant-curvature-local}
  \sketch
  In exponential coordinates, constant curvature gives the radial Riccati equation
  $\nabla_{\partial_r}\operatorname{Hess}r+(\operatorname{Hess}r)^2=-kg_r$, with
  the asymptotics from \cref{rem:pet-ch5-hessian-r-asymptotics}.  The model
  tensor $(\mathrm{sn}_k'/\mathrm{sn}_k)g_r$ satisfies the same equation because
  \begin{align*}
    \nabla_{\partial_r}\!\left(\frac{\mathrm{sn}_k'}{\mathrm{sn}_k}g_r\right)
    +\left(\frac{\mathrm{sn}_k'}{\mathrm{sn}_k}\right)^2g_r
    &=\left(\frac{\mathrm{sn}_k'}{\mathrm{sn}_k}\right)'g_r
    +\left(\frac{\mathrm{sn}_k'}{\mathrm{sn}_k}\right)^2g_r
    +\frac{\mathrm{sn}_k'}{\mathrm{sn}_k}\nabla_{\partial_r}g_r\\
    &=-kg_r-\frac{\mathrm{sn}_k'}{\mathrm{sn}_k}\nabla_{\partial_r}dr^2=-kg_r,
  \end{align*}
  The prescribed initial behavior and ODE uniqueness yield
  $\operatorname{Hess}r=(\mathrm{sn}_k'/\mathrm{sn}_k)g_r$.  Consequently
  \[
    f_k(r)=\begin{cases}\tfrac12r^2,&k=0\\ \tfrac1k-\tfrac1k\mathrm{cs}_k(r),&k\ne0\end{cases}
  \]
  satisfies $\operatorname{Hess}f_k=(1-kf_k)g$, using
  $f_k'=\mathrm{sn}_k$ and $f_k''=\mathrm{cs}_k=1-kf_k$.  The local classification
  for this Hessian equation identifies the metric with the space form $S_k^n$.
\end{proof}

\begin{remark}[on the uniqueness used above]
  \label{rem:pet-ch5-uniqueness-remark-5-5-9}
  \lean{PetersenLib.remark_hessianODE_uniquenessSubtlety}
  \source{petersen:page-0212,petersen:page-0213}
  \uses{thm:pet-ch5-riemann-constant-curvature-local}
  \dcref{ch5:5.5.9}
  Neither of the systems $L_{\partial_r}g_r=2\operatorname{Hess}r$,
  $\nabla_{\partial_r}\operatorname{Hess}r+(\operatorname{Hess}r)^2=-kg_r$, nor the
  variant with $L_{\partial_r}\operatorname{Hess}r-(\operatorname{Hess}r)^2=-kg_r$,
  has unique solutions with $g_r,\operatorname{Hess}r\to0$ at $r=0$ in general — the
  trivial solution $g_r\equiv0$, $\operatorname{Hess}r\equiv0$ always solves them.
  Nor is it clear a priori that $\operatorname{Hess}r=\tfrac{\mathrm{sn}_k'}{\mathrm{sn}_k}g_r$
  solves the second system without knowing in advance that $L_{\partial_r}g_r=
  2\tfrac{\mathrm{sn}_k'}{\mathrm{sn}_k}g_r$; and the initial value problem
  $L_{\partial_r}g_r=2\tfrac{\mathrm{sn}_k'}{\mathrm{sn}_k}g_r$, $\lim_{r\to0}g_r=0$ has
  infinitely many solutions $\lambda\,\mathrm{sn}_k^2(r)\,ds_{n-1}^2$, $\lambda\in
  \mathbb R$, of which only one gives a smooth metric at $p$.
\end{remark}

\section{Riemannian Isometries}

\subsection{Local Isometries}

\begin{definition}[local Riemannian isometry]
  \label{def:pet-ch5-local-isometry}
  \lean{PetersenLib.IsLocalRiemannianIsometry}
  \source{petersen:page-0213}
  \leanok
  A map $F:(M,g_M)\to(N,g_N)$ is a \emph{local Riemannian isometry} if for each
  $p\in M$, $DF_p:T_pM\to T_{F(p)}N$ is a linear isometry. A local coordinate
  system $\varphi:U\to\Omega\subset\mathbb R^n$, with $U$ given the induced metric
  $g$ and $\Omega$ the coordinate representation $(\varphi^{-1})^*g=g_{ij}dx^idx^j$,
  is a trivial example.
\end{definition}

\begin{proposition}[basic properties of local isometries]
  \label{prop:pet-ch5-local-isometry-properties}
  \lean{PetersenLib.localIsometry_mapsGeodesicsToGeodesics,PetersenLib.localIsometry_isGeodesicOn,PetersenLib.localIsometry_expNaturality,PetersenLib.localIsometry_distanceDecreasing,PetersenLib.localIsometry_distancePreserving}
  \source{petersen:page-0213}
  \uses{def:pet-ch5-local-isometry, def:pet-ch5-geodesic, def:pet-ch5-exponential-map, def:pet-ch5-distance, lem:pet-ch5-global-uniqueness}
  \leanok
  \dcref{ch5:5.6.1}
  Let $F:(M,g_M)\to(N,g_N)$ be a local Riemannian isometry.
  \begin{enumerate}
  \item[(1)] $F$ maps geodesics to geodesics: if $\gamma$ is a continuous geodesic
  in $(M,g_M)$, then $F\circ\gamma$ is a geodesic in $(N,g_N)$; likewise on any
  open set of times on which $\gamma$ is a continuous geodesic.
  \item[(2)] The exponential map is natural: for $p\in M$ and $v\in T_pM$, the
  maximal existence domain of the geodesic initial-value problem $(p,v)$ is
  contained in that of $(F(p),DF_p(v))$, and on it
  \[
    F\circ c_{p,v}=c_{F(p),DF_p(v)},
  \]
  where $c_{q,w}$ denotes the maximal geodesic with initial datum $(q,w)$. Since
  $\exp_p(v)=c_{p,v}(1)$, the case $t=1$ is Petersen's
  $F(\exp_p(v))=\exp_{F(p)}(DF_p(v))$ whenever $\exp_p(v)$ is defined.
  \item[(3)] $F$ is distance decreasing.
  \item[(4)] If $F$ is also a bijection, $F$ is distance preserving.
  \end{enumerate}

  \emph{Formalization status.} All four parts are formalized. Two points of
  wording are forced by the formal definitions and are recorded here rather than
  smuggled into the statements. First, part~(1) carries a continuity hypothesis on
  $\gamma$: the formal \texttt{IsGeodesic} predicate constrains only the chart
  reading at the moving foot and so does not by itself imply that $\gamma$ is
  continuous, which is needed to place nearby points of the curve in a common
  chart. Continuity is automatic for Petersen's geodesics, which are smooth; the
  same hypothesis is already recorded explicitly by the formal
  initial-value predicate \texttt{IsGeodesicWithInitialOn} for this exact reason.
  Second, part~(2) is formalized at the intrinsic maximal geodesic rather than at
  the on-disk \texttt{expMap}; see
  \cref{rem:pet-ch5-exp-naturality-intrinsic}. Parts~(1) and~(2) are proved
  directly from the transformation law of the chart Christoffel symbols under an
  isometry; no Levi-Civita connection object occurs anywhere in the proof.
\end{proposition}
\begin{proof}
  \leanok
  \uses{prop:pet-ch5-local-isometry-properties}
  \sketch
  (1) In matched coordinates, the pullback metric has the same Christoffel symbols,
  so $F$ carries the geodesic equation to the target.

  (2) If $J$ is the maximal domain of $c_{p,v}$, then $F\circ c_{p,v}$ has initial
  data $(F(p),DF_pv)$.  Uniqueness
  (\cref{lem:pet-ch5-global-uniqueness}) identifies it with the target maximal
  geodesic on $J$, yielding the exponential identity at $t=1$.

  (3) Norm preservation by $DF$ gives $L(F\circ c)=L(c)$, hence infimizing gives
  $|F(p)F(q)|\le|pq|$.  If $F$ is bijective, apply this to $F^{-1}$ for equality.
\end{proof}

\begin{remark}[why naturality is stated at the intrinsic maximal geodesic]
  \label{rem:pet-ch5-exp-naturality-intrinsic}
  \uses{prop:pet-ch5-local-isometry-properties, rem:pet-ch5-injectivity-radius-chart-artifact}
  Part~(2) of \cref{prop:pet-ch5-local-isometry-properties} is Petersen's
  $F(\exp_p(v))=\exp_{F(p)}(DF_p(v))$, and the formalization at the maximal
  geodesic is faithful to the book: $\exp_p(v)$ \emph{is} the maximal geodesic
  through $(p,v)$ evaluated at time $1$, so the intrinsic form contains the book's
  statement as its $t=1$ case and is in fact stronger, giving the identity at every
  admissible time.

  The formalization does \emph{not} use the on-disk \texttt{PetersenLib.expMap} and
  \texttt{expDomain}, and this is forced rather than a convenience: at those objects
  the statement is \emph{false}, not merely unproven. They are anchored to the chart
  at the initial point, so $\exp_p$-geodesics cannot leave the chart at $p$ while
  $\exp_{F(p)}$-geodesics cannot leave the chart at $F(p)$; a local isometry need
  not match the two atlases, and the inclusion of existence domains fails. This is
  the defect recorded in \cref{rem:pet-ch5-injectivity-radius-chart-artifact}, and
  it is a defect of the on-disk exponential — i.e.\ of
  \cref{def:pet-ch5-exponential-map}, where it is already flagged — not of the
  naturality statement. The formal statement additionally assumes $M$ and $N$
  Hausdorff, which the uniqueness of maximal geodesics requires and which is part of
  Petersen's standing conventions.
\end{remark}

\begin{proposition}[uniqueness of Riemannian isometries]
  \label{prop:pet-ch5-isometry-uniqueness}
  \lean{PetersenLib.localIsometry_uniquelyDeterminedByOnePoint}
  \source{petersen:page-0214}
  \uses{def:pet-ch5-local-isometry, prop:pet-ch5-local-isometry-properties, cor:pet-ch5-uniform-injectivity-radius-compact}
  \leanok
  \dcref{ch5:5.6.2}
  Let $F,G:(M,g_M)\to(N,g_N)$ be local Riemannian isometries, $M$ connected. If
  $F(p)=G(p)$ and $DF_p=DG_p$ for some $p\in M$, then $F=G$ on $M$.
\end{proposition}
\begin{proof}
  \leanok
  \uses{prop:pet-ch5-isometry-uniqueness,prop:pet-ch5-local-isometry-properties,cor:pet-ch5-uniform-injectivity-radius-compact}
  \sketch
  Let $A=\{x\mid F(x)=G(x),\ DF_x=DG_x\}$.  It is closed and contains $p$.  If
  $x\in A$, part (2) of \cref{prop:pet-ch5-local-isometry-properties} gives
  \[
    F\circ\exp_x(v)=\exp_{F(x)}\circ DF_x(v)=\exp_{G(x)}\circ DG_x(v)=G\circ\exp_x(v),
  \]
  The image of a small ball under $\exp_x$ is a neighborhood by
  \cref{cor:pet-ch5-uniform-injectivity-radius-compact}, so $F=G$ near $x$ and their
  differentials agree there.  Thus $A$ is open, and connectedness gives $A=M$.
\end{proof}

\begin{proposition}[completeness is preserved by Riemannian coverings]
  \label{prop:pet-ch5-covering-completeness}
  \lean{PetersenLib.riemannianCovering_completenessEquivalence}
  \source{petersen:page-0214}
  \uses{def:pet-ch5-local-isometry, def:pet-ch5-geodesically-complete}
  \leanok
  \dcref{ch5:5.6.3}
  Let $F:(M,g_M)\to(N,g_N)$ be a Riemannian covering map. $(M,g_M)$ is geodesically
  complete iff $(N,g_N)$ is.
\end{proposition}
\begin{proof}
  \leanok
  \sketch
  Geodesics lift uniquely through a covering, and a local isometry preserves the
  geodesic equation.  Thus completeness of $N$ extends every lifted geodesic to all
  time, while completeness of $M$ pushes each all-time lift down to an extension in
  $N$.
\end{proof}

\begin{lemma}[completeness forces a covering map]
  \label{lem:pet-ch5-local-isometry-covering}
  \lean{PetersenLib.completeLocalIsometry_isCoveringMap}
  \source{petersen:page-0214}
  \uses{def:pet-ch5-local-isometry, def:pet-ch5-geodesically-complete, cor:pet-ch5-uniform-injectivity-radius-compact}
  \dcref{ch5:5.6.4}
  If $F:(M,g_M)\to(N,g_N)$ is a local Riemannian isometry, $M$ geodesically
  complete, $N$ connected, then $F$ is a Riemannian covering map.
\end{lemma}
\begin{proof}
  \uses{lem:pet-ch5-local-isometry-covering,prop:pet-ch5-local-isometry-properties}
  \sketch
  Choose $\varepsilon$ with $\exp_q:B(0,\varepsilon)\to B(q,\varepsilon)$ a
  diffeomorphism. Completeness defines the corresponding exponentials at each
  $p\in F^{-1}(q)$, and \cref{prop:pet-ch5-local-isometry-properties} gives
  $F\circ\exp_p(v)=\exp_q(DF_p(v))$ for $v\in B(0,\varepsilon)$; as $\exp_q$ and
  Hence $F:B(p,\varepsilon)\to B(q,\varepsilon)$ is a diffeomorphism.

  Conversely, lift the unique radial segment from $q$ to $F(x)$ backwards from
  $x$. Completeness and \cref{thm:pet-ch5-local-uniqueness} put $x$ in a unique
  such ball, so these balls evenly cover $F^{-1}(B(q,\varepsilon))$.

  The image is open.  It is closed because
  \cref{cor:pet-ch5-uniform-injectivity-radius-compact} gives one radius for a
  convergent image sequence and its limit. Connectedness of $N$ gives surjectivity.
\end{proof}

\begin{definition}[fixed-point set; totally geodesic submanifold]
  \label{def:pet-ch5-fixed-point-set}
  \lean{PetersenLib.fixedPointSet,PetersenLib.IsTotallyGeodesic}
  \source{petersen:page-0215}
  \uses{def:pet-ch5-exponential-map}
  \leanok
  For a set $S$ of isometries of $(M,g)$, $\operatorname{Fix}(S)=\{x\in M\mid
  F(x)=x\ \forall F\in S\}$. A submanifold $N\subset(M,g)$ is \emph{totally
  geodesic} if for each $p\in N$ a neighborhood of $0\in T_pN\subset T_pM$ is
  mapped into $N$ by $\exp_p$ (for $M$): geodesics of $N$ are geodesics of $M$, and
  any $M$-geodesic tangent to $N$ at a point stays in $N$ for a short time.
\end{definition}

\begin{definition}[the fixed tangent subspace]
  \label{def:pet-ch5-fixed-tangent-subspace}
  \lean{PetersenLib.fixedTangentSubspace}
  \source{petersen:page-0215}
  \uses{def:pet-ch5-fixed-point-set}
  \leanok
  For a set $S$ of isometries of $(M,g)$ and $p\in M$,
  \[
    V_p=\{v\in T_pM\mid DF_p(v)=v\ \text{ for all }F\in S\},
  \]
  the intersection of the $+1$-eigenspaces of the $DF_p$, $F\in S$. It is a linear
  subspace of $T_pM$.
\end{definition}

\begin{proposition}[fixed-point components are totally geodesic]
  \label{prop:pet-ch5-fixed-point-component-totally-geodesic}
  \lean{PetersenLib.fixedPointSetComponent_totallyGeodesic}
  \source{petersen:page-0215}
  \uses{def:pet-ch5-fixed-point-set, def:pet-ch5-fixed-tangent-subspace, lem:pet-ch5-short-geodesics-segments-normal, prop:pet-ch5-local-isometry-properties}
  \leanok
  Let $(M,g)$ be a connected Riemannian manifold whose Riemannian distance induces
  its topology, and let $S\subset\operatorname{Iso}(M,g)$. Then every connected
  component of $\operatorname{Fix}(S)$ is totally geodesic, with tangent
  distribution $V_p$ as in \cref{def:pet-ch5-fixed-tangent-subspace}.

  \emph{Formalization status.} This is the totally-geodesic half of
  \cref{prop:pet-ch5-fixed-point-totally-geodesic}, which is the half that is
  formalized; the submanifold half is not. Being totally geodesic is a chart-free
  property and does not presuppose submanifold structure. The formal statement adds
  two hypotheses to Petersen's: that the Riemannian distance induces the topology
  of $M$ (the chapter's standing metric/length compatibility assumption, under
  which the short-segment rigidity of
  \cref{lem:pet-ch5-short-geodesics-segments-normal} is stated), and that $M$ is
  connected (so that the infimum defining the distance runs over a nonempty set of
  competitors). The statement quantifies over the intrinsic ball
  $\{v\mid|v|_g<\varepsilon\}$ and does not mention the on-disk exponential's
  domain, so the chart artifact of
  \cref{rem:pet-ch5-injectivity-radius-chart-artifact} does not touch it.
\end{proposition}
\begin{proof}
  \leanok
  \uses{def:pet-ch5-fixed-tangent-subspace,lem:pet-ch5-short-geodesics-segments-normal,prop:pet-ch5-local-isometry-properties}
  \sketch
  For $p\in\operatorname{Fix}(S)$ and $v\in V_p$, naturality from
  \cref{prop:pet-ch5-local-isometry-properties} gives, for every $F\in S$,
  $F(\exp_p(tv))=\exp_{F(p)}(tDF_p(v))=\exp_p(tv)$ for every $t$ for which
  Thus the radial geodesic is fixed pointwise.  By
  \cref{lem:pet-ch5-short-geodesics-segments-normal}, small such rays are defined
  and continuous, hence remain in the connected component of $p$.  Therefore
  $\exp_p$ carries a neighborhood of $0\in V_p$ into that component.
\end{proof}

\begin{proposition}[fixed-point sets are totally geodesic]
  \label{prop:pet-ch5-fixed-point-totally-geodesic}
  \source{petersen:page-0215}
  \uses{def:pet-ch5-fixed-point-set, def:pet-ch5-fixed-tangent-subspace, prop:pet-ch5-fixed-point-component-totally-geodesic, cor:pet-ch5-metric-ball-image}
  \dcref{ch5:5.6.5}
  If $S\subset\operatorname{Iso}(M,g)$, each connected component of
  $\operatorname{Fix}(S)$ is a totally geodesic submanifold.

  \emph{Formalization status.} Only the totally-geodesic half is formalized, as
  \cref{prop:pet-ch5-fixed-point-component-totally-geodesic}. The submanifold half
  — that $\exp_p$ restricts to a bijection $V_p\cap B(0,\varepsilon)\to
  \operatorname{Fix}(S)\cap B(p,\varepsilon)$, exhibiting $\operatorname{Fix}(S)$
  near $p$ as the image of a linear subspace and so furnishing a submanifold chart
  — is not formalized, and nothing on disk asserts that
  $\operatorname{Fix}(S)$ or any component of it is a submanifold.
\end{proposition}
\begin{proof}
  \uses{prop:pet-ch5-fixed-point-totally-geodesic,prop:pet-ch5-fixed-point-component-totally-geodesic}
  \sketch
  Proposition \cref{prop:pet-ch5-fixed-point-component-totally-geodesic} gives
  $\exp_p(V_p)\subset\operatorname{Fix}(S)$ locally.  Choose a normal ball using
  \cref{def:pet-ch5-injectivity-radius}.  If a fixed point $q$ lies in it, each
  $F\in S$ carries the unique short segment from $p$ to $q$ to another such segment;
  uniqueness from \cref{cor:pet-ch5-metric-ball-image} fixes the segment pointwise.
  Consequently
  $\exp_p:V_p\cap B(0,\varepsilon)\to\operatorname{Fix}(S)\cap B(p,\varepsilon)$
  is a bijection, furnishing the required submanifold chart.
\end{proof}

\begin{remark}[even codimension of orientation-preserving fixed sets]
  \label{rem:pet-ch5-fixed-point-codimension}
  \lean{PetersenLib.orientationPreservingFixedSet_evenCodimension}
  \source{petersen:page-0216}
  \uses{prop:pet-ch5-fixed-point-totally-geodesic}
  \dcref{ch5:5.6.6}
  If $DF_p:T_pM\to T_pM$ is orientation preserving for $p\in\operatorname{Fix}(F)$,
  the Zariski tangent space at $p$ has even codimension, since the $+1$-eigenspace
  of an element of $SO(n)$ has even codimension; so each component of
  $\operatorname{Fix}(F)$ has even codimension.
\end{remark}

\subsection{Constant Curvature Revisited}

\begin{theorem}[the monodromy map]
  \label{thm:pet-ch5-monodromy-theorem}
  \lean{PetersenLib.monodromyMap}
  \source{petersen:page-0216}
  \uses{thm:pet-ch5-riemann-constant-curvature-local, prop:pet-ch5-isometry-uniqueness}
  \dcref{ch5:5.6.7}
  Let $(M,g)$ be simply connected of dimension $n$ and constant curvature $k$, and
  $L:T_pM\to T_qS^n_k$ a linear isometry. There is a unique local Riemannian
  isometry $F:M\to S^n_k$ (the \emph{monodromy map}) with $DF_p=L$. If $(M,g)$ is
  geodesically complete, $F$ is a diffeomorphism.
\end{theorem}
\begin{proof}
  \uses{thm:pet-ch5-monodromy-theorem,lem:pet-ch5-local-isometry-covering,prop:pet-ch5-isometry-uniqueness}
  \sketch
  By \cref{thm:pet-ch5-riemann-constant-curvature-local}, sufficiently small convex
  balls admit local isometries to $S_k^n$.  Continue the prescribed germ along a
  path by a chain of such balls; on connected overlaps the extensions agree by
  \cref{prop:pet-ch5-isometry-uniqueness}.  The endpoint value is locally constant
  under path deformations, hence depends only on the homotopy class.  Simple
  connectedness makes the continuation path independent and defines the required
  local isometry on all of $M$.  Its germ is unique by the same isometry-uniqueness
  result.  If $M$ is complete, \cref{lem:pet-ch5-local-isometry-covering} makes this
  map a covering of the simply connected space form, hence a diffeomorphism.
\end{proof}

\begin{example}[immersions realize monodromy maps]
  \label{ex:pet-ch5-immersion-monodromy}
  \lean{PetersenLib.example_immersion_monodromy}
  \source{petersen:page-0216}
  \uses{thm:pet-ch5-monodromy-theorem}
  \dcref{ch5:5.6.8}
  For an immersion $M^n\to S^n_k$ with pulled-back metric, the monodromy map $F$ of
  \cref{thm:pet-ch5-monodromy-theorem} is (a version of) this immersion; such maps
  can fold wildly for $n\ge2$ and need not resemble covering maps.
\end{example}

\begin{example}[a diffeomorphism of a punctured space form]
  \label{ex:pet-ch5-folding-immersion}
  \lean{PetersenLib.example_puncturedSpaceForm_diffeomorphism}
  \source{petersen:page-0217}
  \uses{thm:pet-ch5-monodromy-theorem}
  \dcref{ch5:5.6.9}
  If $U\subset S^n_k$ is an open disc with smooth boundary, there is a
  diffeomorphism $F:M=S^n_k-\{p\}\to S^n_k-U$; the pulled-back metric on $M$ has
  constant curvature but looks poor near the missing point.
\end{example}

\begin{example}[$\mathbb{RP}^n$ admits no immersion into $S^n$]
  \label{ex:pet-ch5-rp-n-no-immersion}
  \lean{PetersenLib.example_projectiveSpace_noImmersionIntoSphere}
  \source{petersen:page-0217}
  \uses{thm:pet-ch5-monodromy-theorem}
  \dcref{ch5:5.6.10}
  $M=\mathbb{RP}^n=(\mathbb R^n-\{0\})/(\text{antipodal map})$ is not simply
  connected, hence admits no immersion into $S^n$ realizing the monodromy
  construction.
\end{example}

\begin{example}[monodromy map of a universal cover]
  \label{ex:pet-ch5-universal-cover-monodromy}
  \lean{PetersenLib.example_universalCover_monodromy}
  \source{petersen:page-0217}
  \uses{thm:pet-ch5-monodromy-theorem}
  \dcref{ch5:5.6.11}
  If $M$ is the universal cover of $S^2-\{\pm p\}$, the monodromy map is not
  one-to-one; it is exactly the covering map $M\to S^2-\{\pm p\}$.
\end{example}

\begin{corollary}[closed simply connected constant curvature]
  \label{cor:pet-ch5-closed-simply-connected-constant-curvature}
  \lean{PetersenLib.closedSimplyConnectedConstantCurvature_isSphere}
  \source{petersen:page-0217}
  \uses{thm:pet-ch5-monodromy-theorem}
  \dcref{ch5:5.6.12}
  If $M$ is closed and simply connected with constant curvature $k$, then $k>0$
  and $M=S^n$; in particular $S^p\times S^q$ and $\mathbb{CP}^n$ admit no
  constant-curvature metric.
\end{corollary}
\begin{proof}
  \uses{cor:pet-ch5-closed-simply-connected-constant-curvature}
  \sketch
  Compactness gives completeness by \cref{cor:pet-ch5-compact-manifold-complete},
  so \cref{thm:pet-ch5-monodromy-theorem} identifies $M$ with $S_k^n$.  This space
  is compact only for $k>0$, when it is a rescaled round sphere.
\end{proof}

\begin{corollary}[complete noncompact constant curvature]
  \label{cor:pet-ch5-noncompact-constant-curvature}
  \lean{PetersenLib.completeNoncompactConstantCurvature_universalCoverIsRn}
  \source{petersen:page-0217}
  \uses{thm:pet-ch5-monodromy-theorem}
  \dcref{ch5:5.6.13}
  If $M$ is geodesically complete, noncompact, with constant curvature $k$, then
  $k\le0$ and the universal cover of $M$ is diffeomorphic to $\mathbb R^n$; in
  particular $S^2\times\mathbb R^2$ and $S^n\times\mathbb R$ admit no geodesically
  complete constant-curvature metric.
\end{corollary}
\begin{proof}
  \uses{cor:pet-ch5-noncompact-constant-curvature}
  \sketch
  By \cref{prop:pet-ch5-covering-completeness}, the universal cover is complete.
  The monodromy theorem \cref{thm:pet-ch5-monodromy-theorem} identifies it with
  $S_k^n$.  Positive $k$ would make this cover, and hence its quotient $M$, compact.
  Thus $k\le0$, and $S_k^n$ is diffeomorphic to $\mathbb R^n$.
\end{proof}

\begin{corollary}[classification of complete constant curvature spaces]
  \label{cor:pet-ch5-classification-constant-curvature}
  \lean{PetersenLib.killingHopfClassification}
  \source{petersen:page-0218}
  \uses{thm:pet-ch5-monodromy-theorem, prop:pet-ch5-covering-completeness}
  \dcref{ch5:5.6.14}
  If $(M,g)$ is connected, geodesically complete, with constant curvature $k$, its
  universal cover is isometric to $S^n_k$.
\end{corollary}
\begin{proof}
  \uses{cor:pet-ch5-classification-constant-curvature}
  \sketch
  The universal cover remains complete by
  \cref{prop:pet-ch5-covering-completeness}.  The map supplied by
  \cref{thm:pet-ch5-monodromy-theorem} is both a local isometry and a
  diffeomorphism onto $S_k^n$, hence a global isometry.
\end{proof}

\subsection{Metric Characterization of Maps}

\begin{remark}[recovering the metric from the distance function]
  \label{rem:pet-ch5-polarization-metric-recovery}
  \lean{PetersenLib.metricRecoveredFromDistance}
  \source{petersen:page-0218,petersen:page-0219}
  \uses{def:pet-ch5-distance}
  By polarization, $g(v,w)=\tfrac12(|v+w|^2-|v|^2-|w|^2)$, so it suffices to
  recover $|v|$ from $|\cdot|_g$; for a curve $\alpha$ with $\dot\alpha(0)=v$,
  $|v|=\lim_{t\to0}|\alpha(t)\alpha(0)|/t$. Thus $g$ is determined by
  $(M,|\cdot|_g)$ together with the differentiable structure, motivating synthetic
  (metric-space-only) characterizations of Riemannian maps.
\end{remark}

\begin{definition}[distance coordinates]
  \label{def:pet-ch5-distance-coordinates}
  \lean{PetersenLib.distanceCoordinates}
  \source{petersen:page-0219}
  \uses{cor:pet-ch5-uniform-injectivity-radius-compact}
  \leanok
  Fix $p\in M$ and $U\ni p$ with $B(q,\operatorname{inj}(q))\supset U$ for all
  $q\in U$; then $r_q(x)=|xq|$ is smooth on $U-\{q\}$ for each $q\in U$. Choosing
  $q_1,\dots,q_n\in U-\{p\}$, $n=\dim M$, with $\nabla r_{q_1}(p),\dots,\nabla
  r_{q_n}(p)$ linearly independent, the inverse function theorem makes
  $\varphi=(r_{q_1},\dots,r_{q_n})$ coordinates near $p$.
\end{definition}

\begin{theorem}[distance-preserving bijections are isometries]
  \label{thm:pet-ch5-myers-steenrod-isometry}
  \lean{PetersenLib.distancePreservingBijection_isRiemannianIsometry}
  \source{petersen:page-0219}
  \uses{def:pet-ch5-distance-coordinates, def:pet-ch5-local-isometry}
  \dcref{ch5:5.6.15}
  If $F:M\to N$ is a bijection between Riemannian manifolds with
  $|F(p)F(q)|_{g_N}=|pq|_{g_M}$ for all $p,q\in M$, then $F$ is a Riemannian
  isometry.
\end{theorem}
\begin{proof}
  \uses{thm:pet-ch5-myers-steenrod-isometry}
  \sketch
  Choose distance coordinates at $q=F(p)$ as in
  \cref{def:pet-ch5-distance-coordinates}, with poles $q_i=F(p_i)$.  Then
  $r_{q_i}\circ F(x)=|xp_i|$, which is smooth near $p$; hence $F$ is smooth in
  these coordinates.  For $v\in T_pM$, take the short segment
  $c(t)=\exp_p(tv)$.  Distance preservation gives
  \[
    |t-s|\cdot|v|=|c(t)c(s)|_{g_M}=|F(c(t))F(c(s))|_{g_N},
  \]
  Therefore
  \[
    |DF(v)|=\left|\frac{d(F\circ c)}{dt}\right|_{t=0}
    =\lim_{t\to0}\frac{|F(c(t))F(c(0))|_{g_N}}{|t|}
    =\lim_{t\to0}\frac{|c(t)c(0)|_{g_M}}{|t|}=|\dot c(0)|=|v|.
  \]
\end{proof}

\begin{definition}[submetry]
  \label{def:pet-ch5-submetry}
  \lean{PetersenLib.IsSubmetry}
  \source{petersen:page-0220}
  \uses{def:pet-ch5-metric-balls}
  \leanok
  $F:(\tilde M,g_{\tilde M})\to(M,g_M)$ is a \emph{submetry} if for every
  $\tilde p\in\tilde M$ there is $r>0$ such that $F(B(\tilde p,\varepsilon))=
  B(F(\tilde p),\varepsilon)$ for all $\varepsilon\le r$. Submetries are locally
  distance nonincreasing, hence continuous, and compositions of submetries (or of
  Riemannian submersions) are again submetries (resp.\ Riemannian submersions).
\end{definition}

\begin{theorem}[surjective submetries are $C^1$ submersions]
  \label{thm:pet-ch5-berestovskii-submetry}
  \lean{PetersenLib.surjectiveSubmetry_isC1RiemannianSubmersion}
  \source{petersen:page-0220}
  \uses{def:pet-ch5-submetry, def:pet-ch5-distance-coordinates}
  \dcref{ch5:5.6.16}
  If $F:(\tilde M,g_{\tilde M})\to(M,g_M)$ is a surjective submetry, then $F$ is a
  $C^1$ Riemannian submersion.
\end{theorem}
\begin{proof}
  \uses{thm:pet-ch5-berestovskii-submetry}
  \sketch
  Choose radii below the injectivity radii at $p$ and $\tilde p$.  The ball-lifting
  property supplies a unique closest lift $\bar q\in F^{-1}(q)$ with
  $|\tilde p\bar q|=|pq|$, continuously in $\tilde p$; normalized short segments
  therefore define continuous horizontal lifts.  In distance coordinates from
  \cref{def:pet-ch5-distance-coordinates}, each component $r_i\circ F$ is a scalar
  submetry.  Its gradient is the continuous horizontal lift of $\partial_{r_i}$,
  proving $C^1$ regularity and that $DF$ is an isometry on the horizontal space.
\end{proof}

\begin{remark}[submetries are $C^{1,1}$, sharply]
  \label{rem:pet-ch5-submetry-c11}
  \lean{PetersenLib.remark_submetriesAreC11Sharp}
  \source{petersen:page-0220,petersen:page-0221}
  \uses{thm:pet-ch5-berestovskii-submetry}
  \dcref{ch5:5.6.17}
  Submetries are in fact $C^{1,1}$ (locally Lipschitz derivative), via local
  Lipschitz continuity of $\tilde p\mapsto\bar q$; this cannot be improved in
  general, e.g.\ for $K=[0,1]^2\subset\mathbb R^2$ the level sets of $r(x)=|xK|$ are
  rounded squares with quarter-circle corners, not $C^2$.
\end{remark}

\subsection{The Slice Theorem}

\begin{definition}[proper action; orbit; isotropy group]
  \label{def:pet-ch5-proper-action}
  \lean{PetersenLib.IsProperAction,PetersenLib.orbit,PetersenLib.actionIsotropyGroup}
  \source{petersen:page-0221}
  \leanok
  An action of a topological group $H$ on $M$ is \emph{proper} if
  $H\times M\to M\times M$, $(h,p)\mapsto(hp,p)$, is proper. The \emph{orbit}
  through $p$ is $Hp=\{hp\mid h\in H\}$; $H\backslash M$ carries the quotient
  topology, making $M\to H\backslash M$ continuous and open, second countable, and
  Hausdorff when the action is proper. The \emph{isotropy group} at $p$ is
  $H_p=\{h\in H\mid hp=p\}$; along an orbit $H_{hp}=hH_ph^{-1}$, $H_p$ is compact
  for proper actions, and $H/H_p\to Hp$ is a homeomorphism. $H$ acts \emph{freely}
  if $H_p=\{e\}$ for all $p$.
\end{definition}

\begin{example}[a nonproper isometric action]
  \label{ex:pet-ch5-nonproper-action}
  \lean{PetersenLib.example_nonproperIsometricAction}
  \source{petersen:page-0221}
  \uses{def:pet-ch5-proper-action}
  \dcref{ch5:5.6.18}
  $\operatorname{Iso}(M)$ acts properly on $M$ (Arzelà–Ascoli), but a subgroup
  $H\subset\operatorname{Iso}(M)$ need not act properly unless closed: the action
  $\theta\cdot(z_1,z_2)=(e^{i\theta}z_1,e^{i\alpha\theta}z_2)$ of $\mathbb R$ on
  $S^1\times S^1$ is proper iff $\alpha$ is rational.
\end{example}

\begin{theorem}[orbits are submanifolds; $\operatorname{Iso}(M,g)$ is a Lie group]
  \label{thm:pet-ch5-myers-steenrod-lie-group}
  \lean{PetersenLib.closedIsometrySubgroup_orbitsAreSubmanifolds,PetersenLib.isometryGroup_isLieGroup}
  \source{petersen:page-0221}
  \uses{def:pet-ch5-proper-action, cor:pet-ch5-tubular-neighborhood}
  \dcref{ch5:5.6.19}
  If $H$ is a closed subgroup of $\operatorname{Iso}(M,g)$, the orbits of the
  action are submanifolds; in particular $\operatorname{Iso}(M,g)$ itself is a Lie
  group.
\end{theorem}
\begin{proof}
  \uses{thm:pet-ch5-myers-steenrod-lie-group,cor:pet-ch5-tubular-neighborhood,cor:pet-ch5-uniform-injectivity-radius-compact,prop:pet-ch5-exp-diffeomorphism-properties}
  \sketch
  Exponential coordinates from
  \cref{prop:pet-ch5-exp-diffeomorphism-properties} identify limits of short orbit
  segments with a cone $T_pHp$.  Reparametrization and composition by nearby group
  elements show that this cone is a linear subspace, that
  $T_{hp}Hp=Dh(T_pHp)$, and that it varies continuously along the orbit.
  Minimizing segments from nearby points to the orbit show that the corresponding
  normal directions are orthogonal to this subspace.

  Choose a small submanifold $N$ tangent to $T_pHp$ and tubular coordinates
  $N\times B$ from \cref{cor:pet-ch5-tubular-neighborhood}.  Transversality makes
  projection of the orbit to $N$ locally one-to-one.  If its image had a boundary,
  a smooth domain in the complement would touch the orbit with tangent hyperplane
  containing $T_pHp$; dimension counting would force a nonzero intersection with
  the nearly normal fibre, contradicting transversality.  Thus the orbit is locally
  a continuous graph over $N$.  Continuous tangent spaces identify its derivative,
  so this graph is $C^1$.

  The full isometry group acts freely on the open set of $(n+1)$-tuples whose
  initial directions are independent; its orbit realizes the group as a $C^1$
  submanifold.  The identities $T_{hp}Hp=Dh(T_pHp)$ and
  $Dh=\exp_{hp}^{-1}\circ h\circ\exp_p$ bootstrap $C^k$ regularity of the action and
  orbits to $C^{k+1}$.  Iteration gives a smooth Lie group and smooth orbits for
  every closed subgroup.
\end{proof}

\begin{proposition}[kernel of the orbit map]
  \label{prop:pet-ch5-orbit-map-differential}
  \lean{PetersenLib.orbitMapDifferentialKernel}
  \source{petersen:page-0224}
  \uses{thm:pet-ch5-myers-steenrod-lie-group}
  \dcref{ch5:5.6.20}
  Let $H$ be a closed subgroup of $\operatorname{Iso}(M,g)$ with Lie algebra
  $\mathfrak h=T_eH$, and $\mathfrak h_p=\{v\in\mathfrak h\mid\exp(tv)\in H_p\ \forall t\}$
  the subalgebra of $H_p$. For the orbit map $O_p(h)=hp$,
  $\ker\big((DO_p)|_v\big)=\mathfrak h_p$, and more generally
  $\ker\big((DO_p)|_s\big)=DL_s(\mathfrak h_p)$.
\end{proposition}
\begin{proof}
  \uses{prop:pet-ch5-orbit-map-differential}
  \sketch
  The chain rule and $(h_1j)p=h_1(jp)$ reduce to $s=e$. There,
  $(DO_p)|_v(v)=\frac{d}{dt}(\exp(tv)p)|_{t=0}$ and
  \[
    \frac{d}{dt}(\exp(tv)\cdot p)\Big|_{t=t_0}
    =D(\exp(t_0v))\left(\frac{d}{ds}(\exp(sv)\cdot p)\Big|_{s=0}\right).
  \]
  If the derivative at $0$ vanishes, the displayed identity makes it vanish for
  every $t$, so $\exp(tv)p=p$ and $v\in\mathfrak h_p$. The converse is immediate.
\end{proof}

\begin{theorem}[the free slice theorem]
  \label{thm:pet-ch5-free-slice-theorem}
  \lean{PetersenLib.freeSliceTheorem}
  \source{petersen:page-0224}
  \uses{def:pet-ch5-proper-action, cor:pet-ch5-tubular-neighborhood, thm:pet-ch5-myers-steenrod-lie-group}
  \dcref{ch5:5.6.21}
  If $H\subset\operatorname{Iso}(M,g)$ is closed and acts freely, $H\backslash M$
  has a smooth manifold structure and Riemannian metric making $M\to H\backslash M$
  a Riemannian submersion.
\end{theorem}
\begin{proof}
  \uses{thm:pet-ch5-free-slice-theorem}
  \sketch
  By \cref{thm:pet-ch5-myers-steenrod-lie-group} and
  \cref{prop:pet-ch5-orbit-map-differential}, a free orbit is a properly embedded
  copy of $H$.  The map $(h,v)\mapsto h\exp_p(v)$ identifies
  $H\times T_p^\perp Hp$ isometrically with the normal bundle, and
  \cref{cor:pet-ch5-tubular-neighborhood} supplies a tube near the zero section.

  Shrink to a uniform radius for which the nearest orbit point is unique.  The first
  variation formula makes every minimizing segment to the orbit normal, proving
  injectivity of $\exp^\perp$ on $H\times B(0,\epsilon)$.  Equivariance gives
  nonsingularity, and proper embeddedness of $Hp$ gives properness.

  Slice balls therefore define smooth quotient charts. Identifying
  $T_{Hp}(H\backslash M)$ with $T_p^\perp Hp$ defines a metric independent of the
  representative because $Dh$ is an isometry, and the quotient map is a Riemannian
  submersion.
\end{proof}

\begin{remark}[homogeneous quotients $H/K$]
  \label{rem:pet-ch5-quotient-homogeneous}
  \lean{PetersenLib.remark_homogeneousQuotient}
  \source{petersen:page-0225}
  \uses{thm:pet-ch5-free-slice-theorem}
  \dcref{ch5:5.6.22}
  For $K\subset H$ compact, averaging any metric on $H$ over $K$ (acting by
  $(k,x)\mapsto xk^{-1}$) makes it right-$K$-invariant, giving a free isometric
  action to which \cref{thm:pet-ch5-free-slice-theorem} applies, so $H/K$ is a
  manifold with a Riemannian submersion metric from $H$; if the metric on $H$ is
  also left-invariant, $H$ acts by isometries on $H/K$, making it homogeneous.
\end{remark}

\begin{remark}[closed subgroups and right-invariant metrics]
  \label{rem:pet-ch5-closed-subgroup-quotient}
  \lean{PetersenLib.remark_closedSubgroupRightInvariantQuotient}
  \source{petersen:page-0226}
  \uses{rem:pet-ch5-quotient-homogeneous}
  For $K\subset H$ merely closed, right multiplication by $K$ is still a proper
  action, made isometric by a right-$K$-invariant metric on $H$; but it may not be
  possible to simultaneously keep the metric on $H$ left-invariant, so $H$ need not
  act by isometries on $H/K$ in this more general case.
\end{remark}

\begin{corollary}[orbits of proper isometric actions are properly embedded]
  \label{cor:pet-ch5-proper-action-orbit-embedding}
  \lean{PetersenLib.properIsometricAction_orbitProperlyEmbedded}
  \source{petersen:page-0226}
  \uses{prop:pet-ch5-orbit-map-differential, def:pet-ch5-proper-action}
  \dcref{ch5:5.6.23}
  For a proper isometric action $H\times M\to M$ and $p\in M$, the orbit map
  $H/H_p\to Hp$ is a proper embedding.
\end{corollary}
\begin{proof}
  \uses{cor:pet-ch5-proper-action-orbit-embedding}
  \sketch
  The orbit map is proper and injective.  By
  \cref{prop:pet-ch5-orbit-map-differential}, quotienting by $\mathfrak h_p$ makes
  its differential injective.  It is therefore a proper injective immersion, hence
  a proper embedding.
\end{proof}

\begin{definition}[slice representation]
  \label{def:pet-ch5-slice-representation}
  \lean{PetersenLib.sliceRepresentation}
  \source{petersen:page-0226}
  \uses{cor:pet-ch5-proper-action-orbit-embedding}
  The \emph{slice representation} of a proper isometric action $H\times M\to M$ at
  $p$ is the linear representation $H_p\times T_p^\perp Hp\to T_p^\perp Hp$,
  $(h,v)\mapsto Dh|_p(v)$. Letting $H_p$ act on $H$ on the right, it also acts on
  $H\times T_p^\perp Hp$, and by \cref{cor:pet-ch5-proper-action-orbit-embedding}
  the quotient $H\times_{H_p}T_p^\perp Hp$ has a natural manifold structure.
\end{definition}

\begin{theorem}[the slice theorem]
  \label{thm:pet-ch5-slice-theorem}
  \lean{PetersenLib.sliceTheorem}
  \source{petersen:page-0226}
  \uses{def:pet-ch5-slice-representation, thm:pet-ch5-free-slice-theorem}
  \dcref{ch5:5.6.24}
  Let $H\subset\operatorname{Iso}(M,g)$ be closed. For $p\in M$ there is a map
  $\exp^\perp:H\times_{H_p}T_p^\perp Hp\to M$ that is a diffeomorphism on a uniform
  tubular neighborhood $H\times_{H_p}B(0,\epsilon)$ onto an $\epsilon$-neighborhood
  of $Hp$.
\end{theorem}
\begin{proof}
  \uses{thm:pet-ch5-slice-theorem}
  \sketch
  By \cref{cor:pet-ch5-proper-action-orbit-embedding}, the orbit is properly
  embedded. The equivariant map $(h,v)\mapsto Dh|_p(v)$ descends to
  $H\times_{H_p}T_p^\perp Hp\cong T^\perp Hp$ because
  $D(hk^{-1})Dk(v)=Dh(v)$ for $k\in H_p$. The tubular construction in
  \cref{thm:pet-ch5-free-slice-theorem} then gives a uniform $H$-invariant
  embedding on $H\times_{H_p}B(0,\epsilon)$.
\end{proof}

\begin{corollary}[isotropy near an orbit]
  \label{cor:pet-ch5-isotropy-near-orbit}
  \lean{PetersenLib.isotropyGroupNearOrbit}
  \source{petersen:page-0227}
  \uses{thm:pet-ch5-slice-theorem}
  \dcref{ch5:5.6.25}
  For small $v\in T_p^\perp Hp$, $H_{\exp_p(v)}=\{h\in H_p\mid Dh|_pv=v\}$.
\end{corollary}
\begin{proof}
  \uses{cor:pet-ch5-isotropy-near-orbit}
  \sketch
  By \cref{thm:pet-ch5-slice-theorem}, equivariant slice coordinates are unique.
  Thus $h$ fixes $\exp_p(v)$ exactly when $h\in H_p$ and $Dh|_p(v)=v$.
\end{proof}

\begin{corollary}[constant isotropy type gives a Riemannian quotient]
  \label{cor:pet-ch5-conjugate-isotropy-manifold}
  \lean{PetersenLib.constantIsotropyType_riemannianQuotient}
  \source{petersen:page-0227}
  \uses{cor:pet-ch5-isotropy-near-orbit, thm:pet-ch5-free-slice-theorem}
  \dcref{ch5:5.6.26}
  If $H\subset\operatorname{Iso}(M,g)$ is closed and all isotropy groups are
  conjugate, the orbit space is a Riemannian manifold and $M\to H\backslash M$ a
  Riemannian submersion.
\end{corollary}
\begin{proof}
  \uses{cor:pet-ch5-conjugate-isotropy-manifold}
  \sketch
  By \cref{cor:pet-ch5-isotropy-near-orbit}, conjugate isotropy makes all slice
  models have one stabilizer type.  The charts from
  \cref{thm:pet-ch5-slice-theorem} therefore patch as in
  \cref{thm:pet-ch5-free-slice-theorem}; the normal-space metric descends and makes
  the quotient map a Riemannian submersion.
\end{proof}

\section{Completeness}

\subsection{The Hopf–Rinow Theorem}

\begin{theorem}[equivalence of completeness notions]
  \label{thm:pet-ch5-hopf-rinow}
  \lean{PetersenLib.hopfRinowTheorem}
  \source{petersen:page-0227}
  \uses{def:pet-ch5-geodesically-complete, def:pet-ch5-distance}
  \leanok
  \dcref{ch5:5.7.1}
  For a Riemannian manifold $(M,g)$ the following are equivalent:
  \begin{enumerate}
  \item[(1)] $M$ is geodesically complete.
  \item[(2)] $M$ is geodesically complete at some $p$ (all geodesics through $p$
  are defined for all time).
  \item[(3)] $M$ satisfies the Heine–Borel property: every closed bounded set is
  compact.
  \item[(4)] $M$ is metrically complete.
  \end{enumerate}
\end{theorem}
\begin{proof}
  \leanok
  \uses{thm:pet-ch5-hopf-rinow}
  \sketch
  The implications $(1)\Rightarrow(2)$ and $(3)\Rightarrow(4)$ are immediate.
  For $(4)\Rightarrow(1)$, a finite-endpoint maximal geodesic is Cauchy because it
  has constant speed.  Its limit and the local geodesic flow extend it, contrary to
  maximality; equivalently this contradicts
  \cref{prop:pet-ch5-leaves-compact-set}.

  For $(2)\Rightarrow(3)$, fix a complete point $p$.  The short ball identity in
  \cref{cor:pet-ch5-metric-ball-image} starts a minimizing radial geodesic toward
  any $q$.  Choosing on the first distance sphere a point nearest $q$, and using
  \cref{cor:pet-ch5-regular-curve-approx}, gives the additive equality
  $|pq|=\varepsilon+|c(\varepsilon)q|$.  Set
  \[
    A=\{t\in[0,|pq|]\mid |pq|=t+|c(t)q|\}.
  \]
  Triangle inequalities make $A$ downward closed, and continuity makes it closed.
  A short normal ball at $c(a)$ extends $A$ beyond every $a<|pq|$: equality makes
  the concatenated minimizer a segment, hence a geodesic by
  \cref{cor:pet-ch5-segments-are-geodesics}.  Thus $A=[0,|pq|]$ and $c(|pq|)=q$.
  Therefore every closed metric ball is the exponential image of a closed bounded
  tangent ball and is compact, proving the Heine--Borel property.  These arguments
  are the geometric counterparts of the formal proper-space and complete-geodesic
  implications.
\end{proof}

\begin{lemma}[geodesics are admissible piecewise-smooth competitors]
  \label{lem:pet-ch5-geodesics-are-piecewise-smooth}
  \lean{PetersenLib.IsGeodesicOn.isPiecewiseSmoothCurve, PetersenLib.IsGeodesic.isPiecewiseSmoothCurve}
  \source{petersen:page-0204}
  \uses{def:pet-ch5-geodesic, def:pet-ch5-piecewise-smooth-curve, def:pet-ch5-exponential-map}
  \leanok
  A geodesic is $C^\infty$, hence an admissible competitor for the length
  infimum of \cref{def:pet-ch5-distance}: if $\gamma$ is a geodesic on an open
  time-set $s\supseteq[a,b]$ (in particular, if $\gamma$ is a globally defined
  geodesic), then $\gamma|_{[a,b]}$ is piecewise smooth in the sense of
  \cref{def:pet-ch5-piecewise-smooth-curve}.

  This is a regularity bridge rather than a statement of Petersen's: the geodesic
  equation is naturally solved at $C^1$, whereas
  \cref{def:pet-ch5-piecewise-smooth-curve} — and therefore
  \cref{def:pet-ch5-segment} — demands $C^\infty$ on each piece. It is what lets
  a minimizing geodesic be recognized as a \emph{segment}.
\end{lemma}
\begin{proof}
  \leanok
  \uses{def:pet-ch5-exponential-map,thm:pet-ch5-local-uniqueness}
  \sketch
  Near each $t_0$, uniqueness
  (\cref{thm:pet-ch5-local-uniqueness}) identifies $\gamma$ after affine
  reparametrization with a radial ray of $\exp_{\gamma(t_0)}$, hence with a smooth
  curve.  Compactness of $[a,b]$ supplies a finite partition by such neighborhoods,
  proving piecewise smoothness.
\end{proof}

\begin{corollary}[complete manifolds have segments everywhere]
  \label{cor:pet-ch5-complete-segment-exists}
  \lean{PetersenLib.completeManifold_allPointsJoinedBySegment}
  \source{petersen:page-0229}
  \uses{thm:pet-ch5-hopf-rinow}
  \leanok
  \dcref{ch5:5.7.2}
  If $(M,g)$ is complete (in any of the equivalent senses), any two points of $M$
  are joined by a segment.  The formalization takes completeness in the metric
  sense and produces the segment on $[0,1]$, with constant speed $|pq|$; the
  unit-speed $[0,|pq|]$ form used by the notation $\overline{pq}$ of
  \cref{def:pet-ch5-segment} is \cref{cor:pet-ch5-segment-notation-unitspeed}.
\end{corollary}
\begin{proof}
  \leanok
  \uses{cor:pet-ch5-complete-segment-exists,lem:pet-ch5-geodesics-are-piecewise-smooth}
  \sketch
  The $(2)\Rightarrow(3)$ part of \cref{thm:pet-ch5-hopf-rinow} supplies a
  minimizing geodesic.  By \cref{lem:pet-ch5-geodesics-are-piecewise-smooth} it is
  an admissible competitor, and its length equals the distance; hence it is a
  segment in the sense of \cref{def:pet-ch5-segment}.
\end{proof}

\begin{corollary}[unit-speed form: the notation $\overline{pq}$ is meaningful]
  \label{cor:pet-ch5-segment-notation-unitspeed}
  \lean{PetersenLib.completeManifold_exists_isSegment_unitSpeed, PetersenLib.segmentNotation_spec, PetersenLib.IsSegment.curveLength_eq_self_of_domain}
  \source{petersen:page-0229}
  \uses{cor:pet-ch5-complete-segment-exists, def:pet-ch5-segment}
  \leanok
  Let $(M,g)$ be complete and $p,q\in M$.  Then $p$ and $q$ are joined by a segment
  $\sigma:[0,|pq|]\to M$ with $\sigma(0)=p$ and $\sigma(|pq|)=q$, and any such
  segment is parametrized by arclength: $L(\sigma|_{[0,t]})=t$ for all
  $t\in[0,|pq|]$, i.e. the proportionality constant of \cref{def:pet-ch5-segment}
  is forced to $k=1$.  Consequently the notation $\overline{pq}$ of
  \cref{def:pet-ch5-segment} denotes such a segment rather than a junk value.
\end{corollary}
\begin{proof}
  \leanok
  \uses{cor:pet-ch5-complete-segment-exists}
  \sketch
  Take a segment $\gamma:[0,1]\to M$ from
  \cref{cor:pet-ch5-complete-segment-exists}.  Its proportionality constant is
  $k=L(\gamma)=|pq|$.  If $p\ne q$, reparametrize by $t\mapsto t/|pq|$ to obtain
  a segment on $[0,|pq|]$ with constant $1$; if $p=q$, use the constant curve.
  Conversely, evaluating the length relation at $|pq|>0$ forces $k=1$.
\end{proof}

\begin{corollary}[a proper Lipschitz function forces completeness]
  \label{cor:pet-ch5-proper-lipschitz-complete}
  \lean{PetersenLib.properLipschitzFunction_impliesComplete}
  \source{petersen:page-0229}
  \uses{thm:pet-ch5-metric-topology}
  \leanok
  \dcref{ch5:5.7.3}
  If $(M,g)$ admits a proper Lipschitz function $f:M\to\mathbb R$, then $M$ is
  complete.
\end{corollary}
\begin{proof}
  \leanok
  \uses{cor:pet-ch5-proper-lipschitz-complete}
  \sketch
  A $K$-Lipschitz function sends $\bar B(x,r)$ into
  $[f(x)-Kr,f(x)+Kr]$. Properness makes the preimage compact, hence the closed ball
  compact. Thus the metric is proper and complete; this is the metric half of
  \cref{thm:pet-ch5-hopf-rinow}, using the topology from
  \cref{thm:pet-ch5-metric-topology}.
\end{proof}

\subsection{Warped Product Characterization}

\begin{theorem}[global warped product from a conformal Hessian]
  \label{thm:pet-ch5-tashiro}
  \lean{PetersenLib.tashiroWarpedProductTheorem}
  \source{petersen:page-0229,petersen:page-0230}
  \uses{thm:pet-ch5-hopf-rinow, def:pet-ch5-exponential-map}
  \dcref{ch5:5.7.4}
  Let $(M,g)$ be a complete Riemannian $n$-manifold admitting a nontrivial $f$
  with $\operatorname{Hess}f=\lambda g$. Then $(M,g)$ is isometric to a complete
  warped product metric of one of the forms:
  \begin{enumerate}
  \item[(1)] $M=\mathbb R\times N$, $g=dr^2+\rho^2(r)g_N$;
  \item[(2)] $M=\mathbb R^n$, $g=dr^2+\rho^2(r)ds_{n-1}^2$, $r>0$;
  \item[(3)] $M=S^n$, $g=dr^2+\rho^2(r)ds_{n-1}^2$, $r\in[a,b]$.
  \end{enumerate}
\end{theorem}
\begin{proof}
  \uses{thm:pet-ch5-tashiro}
  \sketch
  Choose a connected regular level component $N$ of $f$.  The local
  conformal-Hessian description makes it a closed hypersurface on which
  $|\nabla f|$ is constant.  Follow the unit normal geodesics
  $c_p(t)=\exp_p(t\nabla f|_p/|\nabla f|)$.

  Fix $a<0<b$ and let $C\subset N$ be the (open, by openness of the regular set) set
  of $p$ for which $f$ is regular at $c_p(t)$ for all $t\in[a,b]$. On
  $U=\{c_p(t)\mid p\in C,t\in[a,b]\}$ the local warped product structure gives
  $g|_U=dr^2+\rho^2(r)g_N$ with $r$ the signed distance to $N$,
  $\nabla r=\nabla f/|\nabla f|$, $f=\int\rho(r)\,dr$. For every $p\in N$,
  $\frac{d^2(f\circ c_p)}{dt^2}=\operatorname{Hess}f(\dot c_p,\dot c_p)=\lambda\circ c_p$
  with $(f\circ c_p)(0)=f_0$, $\frac{d(f\circ c_p)}{dt}(0)=|\nabla f|$; on $U$,
  $\lambda\circ c_p=\lambda(f\circ c_p)$ depends only on the value of
  $f\circ c_p$, so this second-order ODE with initial data independent of
  $p\in C$ forces $f\circ c_p(t)$ to be independent of $p\in C$, and likewise
  $\frac{d(f\circ c_p)}{dt}(t)=|\nabla f|_{c_p(t)}|=\rho(t)$ is independent of
  $p\in C$. Continuity of $|\nabla f|_{c_p(t)}|$ in $(p,t)$ and regularity of $f$
  on $U$ prevent $|\nabla f|_{c_p(t)}|$ from vanishing as $p\to\partial C$ (for
  fixed $t$), so $C$ is also closed in $N$; as $N$ is connected, $C=N$.

  Extend the warped product to the maximal common regular interval $(a,b)$.  If
  both endpoints are infinite, this is case (1).  At a finite endpoint,
  $\rho=|\nabla f|\to0$, so the level collapses to one critical point.  Its incoming
  normal directions form an open and closed subset of the unit sphere, hence all
  directions occur and the level metric is round.  One finite endpoint gives case
  (2), and two finite endpoints give case (3).
\end{proof}

\begin{theorem}[global classification for $\operatorname{Hess}f=(\alpha f+\beta)g$]
  \label{thm:pet-ch5-hessian-classification}
  \lean{PetersenLib.affineHessianGlobalClassification}
  \source{petersen:page-0231}
  \uses{thm:pet-ch5-tashiro}
  \dcref{ch5:5.7.5}
  Let $(M,g)$ be complete of dimension $n$ with a nontrivial $f$,
  $\operatorname{Hess}f=(\alpha f+\beta)g$, $\alpha,\beta\in\mathbb R$. Then
  $(M,g)$ is one of:
  \begin{enumerate}
  \item[(a)] $(I\times S^{n-1},\,dr^2+\mathrm{sn}_k^2(r)ds_{n-1}^2)$: a constant
  curvature space form;
  \item[(b)] $(\mathbb R\times N,\,dr^2+g_N)$: a product metric;
  \item[(c)] $\big(\mathbb R\times N,\,dr^2+(A\exp(\sqrt\alpha r)+B\exp(-\sqrt\alpha r))^2h\big)$,
  $A,B\ge0$.
  \end{enumerate}
\end{theorem}
\begin{proof}
  \uses{thm:pet-ch5-hessian-classification}
  \sketch
  By \cref{thm:pet-ch5-tashiro}, $g=dr^2+\rho^2g_N$ with
  $\rho=f'$ and $f''=\alpha f+\beta$.  If $\alpha=\beta=0$, $\rho$ is constant
  (case (b)); if $\alpha=0\ne\beta$, $\rho$ is affine and completeness gives the
  Euclidean case (a).  For $\alpha\ne0$, shift $f$ by $\beta/\alpha$.  Negative
  $\alpha$ yields sine warping on a compact interval, hence a round sphere.  For
  positive $\alpha$,
  $\rho=Ae^{\sqrt\alpha r}+Be^{-\sqrt\alpha r}$; nonnegative coefficients give
  case (c), while opposite signs reduce to hyperbolic-sine warping, the negative
  curvature case (a).
\end{proof}

\begin{remark}[critical points and transnormal functions]
  \label{rem:pet-ch5-obata-transnormal}
  \lean{PetersenLib.remark_obataTransnormalFunctions}
  \source{petersen:page-0232}
  \uses{thm:pet-ch5-hessian-classification}
  \dcref{ch5:5.7.6}
  In case (a) $f$ has at least one critical point, while in (b), (c) it has none.
  Obata (1961) established the round-sphere case via $\operatorname{Hess}f=(1-f)g$.
  Transnormal functions on a complete manifold give a similar topological
  decomposition, with zero, one, or two critical values, all level sets (including
  critical ones) smooth submanifolds, and $\tilde f=\phi(r)$ for $r$ the signed
  distance to a fixed level set.
\end{remark}

\subsection{The Segment Domain}

\begin{definition}[segment domain and its star interior]
  \label{def:pet-ch5-segment-domain}
  \lean{PetersenLib.segmentDomain,PetersenLib.segmentDomainStarInterior}
  \source{petersen:page-0232,petersen:page-0233}
  \uses{def:pet-ch5-segment, thm:pet-ch5-hopf-rinow}
  \leanok
  For a complete $(M,g)$ and $p\in M$, the \emph{segment domain} is
  $\operatorname{seg}(p)=\{v\in T_pM\mid\exp_p(tv):[0,1]\to M\text{ is a segment}\}$,
  a closed star-shaped set with $M=\exp_p(\operatorname{seg}(p))$ (by
  \cref{thm:pet-ch5-hopf-rinow}). Its \emph{star interior} is
  $\operatorname{seg}^0(p)=\{sv\mid s\in[0,1),\,v\in\operatorname{seg}(p)\}$.
\end{definition}

\begin{proposition}[injectivity of $\exp_p$ on the star interior]
  \label{prop:pet-ch5-segment-domain-injective}
  \lean{PetersenLib.expMap_injectiveOnSegmentDomainStarInterior, PetersenLib.segment_eq_expMap_ray_of_mem_segmentDomainStarInterior', PetersenLib.segmentDomainStarInterior_subset_segmentDomain}
  \source{petersen:page-0233}
  \uses{def:pet-ch5-segment-domain, cor:pet-ch5-segments-are-geodesics}
  \leanok
  \dcref{ch5:5.7.7}
  If $x\in\exp_p(\operatorname{seg}^0(p))$, it is joined to $p$ by a unique
  segment; in particular $\exp_p$ is injective on $\operatorname{seg}^0(p)$.
\end{proposition}
\begin{proof}
  \leanok
  \uses{prop:pet-ch5-segment-domain-injective}
  \sketch
  Concatenate any competing segment to $x=\sigma(t_0)$, $t_0<1$, with the tail of
  $\sigma$.  The result is again minimizing, hence a geodesic by
  \cref{cor:pet-ch5-segments-are-geodesics}.  It agrees with $\sigma$ after $t_0$,
  so geodesic uniqueness propagates equality across the joining point.  Thus the
  initial segments, and consequently their exponential initial vectors, coincide.
\end{proof}

\begin{lemma}[nonsingularity of $\exp_p$ on the star interior]
  \label{lem:pet-ch5-exp-nonsingular-segment-domain}
  \source{petersen:page-0233}
  \uses{prop:pet-ch5-segment-domain-injective, cor:pet-ch5-uniform-injectivity-radius-compact}
  \notready
  \dcref{ch5:5.7.8}
  \notready
  $D\exp_p$ is nonsingular at every point of $\operatorname{seg}^0(p)$.

  \emph{Formalization status.} Not formalized. This lemma — equivalently, that
  there is no conjugate point strictly before the cut point — is the genuine
  dependency of \cref{prop:pet-ch5-segment-domain-open}; see
  \cref{rem:pet-ch5-segment-domain-open-blocker}. Chapter 6 contains local
  Jacobi--$D\exp$ and Hessian identities, but those files import Chapter 5 and
  concern the chart exponential on a normal ball. Using them here requires
  moving shared Jacobi theory below Chapter 5 and extending it to the intrinsic
  exponential along the whole segment up to the cut point.
\end{lemma}
\begin{proof}
  \uses{lem:pet-ch5-exp-nonsingular-segment-domain}
  \sketch
  If the first singular point on a segment ray were $v$, a kernel vector would
  generate a Jacobi field $J$ with $J(t)\to0$ as $t\uparrow1$.  Hence along a
  sequence $t_n\uparrow1$,
  $\operatorname{Hess}r(J,J)/|J|^2\to-\infty$, using
  $\partial_r|J|^2=2\operatorname{Hess}r(J,J)$.

  The segment extends slightly beyond $v$.  By
  \cref{cor:pet-ch5-uniform-injectivity-radius-compact}, the distance
  $\tilde r(x)=|xc(1+\varepsilon)|$ is smooth near $c(1)$.  The triangle inequality
  makes $e=r+\tilde r$ minimal along $c$, so $\operatorname{Hess}e\ge0$.  But
  \[
    \frac{\operatorname{Hess}e(J,J)}{|J|^2}\Big|_{\exp_p(t_nv)}
    =\frac{\operatorname{Hess}r(J,J)}{|J|^2}\Big|_{\exp_p(t_nv)}
    +\frac{\operatorname{Hess}\tilde r(J,J)}{|J|^2}\Big|_{\exp_p(t_nv)}\to-\infty,
  \]
  because the second term is bounded, a contradiction.
\end{proof}

\begin{lemma}[the star boundary of $\operatorname{seg}(p)$]
  \label{lem:pet-ch5-cut-locus-characterization}
  \source{petersen:page-0235}
  \uses{lem:pet-ch5-exp-nonsingular-segment-domain, rem:pet-ch5-segment-domain-chart-artifact}
  \notready
  \dcref{ch5:5.7.9}
  \notready
  If $v\in\operatorname{seg}(p)-\operatorname{seg}^0(p)$, then either (1) there is
  $w\ne v$ in $\operatorname{seg}(p)$ with $\exp_p(v)=\exp_p(w)$, or (2) $D\exp_p$
  is singular at $v$.

  \emph{Formalization status.} Not formalized, and it cannot be formalized against
  the on-disk definitions as they stand: read there, the statement is \emph{false}
  rather than merely unproven, because \texttt{segmentDomain} inherits the chart
  artifact of the on-disk exponential. See
  \cref{rem:pet-ch5-segment-domain-chart-artifact}. The lemma as stated above — at
  the intrinsic exponential, which is Petersen's — is the true statement and is what
  a formalization should target.
\end{lemma}
\begin{proof}
  \uses{lem:pet-ch5-cut-locus-characterization}
  \sketch
  For $t_n\downarrow1$, choose minimizing segments to $\exp_p(t_nv)$ and take a
  convergent subsequence of their initial vectors, with limit $w$ and $|w|=|v|$.
  If $w\ne v$, then $\exp_pw=\exp_pv$, giving (1).  If $w=v$ and $D\exp_p|_v$ were
  nonsingular, local injectivity would identify the nearby initial vectors with
  $t_nv$, making the radial geodesic minimizing beyond $v$.  This contradicts
  $v\notin\operatorname{seg}^0(p)$, so case (2) holds.
\end{proof}

\begin{remark}[the Lean \texttt{segmentDomain} and \texttt{cutLocus} are chart artifacts]
  \label{rem:pet-ch5-segment-domain-chart-artifact}
  \uses{def:pet-ch5-segment-domain, rem:pet-ch5-injectivity-radius-chart-artifact}
  \texttt{PetersenLib.segmentDomain} is defined as the set of $v\in T_pM$ for which
  $t\mapsto\texttt{expMap}(g,p,tv)$ is a segment on $[0,1]$, and
  \texttt{segmentDomainStarInterior} and \texttt{cutLocus} are defined from it.
  All three therefore inherit the defect recorded in
  \cref{rem:pet-ch5-injectivity-radius-chart-artifact}: \texttt{expMap} is built
  from the geodesic vector field of the \emph{single chart at $p$}, which is
  identically $0$ off that chart's source, so a geodesic witnessing membership of
  $\operatorname{seg}(p)$ can never leave the chart at $p$. The predicate
  \texttt{IsSegment} that such a geodesic is required to satisfy is, by contrast,
  chart-free.

  The mismatch is not a gap but a falsity: on a manifold whose chart at $p$ is
  small, there are genuine segments emanating from $p$ that leave that chart, and
  the corresponding $v$ lie outside the on-disk $\operatorname{seg}(p)$ while
  satisfying the hypotheses of \cref{lem:pet-ch5-cut-locus-characterization}. That
  lemma read against these definitions is therefore \emph{false}, not unproven, and
  it must not be given a \texttt{\textbackslash leanok} against them. The repair is
  the one already applied to
  \cref{cor:pet-ch5-uniform-injectivity-radius-compact} and to part~(2) of
  \cref{prop:pet-ch5-local-isometry-properties}: redefine the segment domain
  through the intrinsic maximal geodesic. That work has not been done.
\end{remark}

\begin{proposition}[the star interior is open]
  \label{prop:pet-ch5-segment-domain-open}
  \source{petersen:page-0235}
  \uses{lem:pet-ch5-cut-locus-characterization, lem:pet-ch5-exp-nonsingular-segment-domain}
  \notready
  \dcref{ch5:5.7.10}
  \notready
  $\operatorname{seg}^0(p)$ is open in $T_pM$.

  \emph{Formalization status.} Not formalized, and genuinely blocked; see
  \cref{rem:pet-ch5-segment-domain-open-blocker}.
\end{proposition}
\begin{proof}
  \uses{prop:pet-ch5-segment-domain-open}
  \sketch
  By \cref{lem:pet-ch5-exp-nonsingular-segment-domain}, choose a neighborhood $V$
  of $v\in\operatorname{seg}^0(p)$ on which $\exp_p$ is injective.  For
  $v_i\to v$, select $w_i\in\operatorname{seg}(p)$ with
  $\exp_pv_i=\exp_pw_i$.  Any limit $w\ne v$ contradicts
  \cref{lem:pet-ch5-cut-locus-characterization}; hence $w_i\to v$ and local
  injectivity gives $w_i=v_i$.  The same dichotomy, together with nonsingularity
  on $V$, excludes $v_i$ from the star boundary.  Thus $v_i\in\operatorname{seg}^0(p)$.
\end{proof}

\begin{remark}[what \cref{prop:pet-ch5-segment-domain-open} is really blocked on]
  \label{rem:pet-ch5-segment-domain-open-blocker}
  \uses{prop:pet-ch5-segment-domain-open, lem:pet-ch5-exp-nonsingular-segment-domain}
  Since $\operatorname{seg}(p)$ is star-shaped and $\operatorname{seg}^0(p)$ is its
  set of proper radial multiples, the assertion of
  \cref{prop:pet-ch5-segment-domain-open} amounts to
  $\operatorname{seg}^0(p)=\operatorname{int}\operatorname{seg}(p)$; the inclusion
  $\operatorname{int}\operatorname{seg}(p)\subset\operatorname{seg}^0(p)$ is the
  easy half. The substance of the proposition is the reverse inclusion, which is
  lower semicontinuity of the cut time $v\mapsto\sup\{t\mid tv\in
  \operatorname{seg}(p)\}$.

  That is not obtainable from the star-shapedness of $\operatorname{seg}(p)$ alone.
  The proof above obtains it from
  \cref{lem:pet-ch5-exp-nonsingular-segment-domain} — the statement that $D\exp_p$
  is nonsingular on $\operatorname{seg}^0(p)$, i.e.\ that there is no conjugate
  point strictly before the cut point — which supplies the local diffeomorphism $V$
  the argument runs in. \Cref{lem:pet-ch5-exp-nonsingular-segment-domain} is itself
  unformalized: the relevant local Jacobi and Hessian components live in Chapter 6,
  above this file in the import graph, and have not been transported to the intrinsic
  exponential along an arbitrary segment. This is the genuine dependency, not a
  bookkeeping gap, and the openness cannot be obtained directly from star-shapedness.
\end{remark}

\begin{definition}[cut locus]
  \label{def:pet-ch5-cut-locus}
  \lean{PetersenLib.cutLocus}
  \source{petersen:page-0236}
  \uses{def:pet-ch5-segment-domain}
  \leanok
  The set $\operatorname{seg}(p)-\operatorname{seg}^0(p)\subset T_pM$ is the
  \emph{cut locus} of $p$ (in $T_pM$). By
  \cref{prop:pet-ch5-segment-domain-open} and
  \cref{lem:pet-ch5-exp-nonsingular-segment-domain},
  $r(x)=|px|$ is smooth on the open dense set $U_p=\{p\}\cup\exp_p(\operatorname{seg}^0(p))$
  and fails to be smooth off $U_p$.
\end{definition}

\begin{corollary}[smoothness of $r$ up to the cut point]
  \label{cor:pet-ch5-cut-locus-smoothness}
  \lean{PetersenLib.distanceFunction_smoothUpToCutPoint}
  \source{petersen:page-0236}
  \uses{def:pet-ch5-cut-locus, lem:pet-ch5-cut-locus-characterization}
  \dcref{ch5:5.7.11}
  Let $c:[0,\infty)\to M$ be a geodesic, $c(0)=p$, and
  $\operatorname{cut}(\dot c(0))=\sup\{t\mid c|_{[0,t]}\text{ is a segment}\}$. Then
  $r$ is smooth at $c(t)$ for $t<\operatorname{cut}(\dot c(0))$ but not at
  $x=c(\operatorname{cut}(\dot c(0)))$, the failure being either that $\exp_p$ is
  not one-to-one at $x$ or that $x$ is a critical value of $\exp_p$.
\end{corollary}
\begin{proof}
  \uses{cor:pet-ch5-cut-locus-smoothness}
  \sketch
  Before the cut time, $tv\in\operatorname{seg}^0(p)$; at the cut time it lies on
  the star boundary.  The smoothness assertion follows from
  \cref{def:pet-ch5-cut-locus}, and the two possible failures are exactly those in
  \cref{lem:pet-ch5-cut-locus-characterization}.
\end{proof}

\subsection{The Injectivity Radius}

\begin{remark}[injectivity radius via the cut locus]
  \label{rem:pet-ch5-injectivity-radius-cut-locus}
  \lean{PetersenLib.injectivityRadius_viaCutLocus}
  \source{petersen:page-0236}
  \uses{def:pet-ch5-cut-locus, def:pet-ch5-injectivity-radius}
  In a complete manifold, $\operatorname{inj}(p)$ (the largest $\varepsilon$ with
  $\exp_p:B(0,\varepsilon)\to B(p,\varepsilon)$ a diffeomorphism) equals $|v|$ for
  $v$ the closest point of the cut locus
  $\operatorname{seg}(p)-\operatorname{seg}^0(p)$ to the origin.
\end{remark}

\begin{lemma}[characterization of the injectivity radius]
  \label{lem:pet-ch5-klingenberg}
  \lean{PetersenLib.klingenbergLemma}
  \source{petersen:page-0236}
  \uses{rem:pet-ch5-injectivity-radius-cut-locus, lem:pet-ch5-cut-locus-characterization}
  \dcref{ch5:5.7.12}
  Suppose $v\in\operatorname{seg}(p)-\operatorname{seg}^0(p)$ with
  $|v|=\operatorname{inj}(p)$. Either
  \begin{enumerate}
  \item[(1)] there is exactly one other $v'$ with $\exp_p(v')=\exp_p(v)$, and $v'$
  satisfies $\frac{d}{dt}\big|_{t=1}\exp_p(tv')=-\frac{d}{dt}\big|_{t=1}\exp_p(tv)$
  (the two segments meet to form a smooth geodesic loop at $x=\exp_p(v)$), or
  \item[(2)] $x=\exp_p(v)$ is a critical value of $\exp_p|_{\operatorname{seg}(p)}$.
  \end{enumerate}
\end{lemma}
\begin{proof}
  \uses{lem:pet-ch5-klingenberg}
  \sketch
  Suppose $x$ is regular and two minimizing segments reach it with nonantipodal
  final velocities.  Choose $w\in T_xM$ pairing negatively with both velocities.
  Local inverse branches of $\exp_p$ lift a curve tangent to $w$ to two curves
  $v_i(s)$.  First variation makes both $|v_i(s)|$ decrease for $s>0$, producing
  distinct preimages inside $B(0,\operatorname{inj}(p))$, contrary to injectivity.
  Hence the final velocities are opposite.  If no second segment exists, the
  alternative in \cref{lem:pet-ch5-cut-locus-characterization} makes $x$ a critical
  value.
\end{proof}

\section{Exercises}

\begin{remark}[fixed-time existence implies completeness]
  \label{rem:pet-ch5-ex-1}
  \lean{PetersenLib.exercise5_9_1}
  \source{petersen:page-0237}
  \uses{def:pet-ch5-geodesically-complete}
  \dcref{ch5:5.9.1}
  If $(M,g)$ has the property that all unit-speed geodesics exist for a fixed time
  $\varepsilon>0$, show $(M,g)$ is geodesically complete.
\end{remark}

\begin{remark}[chain rule for velocity and acceleration]
  \label{rem:pet-ch5-ex-2}
  \lean{PetersenLib.exercise5_9_2}
  \source{petersen:page-0237,petersen:page-0238}
  \uses{def:pet-ch5-acceleration}
  \dcref{ch5:5.9.2}
  For $c:I\to M$ and $\phi:J\to I$, show
  $\frac{d(c\circ\phi)}{dt}=\dot c\circ\phi\cdot\frac{d\phi}{dt}$ and
  $\frac{d^2(c\circ\phi)}{dt^2}=\dot c\circ\phi\cdot\frac{d^2\phi}{dt^2}
  +\ddot c\circ\phi\cdot\big(\frac{d\phi}{dt}\big)^2$.
\end{remark}

\begin{remark}[geodesic equation with orthogonal coordinates]
  \label{rem:pet-ch5-ex-3}
  \lean{PetersenLib.exercise5_9_3}
  \source{petersen:page-0238}
  \uses{def:pet-ch5-geodesic}
  \dcref{ch5:5.9.3}
  If the coordinate vector fields of a chart are orthogonal ($g_{ij}=0$, $i\ne j$),
  show the geodesic equations become
  $\frac{d}{dt}\big(g_{ii}\frac{dc^i}{dt}\big)=\frac12\sum_j\frac{\partial g_{jj}}
  {\partial x^i}\big(\frac{dc^j}{dt}\big)^2$.
\end{remark}

\begin{remark}[reparametrization to a geodesic]
  \label{rem:pet-ch5-ex-4}
  \lean{PetersenLib.exercise5_9_4}
  \source{petersen:page-0238}
  \uses{def:pet-ch5-geodesic}
  \dcref{ch5:5.9.4}
  Show a regular curve can be reparametrized to be a geodesic iff its acceleration
  is everywhere tangent to the curve.
\end{remark}

\begin{remark}[a complete open submanifold is everything]
  \label{rem:pet-ch5-ex-5}
  \lean{PetersenLib.exercise5_9_5}
  \source{petersen:page-0238}
  \uses{def:pet-ch5-geodesically-complete}
  \dcref{ch5:5.9.5}
  If $O\subset(M,g)$ is open and $(O,g)$ is complete, show $O=M$.
\end{remark}

\begin{remark}[Misner completeness implies completeness]
  \label{rem:pet-ch5-ex-6}
  \lean{PetersenLib.exercise5_9_6}
  \source{petersen:page-0238}
  \uses{def:pet-ch5-geodesically-complete}
  \dcref{ch5:5.9.6}
  Call $(M,g)$ \emph{Misner complete} if every geodesic $c:(a,b)\to M$ with
  $b-a<\infty$ lies in a compact set. Show Misner completeness implies
  completeness.
\end{remark}

\begin{remark}[length-minimizers reparametrize a segment]
  \label{rem:pet-ch5-ex-7}
  \lean{PetersenLib.exercise5_9_7}
  \source{petersen:page-0238}
  \uses{def:pet-ch5-segment}
  \dcref{ch5:5.9.7}
  For $c\in\Omega_{p,q}$ with $L(c)=|pq|$: (1) show $L(c|_{[a,b]})=|c(a)c(b)|$ for
  all $a,b\in[0,1]$; (2) show there is a segment $\sigma\in\Omega_{p,q}$ and
  monotone $\varphi:[0,1]\to[0,1]$ (not necessarily everywhere smooth) with
  $c=\sigma\circ\varphi$.
\end{remark}

\begin{remark}[completeness under a larger metric]
  \label{rem:pet-ch5-ex-8}
  \lean{PetersenLib.exercise5_9_8}
  \source{petersen:page-0238}
  \uses{def:pet-ch5-geodesically-complete}
  \dcref{ch5:5.9.8}
  If $(M,g)$ is metrically complete and $\bar g\ge g$ another metric on $M$, show
  $(M,\bar g)$ is metrically complete.
\end{remark}

\begin{remark}[homogeneous manifolds are complete]
  \label{rem:pet-ch5-ex-9}
  \lean{PetersenLib.exercise5_9_9}
  \source{petersen:page-0238}
  \uses{def:pet-ch5-geodesically-complete}
  \dcref{ch5:5.9.9}
  A Riemannian manifold is \emph{homogeneous} if its isometry group acts
  transitively. Show homogeneous manifolds are geodesically complete.
\end{remark}

\begin{remark}[completeness of $\mathbb R\times N$ warped products]
  \label{rem:pet-ch5-ex-10}
  \lean{PetersenLib.exercise5_9_10}
  \source{petersen:page-0238}
  \uses{def:pet-ch5-geodesically-complete}
  \dcref{ch5:5.9.10}
  For $(M,g)=(\mathbb R\times N,dr^2+g_r)$ with $(N,g_r)$ complete for all
  $r\in\mathbb R$ (e.g.\ $g_r=\rho^2(r)g_N$, $\rho:\mathbb R\to(0,\infty)$, $(N,g_N)$
  metrically complete), show $(M,g)$ is metrically complete.
\end{remark}

\begin{remark}[completeness on a half-line warped product]
  \label{rem:pet-ch5-ex-11}
  \lean{PetersenLib.exercise5_9_11}
  \source{petersen:page-0238}
  \uses{def:pet-ch5-geodesically-complete}
  \dcref{ch5:5.9.11}
  For $(M,g)=((0,\infty)\times N,dr^2+\rho^2(r)g_N)$, $\rho:(0,\infty)\to
  (0,\infty)$, $(N,g_N)$ complete: give examples that are complete and examples
  that are not.
\end{remark}

\begin{remark}[splitting via a Hessian-flat submersion to $\mathbb R^k$]
  \label{rem:pet-ch5-ex-12}
  \lean{PetersenLib.exercise5_9_12}
  \source{petersen:page-0238}
  \uses{def:pet-ch5-geodesically-complete}
  \dcref{ch5:5.9.12}
  If $F:(M,g)\to(\mathbb R^k,g_{\mathbb R^k})$ is a Riemannian submersion, $(M,g)$
  complete, and each component of $F$ has zero Hessian, show
  $(M,g)=(N,h)\times(\mathbb R^k,g_{\mathbb R^k})$.
\end{remark}

\begin{remark}[filling a gap in the submetry theorem]
  \label{rem:pet-ch5-ex-13}
  \lean{PetersenLib.exercise5_9_13}
  \source{petersen:page-0239}
  \uses{thm:pet-ch5-berestovskii-submetry}
  \dcref{ch5:5.9.13}
  Find and fill in the gap in the proof of \cref{thm:pet-ch5-berestovskii-submetry}.
\end{remark}

\begin{remark}[isotropic implies homogeneous]
  \label{rem:pet-ch5-ex-14}
  \lean{PetersenLib.exercise5_9_14}
  \source{petersen:page-0239}
  \uses{def:pet-ch5-fixed-point-set}
  \dcref{ch5:5.9.14}
  A manifold is \emph{isotropic} at $p$ if $\operatorname{Iso}_p$ acts transitively
  on the unit sphere in $T_pM$. Show isotropy at every point implies homogeneity.
\end{remark}

\begin{remark}[orthonormal coordinates are distance coordinates]
  \label{rem:pet-ch5-ex-15}
  \lean{PetersenLib.exercise5_9_15}
  \source{petersen:page-0239}
  \uses{rem:pet-ch5-distance-function-geodesics}
  \dcref{ch5:5.9.15}
  If coordinates satisfy $g_{ii}=\delta_{ij}$ (i.e.\ $g_{i1}=\delta_{i1}$ for the
  first coordinate direction), show $x^1$ is a distance function.
\end{remark}

\begin{remark}[distance functions persist under a compatible metric change]
  \label{rem:pet-ch5-ex-16}
  \lean{PetersenLib.exercise5_9_16}
  \source{petersen:page-0239}
  \uses{rem:pet-ch5-distance-function-geodesics}
  \dcref{ch5:5.9.16}
  Let $r:U\to\mathbb R$ be a distance function w.r.t.\ $g$. If $\bar g$ satisfies
  $\bar g(\nabla r,v)=g(\nabla r,v)$ for all $v$ (gradient taken w.r.t.\ $g$), show
  $r$ is also a distance function w.r.t.\ $\bar g$.
\end{remark}

\begin{remark}[projective models of $S^n(R)$ and $H^n(R)$]
  \label{rem:pet-ch5-ex-17}
  \lean{PetersenLib.exercise5_9_17}
  \source{petersen:page-0239}
  \uses{ex:pet-ch5-sphere-great-circles, ex:pet-ch5-hyperbolic-hyperbolas}
  \dcref{ch5:5.9.17}
  For $x\in\mathbb R^{n+1}$, $x^{n+1}>0$, the projection to $x^{n+1}=R$ is
  $P(x)=R(x^1/x^{n+1},\dots,x^n/x^{n+1},1)$. (1) Verify this formula; (2) show
  geodesics of $S^n(R)$, $H^n(R)$ are intersections with $2$-planes through the
  origin; (3) show the upper hemisphere of $S^n(R)$ projects onto all of
  $x^{n+1}=R$; (4) show $H^n(R)$ projects onto an open disc of radius $R$; (5) show
  geodesics of $S^n(R)$, $H^n(R)$ project to straight lines.
\end{remark}

\begin{remark}[a conformal change making a manifold complete]
  \label{rem:pet-ch5-ex-18}
  \lean{PetersenLib.exercise5_9_18}
  \source{petersen:page-0239}
  \uses{cor:pet-ch5-proper-lipschitz-complete}
  \dcref{ch5:5.9.18}
  Show any $(M,g)$ admits a conformal change $(M,\lambda^2g)$ that is complete;
  take $\lambda:M\to[1,\infty)$ proper and rapidly growing.
\end{remark}

\begin{remark}[induced versus ambient distance on subsets of $\mathbb R^n$]
  \label{rem:pet-ch5-ex-19}
  \lean{PetersenLib.exercise5_9_19}
  \source{petersen:page-0239}
  \uses{def:pet-ch5-distance}
  \dcref{ch5:5.9.19}
  On $U\subset\mathbb R^n$ compare the Riemannian-induced distance to the ambient
  $\mathbb R^n$ distance: (1) give a non-convex $U$ where they agree; (2) give a
  convex $U$ where they do not.
\end{remark}

\begin{remark}[the $T$-tensor detects totally geodesic submanifolds]
  \label{rem:pet-ch5-ex-20}
  \lean{PetersenLib.exercise5_9_20}
  \source{petersen:page-0239}
  \uses{def:pet-ch5-fixed-point-set}
  \dcref{ch5:5.9.20}
  For $M\subset(\bar M,\bar g)$, using the $T$-tensor, show $T=0$ on $M$ iff
  $M\subset(\bar M,\bar g)$ is totally geodesic.
\end{remark}

\begin{remark}[convexity along geodesics via the Hessian]
  \label{rem:pet-ch5-ex-21}
  \lean{PetersenLib.exercise5_9_21}
  \source{petersen:page-0239,petersen:page-0240}
  \uses{def:pet-ch5-geodesic}
  \dcref{ch5:5.9.21}
  For $f:(M,g)\to\mathbb R$ smooth: (1) compute the first and second derivatives of
  $f\circ c$ along a geodesic $c$; (2) deduce that at a local max/min of $f$ the
  gradient vanishes and the Hessian is nonpositive/nonnegative; (3) show $f$ has
  everywhere nonnegative Hessian iff $f\circ c$ is convex for every geodesic $c$.
\end{remark}

\begin{remark}[leading term of the volume density]
  \label{rem:pet-ch5-ex-22}
  \lean{PetersenLib.exercise5_9_22}
  \source{petersen:page-0240}
  \uses{def:pet-ch5-radial-distance-function}
  \dcref{ch5:5.9.22}
  If the volume form near a point is $\lambda(r,\theta)\,dr\wedge\mathrm{vol}_{n-1}$,
  show $\lambda(r,\theta)=r^{n-1}+O(r^{n+1})$.
\end{remark}

\begin{remark}[segments to a properly embedded submanifold]
  \label{rem:pet-ch5-ex-23}
  \lean{PetersenLib.exercise5_9_23}
  \source{petersen:page-0240}
  \uses{def:pet-ch5-segment, thm:pet-ch5-hopf-rinow}
  \dcref{ch5:5.9.23}
  Let $N\subset M$ be properly embedded, $(M,g)$ complete, $|xN|=\inf\{|xp|\mid
  p\in N\}$, and $\sigma:[a,b]\to M$ a unit-speed segment from $N$ to $x$. (1) Show
  $\sigma$ is also a segment from $N$ to $\sigma(t)$, $t<b$, and $\dot\sigma(a)\perp
  N$; (2) show every point of $M$ has a segment to $N$ when $N$ is closed; (3) show
  there is always an open neighborhood of $N$ where points are joined to $N$ by
  segments; (4) show $r(x)=|xN|$ is smooth on a punctured neighborhood of $N$; (5)
  show integral curves of $\nabla r$ are the geodesics perpendicular to $N$.
\end{remark}

\begin{remark}[cut locus on the square torus]
  \label{rem:pet-ch5-ex-24}
  \lean{PetersenLib.exercise5_9_24}
  \source{petersen:page-0240}
  \uses{def:pet-ch5-cut-locus}
  \dcref{ch5:5.9.24}
  Find the cut locus on the square torus $\mathbb R^2/\mathbb Z^2$.
\end{remark}

\begin{remark}[cut locus on the sphere and $\mathbb{RP}^n$]
  \label{rem:pet-ch5-ex-25}
  \lean{PetersenLib.exercise5_9_25}
  \source{petersen:page-0240}
  \uses{def:pet-ch5-cut-locus}
  \dcref{ch5:5.9.25}
  Find the cut locus on the sphere and on real projective space with the constant
  curvature metrics.
\end{remark}

\begin{remark}[first-order comparison of exponential images]
  \label{rem:pet-ch5-ex-26}
  \lean{PetersenLib.exercise5_9_26}
  \source{petersen:page-0240}
  \uses{def:pet-ch5-exponential-map}
  \dcref{ch5:5.9.26}
  Show $|\exp_p(v)\exp_p(w)|=|v-w|+O(r^2)$ for $|v|,|w|\le r$.
\end{remark}

\begin{remark}[second-order asymptotics of $\operatorname{Hess}r$]
  \label{rem:pet-ch5-ex-27}
  \lean{PetersenLib.exercise5_9_27}
  \source{petersen:page-0240,petersen:page-0241}
  \uses{lem:pet-ch5-normal-coords-first-order}
  \dcref{ch5:5.9.27}
  In exponential normal coordinates $x^i$ at $p$: (1) show
  $(\operatorname{Hess}x^i)_{kl}=\Gamma^i_{kl}=O(r)$; (2) using
  $\tfrac12r^2=\tfrac12\sum_i(x^i)^2$ and $g=\delta_{ij}+O(r^2)$, show
  $\operatorname{Hess}\tfrac12r^2=g+O(r^2)$; (3) show $\operatorname{Hess}r=
  \tfrac1rg_r+O(r)$.
\end{remark}

\begin{remark}[one-sided directional derivatives of distance to a submanifold]
  \label{rem:pet-ch5-ex-28}
  \lean{PetersenLib.exercise5_9_28}
  \source{petersen:page-0241}
  \uses{def:pet-ch5-segment, lem:pet-ch5-first-variation-formula, rem:pet-ch5-ex-26}
  \dcref{ch5:5.9.28}
  Let $(M,g)$ be complete, $K\subset M$ compact (or properly embedded), $r(x)=|xK|$.
  (1) Show that if $r$ is differentiable at $x\notin K$ then $\overline{xK}$
  contains only one vector; (2) for unit-speed $c$, show that if $f=r\circ c$ is
  differentiable at $t$, all vectors $\overline{c(t)K}$ form the same angle with
  $\dot c(t)$; (3) show the upper right Dini derivative satisfies
  $\overline D^+f(t_0)\le g(\dot c(t_0),-\overline{c(t_0)K})$, via the first
  variation formula for a variation toward $K$; (4) show
  $|c(t_0)c(t)|=h+O(h^2)$ for $h=t-t_0$ small; (5)–(6) via a comparison point on a
  segment to $K$ and \cref{rem:pet-ch5-ex-26}, show the lower right Dini derivative
  satisfies $D^+f(t_0)\ge\min_{\overline{c(t_0)K}}g(\dot c(t_0),-\overline{c(t_0)K})$;
  (7) conclude the one-sided derivatives of $f$ exist everywhere.
\end{remark}

\begin{remark}[metric-space length agrees with the Riemannian length]
  \label{rem:pet-ch5-ex-29}
  \lean{PetersenLib.exercise5_9_29}
  \source{petersen:page-0242}
  \uses{def:pet-ch5-curve-length, rem:pet-ch5-ex-26}
  \dcref{ch5:5.9.29}
  In a metric space $(X,|\cdot|)$, define
  $L(c)=\sup\{\sum|c(t_i)c(t_{i+1})|\mid a=t_1\le\dots\le t_k=b\}$. (1) Show a
  curve of finite length is absolutely continuous; (2) show the Cantor step
  function is a counterexample to the converse; (3) show this definition recovers
  the Riemannian curve length for smooth (and absolutely continuous) curves, using
  \cref{rem:pet-ch5-ex-26} to approximate the Euclidean case; (4) show an
  absolutely continuous $c:[a,b]\to M$ of length $|c(a)c(b)|$ factors as
  $c=\sigma\circ\varphi$ for a segment $\sigma$ and monotone $\varphi$.
\end{remark}

\begin{remark}[a synthetic characterization of exponential coordinates]
  \label{rem:pet-ch5-ex-30}
  \lean{PetersenLib.exercise5_9_30}
  \source{petersen:page-0242}
  \uses{rem:pet-ch5-radial-isometry}
  \dcref{ch5:5.9.30}
  If coordinates $x^i$ near $p$ satisfy $x^i(p)=0$ and $g_{ij}x^j=\delta_{ij}x^j$,
  show they must be exponential normal coordinates. Hint: set
  $r=\sqrt{\delta_{ij}x^ix^j}$, show it is a smooth distance function away from
  $p$, and that the gradient's integral curves are geodesics emanating from $p$.
\end{remark}

\begin{remark}[intersections of totally geodesic submanifolds]
  \label{rem:pet-ch5-ex-31}
  \lean{PetersenLib.exercise5_9_31}
  \source{petersen:page-0242}
  \uses{def:pet-ch5-fixed-point-set}
  \dcref{ch5:5.9.31}
  If $N_1,N_2\subset M$ are totally geodesic, show each component of $N_1\cap N_2$
  is a totally geodesic submanifold, with tangent space at $p$ the Zariski tangent
  space $T_pN_1\cap T_pN_2$.
\end{remark}

\begin{remark}[an isometry with differential $-I$]
  \label{rem:pet-ch5-ex-32}
  \lean{PetersenLib.exercise5_9_32}
  \source{petersen:page-0242}
  \uses{prop:pet-ch5-isometry-uniqueness}
  \dcref{ch5:5.9.32}
  Let $F:(M,g)\to(M,g)$ be an isometry fixing $p$. Show $DF|_p=-I$ on $T_pM$ iff
  $F^2=\mathrm{id}_M$ and $p$ is an isolated fixed point.
\end{remark}

\begin{remark}[functional distance equals distance when complete]
  \label{rem:pet-ch5-ex-33}
  \lean{PetersenLib.exercise5_9_33}
  \source{petersen:page-0242}
  \uses{def:pet-ch5-functional-distance, thm:pet-ch5-hopf-rinow}
  \dcref{ch5:5.9.33}
  Show that for a complete manifold, the functional distance $d_F$ equals the
  Riemannian distance $d$.
\end{remark}

\begin{remark}[regular exponential images are local energy minima]
  \label{rem:pet-ch5-ex-34}
  \lean{PetersenLib.exercise5_9_34}
  \source{petersen:page-0242}
  \uses{thm:pet-ch5-energy-minima-geodesics, lem:pet-ch5-exp-nonsingular-segment-domain}
  \dcref{ch5:5.9.34}
  Let $c:[0,1]\to M$ be a geodesic such that $\exp_{c(0)}$ is regular at $tc'(0)$
  for all $t\le1$. Show $c$ is a local minimum of the energy functional. Hint: the
  lift of $c$ via $\exp_{c(0)}$ is a minimizing geodesic in the pullback metric.
\end{remark}

\begin{remark}[bi-invariant metrics and the Lie exponential]
  \label{rem:pet-ch5-ex-35}
  \lean{PetersenLib.exercise5_9_35}
  \source{petersen:page-0242,petersen:page-0243}
  \uses{def:pet-ch5-exponential-map, thm:pet-ch5-hopf-rinow}
  \dcref{ch5:5.9.35}
  For a Lie group $G$ with bi-invariant (pseudo-)Riemannian metric: (1) show
  homomorphisms $\mathbb R\to G$ are precisely the integral curves of left-invariant
  vector fields through $e$; (2) show geodesics through $e$ are precisely such
  homomorphisms, so the Lie exponential coincides with the Riemannian exponential
  at $e$; (3) for a Riemannian (positive-definite) bi-invariant metric, show every
  $x\in G$ has a square root, using metric completeness; (4) show
  $\mathrm{SL}(n,\mathbb R)$ admits no bi-invariant Riemannian metric.
\end{remark}

\begin{remark}[Riemannian submersions are submetries]
  \label{rem:pet-ch5-ex-36}
  \lean{PetersenLib.exercise5_9_36}
  \source{petersen:page-0243}
  \uses{def:pet-ch5-submetry}
  \dcref{ch5:5.9.36}
  Show a Riemannian submersion is a submetry.
\end{remark}

\begin{remark}[completeness descends along submersions and gives a fibration]
  \label{rem:pet-ch5-ex-37}
  \lean{PetersenLib.exercise5_9_37}
  \source{petersen:page-0243}
  \uses{def:pet-ch5-submetry, thm:pet-ch5-hopf-rinow}
  \dcref{ch5:5.9.37}
  Let $F:(M,g_M)\to(N,g_N)$ be a Riemannian submersion. (1) Show $(N,g_N)$ is
  complete if $(M,g_M)$ is; (2) show $F$ is a fibration (locally
  $F^{-1}(U)\cong U\times F^{-1}(p)$) if $(M,g_M)$ is complete, and give a
  counterexample when it is not.
\end{remark}

\begin{remark}[commuting isometries and even-codimension fixed sets]
  \label{rem:pet-ch5-ex-38}
  \lean{PetersenLib.exercise5_9_38}
  \source{petersen:page-0243}
  \uses{rem:pet-ch5-fixed-point-codimension}
  \dcref{ch5:5.9.38}
  Let $S$ be a set of orientation-preserving isometries of $(M,g)$ that pairwise
  commute. Show each component of $\operatorname{Fix}(S)$ has even codimension.
\end{remark}

\begin{remark}[affine maps]
  \label{rem:pet-ch5-ex-39}
  \lean{PetersenLib.exercise5_9_39}
  \source{petersen:page-0243}
  \uses{def:pet-ch5-geodesic, prop:pet-ch5-isometry-uniqueness}
  \dcref{ch5:5.9.39}
  A local diffeomorphism $F:(M,g_M)\to(N,g_N)$ is \emph{affine} if
  $F_*(\nabla^M_XY)=\nabla^N_{F_*X}F_*Y$ for all vector fields $X,Y$. (1) Show
  affine maps carry geodesics to geodesics; (2) show $F$ is determined by $F(p)$
  and $DF|_p$; (3) give an affine $\mathbb R^n\to\mathbb R^n$ that is not an
  isometry.
\end{remark}

\begin{remark}[projective linear groups and isometry groups of space forms]
  \label{rem:pet-ch5-ex-40}
  \lean{PetersenLib.exercise5_9_40}
  \source{petersen:page-0243}
  \uses{rem:pet-ch5-ex-39, ex:pet-ch5-sphere-great-circles, ex:pet-ch5-hyperbolic-hyperbolas}
  \dcref{ch5:5.9.40}
  For $\mathbb{FP}^n$ ($\mathbb F=\mathbb R,\mathbb C$): (1) $\mathrm{GL}(n{+}1,
  \mathbb F)$ acts on $1$-dimensional subspaces of $\mathbb F^{n+1}$; (2) the
  kernel $H$ of this action is $\{\lambda I_{n+1}\}$, a normal subgroup; (3)
  $\mathrm{PGL}(n{+}1,\mathbb F)=\mathrm{GL}(n{+}1,\mathbb F)/H$ acts with each
  element determined by its value and differential at one point; (4) no
  Riemannian metric on $\mathbb{FP}^n$ makes this action isometric; (5) the action
  is affine (\cref{rem:pet-ch5-ex-39}) for the submersion metric; (6) the isometry
  group of $\mathbb{RP}^n$ is $\mathrm{PO}(n{+}1)$; (7) that of $\mathbb{CP}^n$ is
  $\mathrm{PU}(n{+}1)$; (8) that of $H^n(\mathbb R)$ is $\mathrm{PO}(n,1)$; (9) for
  $G=\{[\begin{smallmatrix}O&v\\0&1\end{smallmatrix}]\mid O\in O(n),v\in\mathbb R^n\}
  \subset\mathrm{GL}(n{+}1,\mathbb R)$ (as in Exercise 1.6.9), show $\mathrm{PG}=G$.
\end{remark}

\begin{remark}[the compact-open topology on $\operatorname{Iso}(M,g)$]
  \label{rem:pet-ch5-ex-41}
  \lean{PetersenLib.exercise5_9_41}
  \source{petersen:page-0244}
  \uses{thm:pet-ch5-myers-steenrod-lie-group, def:pet-ch5-proper-action}
  \dcref{ch5:5.9.41}
  Let $\operatorname{Diff}(M;K,O)=\{F\in\operatorname{Diff}(M)\mid F(K)\subset O\}$.
  (1) Finite intersections of such sets ($K$ compact, $O$ open) form a topology
  base — the compact-open topology; (2) it is second countable; (3) on a
  Riemannian manifold, convergence in it is uniform convergence on compact sets;
  (4) a sequence in $\operatorname{Iso}(M,g)$ converges in it iff it converges
  pointwise (Arzelà–Ascoli); (5) $\operatorname{Iso}_0(M,g)$ is locally compact
  (Arzelà–Ascoli); (6) $\operatorname{Iso}_0(M,g)$ is compact; (7)
  $\operatorname{Iso}(M,g)$ acts properly on $M$; (8) $F\mapsto(F(p),DF|_p)$ is
  continuous on $\operatorname{Iso}(M,g)$; (9) this evaluation map is a
  homeomorphism onto its image.
\end{remark}

\begin{remark}[Bianchi identities and the curvature tensor in normal coordinates]
  \label{rem:pet-ch5-ex-42}
  \lean{PetersenLib.exercise5_9_42}
  \source{petersen:page-0244,petersen:page-0245}
  \uses{lem:pet-ch5-normal-coords-first-order}
  \dcref{ch5:5.9.42}
  In exponential normal coordinates at $p$, all computed at $p$: (1) taking three
  derivatives of $x^i=\sum_sg_{is}x^s$ (as in
  \cref{lem:pet-ch5-normal-coords-first-order}), show
  $\partial_i\partial_kg_{jl}+\partial_j\partial_ig_{kl}+\partial_k\partial_jg_{il}=0$;
  (2) using all four cyclic versions, show
  $\partial_i\partial_jg_{kl}=\partial_k\partial_lg_{ij}$; (3) using the curvature
  tensor formula in normal coordinates, show
  $R_{ikjl}=\partial_i\partial_lg_{kj}-\partial_i\partial_jg_{kl}$; (4) combine (1)
  and (3) to get $\partial_i\partial_jg_{kl}=\tfrac13(R_{ikjl}+R_{jkil})$.
\end{remark}

\begin{remark}[Riemannian volume expansion (A. Gray)]
  \label{rem:pet-ch5-ex-43}
  \lean{PetersenLib.exercise5_9_43}
  \source{petersen:page-0246}
  \uses{rem:pet-ch5-ex-42}
  \dcref{ch5:5.9.43}
  With notation as in \cref{rem:pet-ch5-ex-42}: (1) show
  $\sqrt{\det(g_{ij})}=1-\tfrac16\mathrm{Ric}_{ij}x^ix^j+O(|x|^3)$; (2) (A. Gray)
  show $\operatorname{vol}B(p,r)=\omega_nr^n\big(1-\tfrac{\mathrm{scal}(p)}{6(n+2)}r^2
  +O(r^3)\big)$, $\omega_n=\operatorname{vol}(B(0,1)\subset\mathbb R^n)$, by
  expanding the volume integral in polar coordinates using (1).
\end{remark}
